\documentclass[12pt, unicode, a4paper]{article}
\usepackage{amsmath,amssymb,ascmac,amsthm,graphicx}
\usepackage{zxjatype}
\usepackage{graphics}
\usepackage[ipaex]{zxjafont}
\usepackage{setspace}
\usepackage[top=25truemm,bottom=25truemm,left=25truemm,right=25truemm]{geometry}
\usepackage{url}
\usepackage{newtxtext,newtxmath}
\usepackage{lscape}
\usepackage{adjustbox}
\usepackage{threeparttable}
\usepackage{tabularx}
\usepackage{siunitx}
%\usepackage{fontspec}
\usepackage[longnamesfirst]{natbib}
\usepackage{multibib}
% Both the main auxiliary file and app.aux require BibTeX; latexmk handles both.
\newcites{app}{References for the Appendix}
\usepackage[hidelinks]{hyperref}
\newtheorem{hypothesis}{Hypothesis}
\newtheorem{prediction}{Prediction}
\newtheorem{lemm}{Lemma}
\newtheorem{cor}{Corollary}
\newtheorem{res}{Result}
\newtheorem{ass}{Assumption}
\newtheorem{rem}{Remark}
\newtheorem{prop}{Proposition}
\newtheorem{defi}{Definition}
\newtheorem{exam}{Example}
\newtheorem{fact}{Fact}
\newtheorem{result}{Result}
\newtheorem{obs}{Observation}
\newcommand{\1}{\mbox{1}\hspace{-0.25em}\mbox{l}} % Indicator function
\usepackage{subfig}
\usepackage{tikz}
\usepackage{epigraph}
\usepackage{booktabs}
\usepackage{floatrow}
\usepackage{multirow}
\usepackage{comment}
\usepackage{titletoc}
\usepackage{adjustbox}
\usepackage{ulem}
\usepackage{makecell}
\usepackage{booktabs}
\usepackage{longtable}
\usepackage{array}
\setlength \epigraphwidth {0.9\linewidth}
\setlength \epigraphrule {0pt}
\newcommand{\argmax}{\mathop{\mathrm{arg~max}}\limits}


\addtolength{\footskip}{5mm} % Lower the page number a bit. 
\linespread{1.2} % Slightly more line space for readability

% Use this command to comment in the document. 
\newcommand{\yama}[1]{{\color{blue}{#1}}}
\newcommand{\aiba}[1]{{\color{red}{#1}}}



% Color definitions and commands. 
\definecolor{tpurple}{rgb}{.32,.18,.5}
\definecolor{DarkBlue}{rgb}{.1,.1,.5}
\definecolor{DarkGreen}{rgb}{0,.4,0}
\definecolor{Princeton}{rgb}{1.0, 0.56, 0.0}
\definecolor{darkviolet}{rgb}{0.58, 0.0, 0.83}
\definecolor{dpink}{rgb}{0.93, 0.23, 0.51}

% Cite color
%\usepackage{float}
%\usepackage{caption}
%\usepackage{xcolor}
\usepackage{hyperref}
\hypersetup{
bookmarksnumbered=true,
bookmarkstype=true,
colorlinks,
linkcolor=DarkBlue,
citecolor=DarkBlue,
urlcolor=DarkBlue}
\urlstyle{rm}

\begin{document}
%\setlength{\baselineskip}{20pt} 
%\doublespacing

\title{Bases and the City: \\ Spatial Economics of US Military Bases in Okinawa
\thanks{
An earlier version of the paper was circulated under the title ``Urbanization without Industrialization: Evidence from US Bases in Okinawa."
We are grateful for the helpful comments from Kristian Behrens, Jan Brueckner, Dan Bogart, Juan Pablo Chauvin, Sean Chen, Don Davis, Ellora Derenoncourt, Mouhcine Guettabi, Marco Gonzalez-Navarro, Jun Goto, Shizuka Inoue, Takuma Kamada, Hans Koster, D\'{a}vid Nagy, Shimeng Liu, Jun Oshiro, Steve Redding, Yasuhiro Sato, Takeru Sugasawa, Kohei Takeda, Kensuke Teshima, Chamna Yoon, Miguel Zerecero, Junfu Zhang, Ben Zou, and numerous seminar and conference participants. We thank Tatsuya Kusumi for his tremendous research support. We also thank Tsuneaki Nakajima and Koichiro Takagi for their research assistance and Shinji Nakano for his assistance in digitizing the GIS data. 
We acknowledge the microdata access to the Person Trip Survey by the Okinawa Prefectural Government and the custom-made tabulation provided by the Ministry of Internal Affairs and Communications.
Aiba appreciates financial support from the Center for Real Estate Innovation, the University of Tokyo, and JSPS KAKENHI Grant Numbers 21J13197 and JP26K16348. Yamagishi appreciates the financial support from the Institute of Economic Research, Hitotsubashi University, the Funai Foundation for Information Technology, and JSPS KAKENHI Grant Number JP25K00629. All remaining errors are ours.
}}
\author{
Ikuto Aiba\thanks{Faculty of Economics, Seikei University \& Center for Real Estate Innovation, the University of Tokyo. E-mail: \href{mailto:ikuto-aiba@econ.seikei.ac.jp}{\color{blue} ikuto-aiba@econ.seikei.ac.jp}.} 
\and 
Atsushi Yamagishi\thanks{Graduate School of Economics and Management, Tohoku University \& IZA. E-mail: \href{mailto:atsushiyamagishi.econ@gmail.com}{\color{blue} atsushiyamagishi.econ@gmail.com}.}
}


\date{ \today \\ }


\maketitle
\begin{abstract}
We examine how external expenditure and employment shape urbanization and economic structure. Using newly digitized data from Okinawa, Japan, we document urbanization without industrialization around US military bases, with service-sector expansion but little manufacturing growth. We calibrate a quantitative spatial model to 1970. In a counterfactual that removes base-sourced expenditure, land-rent income, and employment, base-exposed municipalities lose residents and shift from services toward agriculture and manufacturing. Okinawa-wide welfare falls by 7.25\% and GNI by 11.99\%, although base-related consumption demand and employment accounted for roughly 30\% of GNI. Sectoral reallocation mitigates these losses alongside sizable migration. The model-implied effects of consumption demand and employment match the observed pattern qualitatively but fall short quantitatively. Evidence on infrastructure and early commercial access suggests that past investment in infrastructure and durable capital may help explain this gap.\end{abstract}

\noindent \textbf{Keywords}: Military bases, Service sector, Quantitative spatial model, Geoeconomics, Urbanization 

\noindent \textbf{JEL classification}: O14, O18, N40, R11, R12

\newpage
%%%%%%%%%%%%%%%%%%%%%
% Introduction
%%%%%%%%%%%%%%%%%%%%%



\begin{quote}
\begin{flushright}
\textit{Dollars flowed into Koza, then on to Naha—and Naha prospered too.} \\
Seiji Tokutomi\footnote{
Quotes are taken from \url{https://www.peace-museum.okinawa.jp/yogawari/testimony/r413/} (in Japanese, last accessed on July 13, 2026). 
Tokutomi ran a restaurant in Koza city near the Kadena Air base in Okinawa. Naha is Okinawa's largest city. The original sources of the quotation are freely translated from Japanese by the authors.
% https://www.peace-museum.okinawa.jp/yogawari/testimony/r413/
}
\end{flushright}
\end{quote}



\section{Introduction}
As military spending rises amid geopolitical tensions, its economic effects are attracting growing attention in policy discussions. In particular, military bases are expected to support local economies through consumption demand generated by military personnel and direct employment of local residents.\footnote{
In reviewing the US base-closure process, the U.S. General Accounting Office considers local purchases by military personnel, trainees, and civilian employees when assessing employment effects \citep{USGAO1995}. A congressional hearing on Fort Ord likewise highlights the loss of civilian base jobs and secondary employment following closure \citep{USHouse2002FortOrd}.}
For instance, preparations for a German brigade in Lithuania have raised expectations that ``new jobs will emerge, businesses will expand, demand for housing and services will increase.''\footnote{Alexander Welscher, \url{https://www.ahk-balt.org/en/publikationen/baltic-business-quarterly/bbq-autumn-2025/getting-ready-for-the-german-brigade}, last accessed on October 2, 2026.} Conversely, reporting on possible US troop withdrawals from Germany, Reuters describes a base town's dependence on ``thousands of well-paying jobs for local people and steady stream of customers for local businesses.''\footnote{\url{https://www.reuters.com/world/small-town-germany-braces-end-decades-life-with-us-troops-2026-05-07/}, last accessed on September 23, 2026.}
Yet the impact of the local spending and base employment on the host economy is not straightforward to evaluate in the spatial economy. Workers can move across locations and sectors, while commuting and shopping transmit effects beyond host municipalities. How do base-related consumption demand and employment reshape the spatial and sectoral distribution of economic activity, and what are their aggregate effects on income and workers' welfare?


This paper addresses these questions using the distinctive case of Okinawa. The U.S. military constructed many bases for strategic reasons during and immediately after World War II, generating substantial spatial variation in base exposure. The bases brought large and spatially uneven income inflows through purchases of goods and services, employment of Okinawan residents, and rental payments to local landowners. Okinawa's GNI statistics separately record these three components and indicate that around 40\% of aggregate income came from the US bases in 1955. These separate accounts allow us to discipline the scale and composition of base-related flows in a quantitative spatial model and assess their effects on local development and aggregate outcomes.

We collect and digitize a comprehensive dataset of the Okinawan economy, allowing us to observe various economic variables within the Okinawa island at the municipality level. Using these data, we first document two stylized facts on the relationship between US bases and local economic outcomes. First, a 10 percentage-point higher base-area share is associated with a roughly 36\% higher ratio of 1970 to 1938 population density, implying substantial urbanization around the US bases.
Second, areas that were highly exposed to the US bases experienced a large increase in the service employment share.
In stark contrast, changes in the manufacturing employment share show little association with base exposure. Taken together, our stylized facts suggest that areas around the US bases were characterized by ``urbanization without industrialization."

%To interpret these patterns, the paper develops a quantitative spatial model in which military bases act as consumers, employers, and land users within a local economy. The model is calibrated to Okinawa in 1970 and used to conduct a counterfactual that shuts down U.S.-base-sourced consumption and employment. The results indicate that the bases were a central force behind urbanization without industrialization around their locations, with sizable positive effects on aggregate income and worker welfare. More broadly, the analysis shows how geopolitical infrastructures can reshape local development by altering the spatial distribution of demand, employment, and land use, and by favoring non-tradable services over tradable production.


Motivated by these stylized facts, we construct a quantitative spatial model that explicitly incorporates US-base sourced consumption demand and employment. 
Combined with the GNI statistics recording the bases' separate contribution to Okinawa's aggregate economy, the model allows us to quantify the effects of their consumption demand and employment while holding other local fundamentals fixed, including those that may reflect infrastructure and durable capital, noise, or pollution.
Moreover, our quantitative spatial model is a general-equilibrium one, allowing us to investigate how the US-base sourced consumption and employment, which propagate through within-island migration, commuting, and cross-border shopping, translate into the aggregate impact on income and workers' welfare.

Our model has many locations, embedded in a larger economy, and workers choose their residence, workplace, and sector to maximize utility, paying commuting costs when residence and workplace differ. Each location has four sectors: agriculture, manufacturing, services, and the US military base sector. We assume that agriculture and manufacturing goods are tradable across locations while non-tradable services are not, although we allow workers to consume differentiated non-tradable services at different locations through cross-border shopping. Okinawan workers at the US military base do not produce any good in the Okinawan economy as their work has a militaristic value that is not traded in a market. The US base takes up a certain fraction of land in a location and pays rents for it, employs Okinawan residents, and demands tradable and non-tradable goods. In addition, developers provide floor space using tradable goods and land, implying that the presence of US bases may increase floor space prices by reducing available land for the private sector.

We calibrate our model to the Okinawan economy in 1970. Our calibration goes in steps. First, we calibrate several model parameters from existing studies. Second, we estimate gravity equations for commuting and shopping to obtain the travel cost parameters. Third, using residence and workplace employment by sector, we recover normalized spatial profiles of wages and amenities. Fourth, we compute the local non-tradable price level using the non-tradable goods market clearing condition, matching base land rents and US personnel's purchases to their respective GNI shares. Finally, we recover the fundamental location-specific amenities and productivity. Using this calibration procedure, our model exactly matches the observed population and employment distribution by sector. The model also successfully replicates untargeted moments: observed bilateral commuting probabilities and housing prices. 

To understand the effects of US-base-sourced expenditure and employment, we use our calibrated model to perform a counterfactual analysis that jointly removes expenditure and land-rent income associated with the bases and employment of Okinawan workers at the bases, while keeping base land unavailable for private use and Okinawa's aggregate worker population fixed.
The counterfactual produces a sharp spatial reversal. A 10 percentage-point increase in the municipal base-area share is associated with a 4.77 percentage-point larger population decline in the counterfactual scenario. The corresponding decline in the  tertiary sector is 2.16 percentage points larger, while the agricultural and manufacturing shares increase. Thus, US-base-sourced expenditure and employment contributed substantially to urbanization without industrialization around the bases. 
The different tradability of goods by sector drives these results. While US-base-sourced consumption and employment increase demand for both tradable and non-tradable goods, non-tradable demand must be satisfied locally, whereas tradable goods can be supplied from outside Okinawa at fixed prices in the benchmark. As a result, areas around the US bases have higher service employment shares to satisfy the larger service demand. Separating the channels, we find that consumption has larger effects on civilian service employment, whereas base employment has larger effects on population concentration. 

In the aggregate, the benchmark counterfactual reduces model-implied workers' welfare by 7.25\% and gross national income (GNI) by 11.99\%, accompanied by the growth of manufacturing and agricultural employment shares and the shrinkage of the non-tradable sector.
Although the losses are large in absolute terms, they are modest relative to the scale of the initial shock: base-related consumption and employment accounted for roughly 30\% of Okinawa's GNI in 1970. In this sense, our findings are broadly consistent with the ``Dutch disease'' hypothesis that a large facility generating substantial demand and employment can draw workers away from manufacturing and reduce its agglomeration economies \citep{gollin2016urbanization,allcott2018dutch}. Once base-related demand and employment disappear, sectoral and spatial mobility can substantially cushion the economic cost. We find that sectoral reallocation is the primary adjustment margin, despite a sizable within-island migration response. Finally, returning base land to private use dampens population outflows from base-exposed municipalities and reduces the benchmark welfare loss from 7.25\% to 1.95\%, while leaving the economy's sectoral composition largely unchanged.

We also quantify direct and spillover effects using local and aggregate elasticities. Base-related consumption generates relatively localized service employment gains, while manufacturing losses spread more widely; additional base hiring likewise causes most civilian employment displacement outside the municipality receiving the new jobs.
Finally, by modifying our structural model, we assess the external validity of our results under alternative assumptions about external migration and manufacturing import supply. Allowing external migration leaves workers' welfare and sectoral reallocation close to the benchmark but increases the GNI loss through population outflows. Less elastic manufacturing import supply weakens, and can reverse, the aggregate manufacturing crowding-out effect because base-related demand then raises manufacturing prices and supports local production. 

We close our paper by highlighting the potential importance of investment in infrastructure and durable capital accompanying the base development. Although our counterfactual analysis predicts that US-base-sourced consumption demand and employment induce ``urbanization without industrialization," which we indeed observe in the data, the predicted tertiary-share response is smaller than the observed association. 
This comparison helps assess how much of the historical association can be accounted for by the ongoing demand and employment channels emphasized in the policy discussions above.
We highlight one possible factor that could fill in this gap: investment in infrastructure and durable capital around the bases. Indeed, we find that infrastructure quality, proxied by water supply and road quality, is positively correlated with exposure to the US bases. We also find that early designation of road corridors for authorized markets serving US military personnel is correlated with base exposure, suggesting that early access to military customers could have encouraged investment in infrastructure by the government or in durable capital by the private sector, yielding persistent advantages around the US bases. Conditional on base exposure and geographic controls, water supply is also positively associated with calibrated amenities, while road quality and early commercial access are positively associated with productivity. Overall, although the effects of US-sourced consumption demand and employment qualitatively match the observed pattern of spatial economic development in Okinawa, our structural analysis suggests that attributing all observed development to these ongoing flows, as a naive observation might suggest, could overstate their effects.


%%%%%%%%%%%%%%%%%%%%%
% Related literature
%%%%%%%%%%%%%%%%%%%%%



This paper contributes to several strands of literature. Most directly related is research on the local economic effects of military activity.
Existing studies estimate how military spending and employment affect local output and employment (e.g., \citealt{hooker1997effects}, \citeyear{hooker2001measuring}; \citealt{paloyo2010regional}; \citealt{nakamura2014fiscal}; \citealt{lee2018regional}; \citealt{zou2018local}; \citealt{biolsi2019local}; \citealt{dahlberg2024coping}; \citealt{brunet2026stimulus}).\footnote{
Existing quantitative assessments of US bases in Okinawa have largely relied on input--output analysis to estimate demand-induced output and employment effects, including studies of base employment and land redevelopment (e.g., \citealt{okinawa07}; \citealt{kusuyama12}; \citealt{tomikawa18}).}
Using newly digitized municipal data, we document ``urbanization without industrialization'' around US bases in Okinawa: greater base exposure is associated with population concentration and service-sector expansion, but has little association with changes in the manufacturing employment share. We combine these data with GNI accounts that separately report base-related purchases, labor income, and land rents. Matching their shares in aggregate income allows us to discipline the scale and composition of base-related flows in a quantitative spatial model. By tracing sectoral reallocation and spatial spillovers, we connect local employment responses to aggregate income and workers' welfare. We also distinguish the effects of ongoing flows from the broader historical association between bases and local development, providing suggestive evidence on the contribution of past investment in infrastructure and durable capital.

Our paper also relates to the literature on urbanization and structural transformation. Urbanization has historically accompanied the movement of labor out of agriculture (\citealt{kuznets1973modern}; \citealt{michaels2012urbanization}; \citealt{bryan2025spatial}), but need not involve sustained manufacturing growth (\citealt{jedwab2015urbanization}; \citealt{henderson2020urbanization}). Related research explains service-led growth and premature deindustrialization through technological change and trade (\citealt{rodrik2016premature}), productivity growth in consumer services (\citealt{fan2023india}), and differences in sectoral productivity and technology adoption (\citealt{huneeus2023heterogeneous}; \citealt{fujiwara2024premature}). Most closely related, \citet{gollin2016urbanization} show how resource income supports consumption cities dominated by non-tradable services. Building on this mechanism, we examine military bases as a source of externally financed expenditure and direct employment. The Okinawan setting allows us to quantify the distinct roles of spending and hiring in population concentration and sectoral reallocation. Spending sustains civilian services, while direct hiring plays a larger role in population concentration. We distinguish these local patterns from the islandwide manufacturing response, which our model shows depends on the elasticity of manufacturing import supply.


% Methodologically, our paper relates to the literature quantitative spatial model (e.g., \citealt{michaels2012urbanization}; \citealt{desmet2014spatial}; \citealt{faber2019tourism}; \citealt{hao2020china}; \citealt{fajgelbaum2022argentina}; \citealt{eckert2023ssc}; \citealt{takeda2023structural};  \citealt{coeurdacier2025structural}). 
% Building on this literature, we propose a new quantitative spatial model that accommodates a facility (the US base in our context) that brings in income from the outside economy through purchasing goods, employing labor, and renting land for its operation.\footnote{
% Other military bases around the world, including the US bases in the US and other countries such as Germany, Italy, and South Korea, can be direct application of our model. A university or a prison could also be an example that fits this situation (c.f., \citealt{liu2015spillovers}; \citealt{chirakijja2024local}). To evaluate different types of public facilities,  \cite{loumeau2023locating} and \cite{ducruet2024all} develop a quantitative spatial model to evaluate the welfare effect of schools and ports, respectively.
% }  
% Calibrating this model to the case of Okinawa with rich information on the contribution of the US bases to the aggregate economy, we analyze how external consumption demand and employment opportunities might lead to urbanization without industrialization. 
Methodologically, our framework builds on quantitative spatial models with commuting (e.g., \citealt{ahlfeldt2015economics}; \citealt{monte2018commuting}; \citealt{heblich2020making}; \citealt{miyauchi2022economics}; \citealt{dingel2023spatial}; \citealt{tsivanidis2019evaluating}; \citealt{takeda2024economic}; see \citealt{redding2024survey} for a survey). We extend this framework to incorporate a facility that purchases goods and services, employs local workers, and occupies land. Commuting, residential mobility, and shopping across municipalities transmit its effects beyond the host location, allowing us to distinguish local gains from aggregate effects. The framework can also be applied to other externally financed facilities, such as public offices and universities, that combine local spending, employment, and land use.

The rest of this paper is organized as follows. Section \ref{sec:inst_data} introduces the institutional background of Okinawa and our data. Section \ref{sec:reducedform} presents our stylized facts on the relationship between the US bases and urbanization without industrialization. Section \ref{sec:model} introduces our quantitative spatial model that incorporates the US-base sourced consumption and employment at the US bases. Section \ref{sec:calibration} calibrates our model using the data from Okinawa in 1970. Section \ref{sec:counterfactuals} presents the counterfactual effects of US-base-sourced consumption and employment, reports direct and spillover elasticities, and examines how the benchmark results depend on external mobility and manufacturing import supply. Section \ref{sec:discussion_extensions} examines past investment in infrastructure and durable capital as a possible additional contributor to economic development. Section \ref{sec:conclusion} concludes.


%%%%%%%%%%%%%%%%%%%%%
% Background and data
%%%%%%%%%%%%%%%%%%%%%


\section{Institutional Background and Data}\label{sec:inst_data}

Section \ref{sec:instutitional} introduces the institutional background of this study and Section \ref{sec:data} introduces our data.

\subsection{Institutional background: US Bases in Okinawa}\label{sec:instutitional}

To facilitate the understanding of our study context, we briefly explain the institutional background of Okinawa and its relationship to the US bases.\footnote{
See, for instance, \cite{okinawa2000}, \cite{taira2012}, \cite{toriyama2013},  \cite{franks2023}, and \cite{oshiro23} for a more detailed description.
}
The Okinawa island is located to the south-west of the main island of Japan.\footnote{
Okinawa was also historically called ``Ryukyus.''  In this paper, for simplicity, we use the name ``Okinawa'' throughout the paper. 
} 
Although Okinawa Prefecture of Japan contains islands other than the main island, we refer to the Okinawa main island by saying ``Okinawa'' unless explicitly noted. 
As of 2020, the population is around 1.3 million and its size is around 1,200 square kilometers.
Figure \ref{fig:pop_longrun} plots the long-run trend of the total population of Okinawa island, showing the increasing trend throughout the post-WW2 period.
The growth in the total population accompanied substantial urbanization, especially in the central and southern parts of the island, such as Naha city and Koza city. 
Figure \ref{fig:empshare_longrun} documents that such an urbanization process entailed substantial growth in the service employment share. Notably, the manufacturing employment share did not see substantial growth and hit a plateau as early as in the 1960's. 
This development process is quite different from that of other Japanese regions. The post-war economy of Mainland Japan experienced a major shift of manufacturing sectors from light to heavy industries between the mid 1950's and the early 1970's, and then shifted to a service economy after the mid 1970's (\citealt{fujita1997regional}), which is in line with Petty's law and is a widely observed pattern of economic development. Relative to this, Figure \ref{fig:empshare_longrun} highlights the uniqueness of Okinawa's development process.
 

\begin{figure}[t]
    \subfloat[Total population\label{fig:pop_longrun}]{%
      \includegraphics[width=0.475\textwidth]{figure/total_pop_plot.pdf}
    }
    \hfill
    \subfloat[Employment share by sector \label{fig:empshare_longrun}]{%
      \includegraphics[width=0.475\textwidth]{figure/industry_share_plot.pdf}
    }
    \vspace{0.2in}
    \caption{Long-run dynamics of the total population and employment share in Okinawa}
    \vspace{-0.2in}
    \floatfoot{\textbf{Note}: Figures \ref{fig:pop_longrun} and \ref{fig:empshare_longrun} plot the long-run trends of the total population and the sectoral employment share from 1934 to 2015 in the Okinawa main island, respectively. In Figure \ref{fig:empshare_longrun}, the lines with red circles, blue squares, and green triangles represent the employment shares of the tertiary, secondary, and primary sectors, respectively.
    }
    \label{fig:total_longrun}
\end{figure}

The substantial presence of US military bases is a distinctive feature of Okinawa's postwar economy. Their spatial distribution reflects, at least in part, the sequence of occupation and military construction during the Battle of Okinawa in 1945. Figure \ref{fig:okinawa_map} places the geography of the landings and early occupation alongside the distribution of US bases in 1970. 
%Proximity to the landing beaches and the geography of Japanese defenses jointly influenced when different parts of the island became available for military use. Because construction began before the battle ended, these differences in occupation timing helped shape where military facilities were established.
Panel (a) marks the actual landing point of the US forces (Toguchi Beach) on April 1, 1945. The figure also presents two coastal sectors in an alternative plan for landing: a Marine landing between Minatogawa and Chinen Point, followed by an Army landing between Kuba and Yonabaru. The selection criteria for the landing point were driven primarily by military and logistical requirements, rather than by the local economy's development potential (see Appendix \ref{appsec:base_location_endogeneity} for more details).
After the landing, US forces advanced rapidly through central Okinawa and into the north during the first two weeks, occupying the colored areas in Panel (a).
In contrast, progress southward was much slower because Japanese forces concentrated their defenses in the southern part of the island, including the positions around Shuri. Resistance there delayed the occupation of southern locations even when they were relatively close to the landing beaches, and fighting in the south continued for several months. The early availability of land for base construction therefore depended on both proximity to the landing area and the geography of Japanese resistance.

Importantly, the timing of the occupation, at least in part, predicts the local presence of the US bases because the US military started constructing the bases even during the battle of Okinawa.
Indeed, they started constructing the bases in an early-occupation area as they needed facilities to support the ongoing Okinawa campaign and to prepare for further operations against the Japanese mainland.
Panel (b) shows the resulting concentration of US bases in central and northern Okinawa in 1970, which visually exhibits a strong correlation with the early-occupation areas in Panel (a).\footnote{ 
Post-war land acquisition and construction, including during the 1950s, also contributed to the base construction. During the 1950s, the US administration used forcible land acquisition without seriously considering its economic consequences to expand military facilities for the Korean War and subsequent strategic needs  (\citealt{taira2012}; \citealt{toriyama2013}; \citealt{franks2023}).
}
Overall, these historical accounts suggest that the US base locations were determined largely for non-economic reasons.\footnote{
Appendix \ref{sec:reduced_additional} examines prewar trends. Consistent with this, we do not find evidence of the pre-war trend that could explain a large part of post-war economic development.
}

\begin{figure}[t]
\centering
  \includegraphics[width=\textwidth]{figure/section2_wartime_base_locations.pdf}
    \caption{Wartime occupation and US base locations in Okinawa}
    \label{fig:okinawa_map}
    \floatfoot{\textbf{Note}: The stars labelled ``Landing point'' mark the reference point at Toguchi Beach in both panels. In panel (a), the northern and southern purple sectors labelled ``Alternative landing area'' correspond to Kuba--Yonabaru and Minatogawa--Chinen Point, respectively. Panel (a) also indicates the broad concentration of Japanese defenses in the south. Panel (b) shows US bases as of 1970. Both panels use the same scale and shade the early-occupation municipalities described in Appendix \ref{appsec:base_location_endogeneity}; municipal boundaries are not precise wartime fronts or defensive lines. Only the main island is displayed. Sources: project GIS data and \href{http://gis.pref.okinawa.jp/OpenData/}{Okinawa Prefecture base-boundary data}; the \href{https://www8.cao.go.jp/okinawa/okinawasen/gaiyou/gaiyou.html}{Cabinet Office account}; the US Army, Marine Corps, and Navy histories cited in the text. Candidate-sector locators use the \href{https://doi.org/10.20676/00000448}{CODH historical place-name dataset} and the Okinawa Prefectural Archives.}
\end{figure}



Even after the war, the US bases have exhibited their substantial presence in the Okinawa island.
The US bases constitute a large share of the geographical area of Okinawa Island, reaching up to about 20\% in 1970, and it stabilizes at 15\% today.
In addition to their large geographical size, the US bases are an important economic actor that has greatly affected the Okinawan economy. 
Broadly speaking, the US bases are important as consumers, employers, and tenants. 
As consumers, the US military personnel consume goods and services provided by the Okinawan people, such as foods, restaurants, and daily commodities.\footnote{
Military contractors often involve large spending on weapons. However, the available data on the spending category of the US indicates that the US sourced virtually no weapons from Okinawa (\citealt{ReversionIssuesStudyGroup}). 
The only primary contract for military activities in Okinawa is construction, which we classify as manufacturing following the Japan Standard Industrial Classification.
%https://www3.archives.pref.okinawa.jp/RDA/data01/R00004916B/index.html?title=%E7%B5%8C%E6%B8%88%E5%A7%94%E5%93%A1%E4%BC%9A%E8%B3%87%E6%96%99%E3%81%AB%E9%96%A2%E3%81%99%E3%82%8B%E6%9B%B8%E9%A1%9E%20&page=481
}
As employers, the US bases employ Okinawan people for their daily operations, such as cleaning and babysitting. As tenants, the US bases pay the land rents to landowners.
In all these three dimensions, the US bases created large consumption demand and employment around the US bases.

Figure \ref{fig:okinawa_GDPshare} shows the contribution of the US bases as the share of Okinawa's Gross National Income (GNI), separately for goods and service consumption, labor income payments for employment at US bases, and land rent payments for the US base area to local landlords.
Note that these data present a distinctive feature of Okinawa for measuring the bases' economic impacts: In other contexts such as the US and mainland Japan, their GNI statistics do not separately record the contribution of the US bases.
The US bases constituted around 40\% of Okinawa's GNI in 1955. The share declined sharply in the 1960's and 1970's.\footnote{This is due to the growth of the Okinawan economy and the shrinkage of the US-base sourced income due to reversion of Okinawa and the US's withdrawal from the Vietnam war.}
The share has been around 5\% for the past 40 years. 
Unpacking the US-base sourced income before the return of Okinawa to Japan in 1972, the income from consumption of goods and services took up the highest share, the labor income of Okinawan residents at the US bases came next, and the land rent payment had the lowest share. After the return, however, the GNI shares of these three sources of income are close to each other due to the substantial drop in the first two. 
Given the very large share of the US-base sourced income in the GNI, reaching up to 40\% in the 1950's, we may reasonably conjecture that the US bases have influenced the post-war economic development of Okinawa. 
As a first step for understanding such an impact, the next section documents stylized facts on how areas that were highly exposed to the US bases are characterized by their urbanization and industrial structure. 


\begin{figure}[t]
  \includegraphics[width=0.66\textwidth]{figure/USshare_GNI.pdf}
    \caption{GNI share of the US bases in Okinawa island}
    \label{fig:okinawa_GDPshare}
    \floatfoot{Note: This figure plots the income sourced from the US military as the share of Okinawa's Gross National Income from 1955 to 2020. We used statistics from the Ryukyus Statistical Yearbook (\textit{Ryukyus Toukei Nenkan}) and Okinawa Statistical Yearbook (\textit{Okinawa Toukei Nenkan}). The black solid line is the sum of the three sources of US-base income: goods and service consumption by the US, labor income of Okinawans who are employed at the US bases, and land rent the US pays to Okinawan landowners for their bases. 
}
\end{figure}


\subsection{Data}\label{sec:data}

We collect, and newly digitize when needed, the municipality-level data on various outcomes.
Below we briefly describe the data used for our baseline analysis. Appendix \ref{appsec:datadetails} provides more details on the data description and Table \ref{tab:data_sources} summarizes our data sources.

\paragraph{Spatial unit and scope of the analysis.}

We focus on the data of Okinawa Prefecture. Okinawa Prefecture consists of the main island, which accounts for about 50\% of the total area and about 90\% of the population as of today, and the other smaller islands. 
Since most US bases are located in the Okinawa main island and a very large part of economic activity takes place there, we focus on the Okinawa main island in our main analysis unless explicitly noted.

Throughout the paper, our analysis is conducted at the municipality level. 
There are 35 municipalities on Okinawa Island as of 1970, although the number of municipalities changes over time due to mergers and divisions (see Figure \ref{fig:cf_pop_change_map} for an illustration of municipal boundaries in 1970).  
To address the mergers and divisions of municipalities during our sample period, we aggregate municipalities to the largest definitions of municipalities when our analysis involves different years (see Appendix \ref{appsec:datadetails} for details). 

\paragraph{Population.}

For population at the municipality level after 1945, we use the Population Census data, which have been available every five years since 1950. The US-occupation-period data are newly digitized.
The population includes only the Japanese population, implying that the US citizens are not covered by the Census. 
The population census also reports the number of workers by industry category living in each municipality.
The census also tells us the number of people working for the US bases.
In addition, we obtain population data by industry category in 1934 and 1938 from the Okinawa Prefectural Statistical Yearbook (\textit{Okinawa-ken toukei sho}).

\paragraph{Employment.} At the workplace level, we use the 1970 Population Census to obtain the number of workers by industry category in each municipality. 
The census also reports workplace counts of Okinawan workers employed at the US bases, which we treat as a separate sector in the model.
Although sourced from the same dataset, workplace employment by sector differs from the number of resident workers in that sector because of commuting across municipalities.

\paragraph{Size of the US bases.}

In our reduced-form analysis, we use the area share of the US bases in a municipality taken from official statistics, which we use as an exposure index for US-base sourced consumption and employment.
This measure is motivated by the assumption that the amount of demand of the US bases is proportional to the size of each base, and each base spends its money locally. 
We also observe in \cite{Komeito69} the number of US military personnel and Okinawan workers at each base. We use these data to validate our assumption that the US base area share indeed captures the base's  economic importance (see Section \ref{sec:reduced_additional}).

\paragraph{GNI and US-base-sourced income.}
Gross National Income (GNI) is reported annually from 1955 in the Ryukyus (Okinawa) Statistical Yearbook. We use the prefecture-wide series as an approximation for the main island. The statistics distinguish base-related purchases of goods and services, salaries of Okinawan base workers, and land rents.\footnote{US-base-sourced income includes receipts from base-related transactions, including payments financed by the Japanese government, but excludes economic aid and intergovernmental transfers, including base-related grants to local governments.} Figure \ref{fig:okinawa_GDPshare} shows these income shares over time, and Appendix \ref{appsec:datadetails} provides further details.

\paragraph{Road network and trip data.} We use the travel-level microdata of the Okinawa Main Island Central and Southern Urban Area Person Trip Survey in 1977.\footnote{The data do not cover the northern part of the Okinawa island. 
We have also used the same survey conducted in 1989 and 2006 and found that our results are hardly affected.}  For each respondent, the survey asks for details of the trips they make in a day. 
The survey asks about the purpose of each trip and the worker's industry, allowing us to identify commuting and shopping trips and to classify commuters by sector. 
Based on these data, we construct separate bilateral commuting matrices for workers in agriculture, manufacturing, and non-tradable services, as well as a bilateral shopping matrix, and use them to estimate gravity models. The bilateral travel distance is calculated based on the newly-digitized road network data in 1973, which are taken from the Geospatial Information Authority of Japan Map.\footnote{
Digitizing the road network data in 1930 and 2005, we find that using these alternative road network data hardly affected the bilateral travel distance, suggesting that the road network did not experience a large change.
}





%%%%%%%%%%%%%%%%%%%%%
% Reduced-form
%%%%%%%%%%%%%%%%%%%%%


\section{Stylized Facts}\label{sec:reducedform}

We first document the cross-sectional relationship between the economic outcomes and the area share of the US bases in 1970, when the US-base sourced income accounted for around 30\% of the Gross National Income of Okinawa (Figure \ref{fig:okinawa_GDPshare}). 
Figure \ref{fig:reduced_crosssection1970} presents the scatterplot of each municipality, where the horizontal axis is the share of land occupied by the US bases and the vertical axis shows the changes in various economic outcomes from 1938 to 1970.
Here, we take the difference from 1938, the pre-WW2 period with no US base, to investigate economic development relative to the period without the US bases. That said, Appendix \ref{sec:facts_1970level} shows that our stylized facts also hold without taking the difference.
We also show a corresponding fitted line and report its slope, along with the standard error.


We measure urban concentration as population density per unit of land outside US bases. Figure \ref{fig:popdens_crosssection1970} relates base exposure to the log change in population density, while Figure \ref{fig:pop_crosssection1970} reports the corresponding change in population as a complementary measure of the number of residents. Population density adjusts for the amount of land available for civilian use: its postwar change reflects both population change and the reduction in available land caused by base occupation.\footnote{Prewar population density divides population by municipal area; postwar population density divides population by municipal area minus US-base area. } Both measures show that municipalities with greater base exposure experienced stronger urbanization.
The relationship is also economically strong: the estimated $\beta_t$ for population density in Figure \ref{fig:popdens_crosssection1970} is 3.045, implying that a 10 percentage-point increase in the US-base area share is associated with a roughly 36\% higher ratio of 1970 to 1938 population density.  Overall, we establish the following Fact \ref{result1}: 



\begin{figure}[]
    \subfloat[Population density\label{fig:popdens_crosssection1970}]{%
      \includegraphics[width=0.4\textwidth, page=5]{figure/loggrowth3870_scatter_areaUS.pdf}
    }
    \hfill
    \subfloat[Population level\label{fig:pop_crosssection1970}]{%
      \includegraphics[width=0.4\textwidth, page=4]{figure/loggrowth3870_scatter_areaUS.pdf}
    } \\  
    \subfloat[Service share\label{fig:serv_crosssection1970}]{%
      \includegraphics[width=0.4\textwidth, page=1]{figure/sharediff3870_scatter_areaUS.pdf}
    }
    \hfill
    \subfloat[Manufacturing share\label{fig:manu_crosssection1970}]{%
      \includegraphics[width=0.4\textwidth, page=2]{figure/sharediff3870_scatter_areaUS.pdf}
    }
    \hfill \\
    \subfloat[Agriculture share\label{fig:agri_crosssection1970}]{%
      \includegraphics[width=0.4\textwidth, page=3]{figure/sharediff3870_scatter_areaUS.pdf}
    }
    \vspace{0.2in}
    \caption{Relationship between US bases, Urbanization, and Industry Structure (1970)}
    \vspace{-0.2in}
    \floatfoot{\textbf{Note}: 
  We present a scatterplot of each municipality using the 1970 data. The horizontal axis is the US base area share. Each dot represents a municipality.
    As the vertical axis, we plot the values relative to those of 1938, a pre-war period without any US base.     
    Figure \ref{fig:popdens_crosssection1970} takes the change in the log population density, Figure \ref{fig:pop_crosssection1970} takes the change in the log population, Figure \ref{fig:serv_crosssection1970} takes the change in the service employment share, Figure \ref{fig:manu_crosssection1970} takes the change in the manufacturing employment share, and Figure \ref{fig:agri_crosssection1970} takes the change in the agriculture employment share. The linear fitted lines are also shown. 
    }
    \label{fig:reduced_crosssection1970}
\end{figure}

\begin{fact}[\textbf{Urbanization around the US bases}]\label{result1}
Municipalities with a higher share of land occupied by US military bases tend to have higher urbanization rates, whether urbanization is measured by population size or population density.
\end{fact}

Such rapid urbanization around the US bases accompanied significant growth of the service sector, but not manufacturing. Figure \ref{fig:serv_crosssection1970}  shows that  municipalities with higher US base shares experienced rapid growth in the service sector. The association is strong: an approximately 10 percentage-point increase in the US base share is associated with a 5 percentage-point increase in the service employment share. In stark contrast, Figure \ref{fig:manu_crosssection1970} shows no association between manufacturing employment share and the US base area share.\footnote{
Following the Japan Standard Industry Classification, we include the construction sector as manufacturing. We expect this tends to create a positive association between manufacturing and US base area share because the US military demands construction for the base operation. Despite this, we find little evidence of a positive association between these two. 
% NOTE: Oshiro and Sato (2021) RSUE treats construction sector separately, leading to a four-sector model
}
As a result, almost the entire decline in the agricultural employment share around the US bases, shown in Figure \ref{fig:agri_crosssection1970}, was absorbed by the increase in the service employment share. 
These patterns can be summarized by the following Fact \ref{result2}:

\begin{fact}[\textbf{Service-sector growth without industrialization around the US bases}]\label{result2}
Municipalities with a higher area share of the US bases were characterized by the rapid growth in the employment share of the service sector and the decline of the agriculture employment share. In contrast, the change in the manufacturing employment share had little association with the US base area share.
\end{fact}


Appendix \ref{appsec:reducedform} reports additional analyses of these stylized facts and their robustness. First, we find that the stylized facts also hold when the vertical axis is replaced with the corresponding levels in 1970, rather than the 1970--1938 difference.
Second, we consider the potential endogeneity of the US base locations by (i) controlling for topographic conditions and (ii) using the landing point and the occupation timing shown in Figure \ref{fig:okinawa_map}(a) as an instrumental variable, both of which do not change the main conclusion.
Third, we do not find statistically-significant evidence that pre-war trends in population density and sectoral employment shares drive our results. Population density exhibits some positive pre-trend but its magnitude is far from explaining the post-war growth in population density, Sectoral employment shares also show no statistically significant pre-trends.
Fourth, using the predetermined 1956 base share to address the potential endogeneity arising from base returns yields the same qualitative conclusions. 
Fifth, our results are robust to alternative measures of base importance based on the densities of Okinawan base workers or U.S. military personnel. 
Sixth, the association between U.S. bases and local economic outcomes remains qualitatively similar across different years, although the magnitude generally declines over time as the aggregate importance of the bases in the Okinawan economy diminishes.
Finally, supplementary sectoral evidence suggests that the positive employment effects are particularly pronounced in consumer-oriented service activities such as wholesale and retail and intangible services. 


These stylized facts document a strong relationship between US-base exposure and urbanization without industrialization. They do not, however, establish how much ongoing base-related expenditure and employment contribute to this pattern or reveal their aggregate income and welfare implications. Municipal comparisons may reflect both current base-related flows and the legacy of past investment, while migration, commuting, and shopping transmit effects across municipal boundaries. To address these questions, we develop a quantitative spatial model disciplined by GNI data on base-related purchases, labor income, and land rents. We evaluate a counterfactual that removes base-related expenditure, land-rent income, and employment while holding other local fundamentals fixed at their calibrated 1970 levels. This exercise quantifies how ongoing base-related consumption demand and employment shape urbanization and sectoral composition and affect aggregate income and welfare. The next section introduces the model.

%%%%%%%%%%%%%%%%%%%%%
% Model
%%%%%%%%%%%%%%%%%%%%%



\section{Model}
\label{sec:model}
We consider an economy (the Okinawa main island) consisting of several locations (municipalities). The total number of locations in this economy is $I$, and each location is indexed by $i=1,2,\cdots,I$. The economy as a whole is inhabited by a measure $L$ of workers, which is exogenously given as we assume a closed economy in the baseline model.\footnote{Section \ref{sec:counterfactual_external_validity} examines the sensitivity of the results to relaxing this fixed-population assumption.}
Workers inelastically supply a unit of labor and earn the wage. Each worker is allowed to reside in a location different from their workplace, implying that the mass of residents in location $i$ ($R_i$) may differ from that of workers working in location $i$ ($L_i$).

The present model has four sectors: agriculture ($A$), manufacturing ($M$), non-tradable services ($N$), and the US military base sector ($B$).
%Agricultural and manufacturing goods are produced in a perfectly competitive market. We assume that Okinawa is a small open economy and both goods are freely traded, indicating that the prices of agricultural and manufacturing goods are
The agricultural and manufacturing sectors are perfectly competitive and produce homogeneous goods. We assume that Okinawa is a small open economy and both goods are freely traded, indicating that the prices of agricultural and manufacturing goods are common across all locations in Okinawa. The non-tradable service sector is monopolistically competitive, implying that firms in this sector produce horizontally differentiated goods. These goods cannot be traded across locations and are consumed only locally, although we allow workers to make a trip to consume non-tradable services outside their residence. 
The US military base sector does not produce any good, but some labor force is required to support the operation of the US army stationed in Okinawa.

Each location is endowed with a fixed amount of land, $S_i$. However, some parts of the land may be occupied by the US military bases. Let $m_i \in [0,1]$ represent the share of land used as US military bases in location $i$. We posit that $m_i$ is determined for a military reason and is exogenously given to the model. Then, a fraction $1-m_i$ of the land is available for developing floor space that can be used for private purposes, including housing and the production of both tradable goods and non-tradable services. The floor space is supplied by perfectly competitive developers, and we let $H_i$ represent the amount of floor space in location $i$.

\subsection{Consumption}

\paragraph{Workers.}

Workers derive utility from consuming agricultural and manufacturing goods, non-tradable services and housing. We assume that the preference of a worker $\nu$ residing in location $i$ and working in sector $s \in \{ A,M,N,B \}$ in location $j$ is represented by a Cobb-Douglas function.\footnote{Preferences are homothetic: at given prices, expenditure shares do not vary with income, as in \cite{faber2019tourism} and \cite{fajgelbaum2022argentina}.} 
The indirect utility is given by
\begin{equation}
\label{eq:utility}
    U_{ijs}(\nu)=\frac{\varepsilon_{ijs}(\nu) w_{js}}{\tau_{ij}\mathbb{P}_{i}^\alpha Q_i^{1-\alpha}},
\end{equation}
where $\alpha \in (0,1)$ is the expenditure share on goods and services; $\mathbb{P}_{i}$ is the price index of goods and services in location $i$; $Q_i$ is the floor-space rent in location $i$; $w_{js}$ is the wage earned in sector $s$ in workplace $j$; and $\tau_{ij}$ is the utility cost of commuting from residence $i$ to workplace $j$. The term $\varepsilon_{ijs}(\nu)$ is an idiosyncratic preference shock for the residence--workplace--sector choice, which captures idiosyncratic reasons for living and working in particular locations and sectors. We assume that these shocks are independently drawn from a Fr\'{e}chet distribution, $\Pr\!\left(\varepsilon_{ijs}\leq x\right)=\exp\!\left(-A_{is}x^{-\theta}\right)$, where $A_{is}>0$ governs the amenity of residing in location $i$ and working in sector $s$, and $\theta>1$ is the (inverse) measure of the dispersion of idiosyncratic preferences. Workers observe the shocks before choosing where to live, where to work, and in which sector to work.
%$A_{i}$ controls the average amenities from residing in location $i$; and $\varepsilon_{ijs}(\nu)$ is an idiosyncratic preference shock for the pair of residence, workplace, and sector of employment, which captures idiosyncratic reasons for living and working in particular locations and sectors. We let $\varepsilon_{ijs}(\omega)$ be drawn from a Fr\'{e}chet distribution, $F(\varepsilon)=\exp(-\varepsilon^{-\theta})$, where $\theta>1$ is the (inverse) measure of the dispersion of idiosyncratic preferences. Workers are assumed to realize their own values of $\varepsilon_{ijs(\nu)}$ before choosing the locations to live and work and the sector to work.
%We assume that $A_{ij}$ is multiplicatively separable into a residence component ($A_i^R$) and a workplace component ($A_j^L$), implying that $A_{ij}=A_i^RA_j^L$ holds. 

The price index, $\mathbb{P}_{i}$, is defined as
\begin{equation}
\label{eq:price_index}
    \mathbb{P}_{i} \equiv
    \left[
    \left( P_{i}^T \right)^{1-\kappa}+
    \left( \mathbb{P}_{i}^N \right)^{1-\kappa}
    \right]^\frac{1}{1-\kappa},
\end{equation}
where $P_i^T$ is the price index of tradable goods (i.e., agricultural and manufacturing goods); $\mathbb{P}_{i}^N$ is the price index of non-tradable services faced by residents in location $i$;  and $\kappa>0$ is the elasticity of substitution between tradable goods and non-tradable services.\footnote{Demand weights are normalized to one because they are not separately identified from sectoral productivity; calibrated productivity therefore incorporates these weights. Appendix \ref{appsec:calibration_supplement} provides details.}

The price index of tradable goods, $P_i^T$, is defined over the prices of agricultural goods in location $i$ ($P_i^A$) and manufacturing goods in location $i$ ($P_i^M$) as follows: $P_i^T \equiv \Upsilon(P_i^A,P_i^M)$, where $\Upsilon(\cdot,\cdot )$ is a homothetic aggregator. Since the prices of agricultural and manufacturing goods are common across locations because of perfect competition and free trade, we can suppress the location subscript and write $P_i^A=P^A$ and $P_i^M=P^M$. It then follows that the price index of tradable goods is also common across locations: $P_i^T=P^T=\Upsilon(P^A,P^M)$. In our quantitative analysis, we normalize the price index of tradable goods to one, i.e., $P^T=1$.\footnote{
In Section \ref{sec:counterfactual_external_validity}, we endogenize the manufacturing goods prices, which requires us to specify $\Upsilon(P^A,P^M)$ to model the endogenous response of $P^T$. We assume that $\Upsilon(P^A,P^M)$ takes a Cobb-Douglas form in this extension.
}

We assume that workers can consume non-tradable services not only at their residence but also in other locations. That is, residents in location $i$ can make trips for shopping and enjoy non-tradable services provided in locations other than their residence. For simplicity, it is assumed that the set of destinations and the itinerary of shopping trips are identical across all residents in location $i$.
%That is, workers residing in location $i$ and working in location $j$ consume non-tradable services produced in both locations $i$ and $j$.
The non-tradable price index $\mathbb{P}_{i}^N$ takes the following form:
\begin{equation}
\label{eq:price_index_N}
    \mathbb{P}_{i}^N \equiv \left[
    \sum_j \xi_{ij} (P_j^N)^{1-\sigma}
    \right]^{\frac{1}{1-\sigma}}
\end{equation}
\begin{equation}
\label{eq:price_index_N_2}
    \mbox{with} \quad
    P_j^N \equiv \left[
    \int_0^{n_j} (p_j^N(\iota))^{1-\sigma} \mathrm{d}\iota
    \right]^{\frac{1}{1-\sigma}}
    ,
\end{equation}
where $P_j^N$ is the price index of non-tradable services in location $j$; $p_j^N(\iota)$ is the price of a non-tradable service indexed by $\iota$ and produced in location $j$; $\sigma>1$ is the elasticity of substitution among non-tradable varieties; $\xi_{ij} \in [0,1]$ satisfies $\sum_j \xi_{ij}=1$ and measures the extent to which residents in location $i$ spend on consumption activities in location $j$; and $n_j$ is the mass of varieties produced in location $j$. It is straightforward to assume that the farther away the locations $i$ and $j$, the smaller $\xi_{ij}$. Note that this formulation nests the standard case of purely local consumption of non-tradables by letting $\xi_{ii}=1$ and $\xi_{ij}=0$ for $j \neq i$.

By Roy's identity, the shares of expenditure on agricultural goods ($\chi_{i}^A$), manufacturing goods ($\chi_{i}^M$) and non-tradable services ($\chi_{i}^N$) are derived  as follows:
\begin{equation*}
\label{eq:share_A}
    \chi_{i}^A=\alpha \left( \frac{P^T}{\mathbb{P}_{i}} \right)^{1-\kappa} \frac{\partial \ln{\Upsilon(P^A,P^M)}}{\partial \ln{P^A}},
\end{equation*}
\begin{equation*}
\label{eq:share_M}
    \chi_{i}^M=\alpha \left( \frac{P^T}{\mathbb{P}_{i}} \right)^{1-\kappa} \frac{\partial \ln{\Upsilon(P^A,P^M)}}{\partial \ln{P^M}},
\end{equation*}
\begin{equation}
\label{eq:share_N}
    \chi_{i}^N=\alpha \left( \frac{\mathbb{P}_{i}^N}{\mathbb{P}_{i}} \right)^{1-\kappa}.
\end{equation}
The homotheticity of $\Upsilon$ implies that $\partial \ln{\Upsilon}/\partial \ln{P^A}+\partial \ln{\Upsilon}/\partial \ln{P^M}=1$ holds and that the share of expenditure on tradable goods is $\chi_{i}^A+\chi_{i}^M=\alpha \left( P^T/\mathbb{P}_{i} \right)^{1-\kappa}$. %In particular, if we let $\Upsilon$ be in the Cobb-Douglas form, both $\partial \ln{\Upsilon}/\partial \ln{P_i^A}$ and $\partial \ln{\Upsilon}/\partial \ln{P_i^M}$ become constant.
Note also that the share of expenditure on non-tradable services provided in location $j$ and spent by residents in location $i$ is $\xi_{ij} \chi_{i}^N \left( P_{j}^N/\mathbb{P}_{i}^N \right)^{1-\sigma}$.

\paragraph{Landowners.}
We assume that the land is owned by local landowners.
Unlike workers, local landowners neither supply labor to any sector nor commute to other locations. The total income of local landowners in location $i$ is the aggregate land rents from the land owned by them, which we denote by $\Omega_i$. Let $\mathbb{Q}_i^{US}$ represent the land rent for the US military bases in location $i$, which is exogenously given to the model. 
Note that $\mathbb{Q}_i^{US}$ can differ from the land rent in the private market, which we denote by $\mathbb{Q}_i$, because landowners were forced to lend the land even though $\mathbb{Q}_i^{US} < \mathbb{Q}_i$ holds, which is the empirically relevant case. 
Since location $i$'s total land endowment is $S_i$ and a fraction $m_i$ of it is occupied by the US bases, $\Omega_i$ is determined as follows: 
\begin{equation}
\label{eq:landowner_income}
    \Omega_i=[(1-m_i)\mathbb{Q}_i+m_i\mathbb{Q}_i^{US}]S_i.
\end{equation}

We assume that local landowners consume the same set of goods and services as workers except for housing and that the expenditure shares are also the same, implying that they spend a fraction $\chi_{i}^A/\alpha$ of income on agricultural goods, a fraction $\chi_{i}^M/\alpha$ on manufacturing goods, and a fraction $\chi_{i}^N/\alpha$ on non-tradable services.

\paragraph{The US military personnel.}
The US military bases are extensively distributed on Okinawa Island. Each base hosts some American military personnel for operation, and the mass of those working in the base in location $i$ is denoted by $L_i^{US}$, which is determined by military considerations and exogenously given to the model.

We assume that they are endowed with a fixed amount of housing in the base area where they work.\footnote{The US military personnel are required to live in housing within the US base area. As an exception, some are eligible to live outside bases, but they are still required to live in housing authorized by the military and such housing tends to be located around the bases (\citealt{tomochi09}).
} 
This implies that they do not account for housing when they make a consumption decision and that their residence and workplace are in the same location. We also assume that they consume both tradable goods and non-tradable services and that their expenditure shares are the same as those of workers. Therefore, they spend a fraction $\chi_{i}^A/\alpha$ of income on agricultural goods, a fraction $\chi_{i}^M/\alpha$ on manufacturing goods, and a fraction $\chi_{i}^N/\alpha$ on non-tradable services.
We denote by $e_{US}$ the individual expenditure of the US military personnel on Okinawa Island.\footnote{The US personnel may also receive income spent elsewhere (e.g., remittance to the US mainland). We do not consider such income to focus on the Okinawan economy.}


\subsection{Production}

In location $i$, sector $s \in \{ A,M,N \}$ hires a measure $L_{is}$ of workers and uses a measure $H_{is}$ of floor space as production factors. 
%Agricultural and manufacturing goods are produced in a perfectly competitive market. We assume that Okinawa is a small open economy and both goods are freely traded, indicating that the prices of agricultural and manufacturing goods are exogenously determined in the outside market, which we denote by $P^A$ and $P^M$, respectively. Since both goods are also assumed to be costlessly traded across locations within Okinawa, $P_i^A=P^A$ and $P_i^M=P^M$ hold for any $i$.

\paragraph{Agricultural sector.}
We assume that the agricultural sector's production technology is represented by a Cobb-Douglas function and exhibits constant returns to scale. Then, profit maximization and free entry imply that the unit cost of production is equal to the price of agricultural goods:
\begin{equation}
\label{eq:priceA}
    P^A=\frac{w_{iA}^{1-\gamma^A} Q_i^{\gamma^A}}{b_i^A},
\end{equation}
where $b_i^A$ is productivity of agricultural goods in location $i$ and $\gamma^A \in [0,1]$ is the agricultural sector's cost share of floor space. We do not consider agglomeration economies in the agricultural sector, indicating that $b_i^A$ is an exogenous parameter that reflects location $i$'s geographic conditions such as soil and terrain.

\paragraph{Manufacturing sector.}
As in the case of the agricultural sector, we assume that the manufacturing sector's production technology is given by a Cobb-Douglas function and exhibits constant returns to scale. Then, the unit cost of production is equal to the price of the manufacturing goods:
\begin{equation}
\label{eq:priceM}
    P^M=\frac{w_{iM}^{1-\gamma^M} Q_i^{\gamma^M}}{B_i^M},
\end{equation}
where $B_i^M$ is productivity of manufacturing goods in location $i$ and $\gamma^M \in [0,1]$ is the manufacturing sector's cost share of floor space.

To incorporate the agglomeration economies in the manufacturing sector, we assume that $B_i^M$ takes the following form:
\begin{equation}
\label{eq:productivityM}
    B_i^M=b_i^M L_{iM}^\eta,
\end{equation}
where $b_i^M$ is the exogenous production fundamental that captures the geographic features such as good access to rivers or the sea, and $\eta>0$ measures the strength of agglomeration forces.\footnote{We assume that there is no spillover from the US military personnel and Okinawans working for the US bases. We believe that no spillover is natural given that the military activities are largely unrelated to private economic activities. Indeed, Okinawan residents were not even permitted to enter the US bases except in special circumstances.
}

\paragraph{Non-tradable service sector.}
The non-tradable service sector is characterized by monopolistic competition, and each variety is produced by a single firm under increasing returns to scale. To start production, $f_i^N$ units of a Cobb-Douglas composite of labor and floor space are needed as a fixed input. Then, producing a unit of output requires $1/b_i^N$ units of the same Cobb-Douglas composite of labor and floor space, where $b_i^N$ is an exogenous parameter that measures productivity of non-tradable services in location $i$. We let $\gamma^N \in [0,1]$ represent the non-tradable service sector's cost share of floor space, implying that the total costs of producing $x_i^N(\iota)$ units of output of variety $\iota$ are given by
\begin{equation}
\label{eq:costN}
    c_i^N(\iota)=\left( f_i^N+\frac{x_i^N(\iota)}{b_i^N} \right) w_{iN}^{1-\gamma^N} Q_i^{\gamma^N}.
\end{equation}
Since consumers' preferences are characterized by the constant elasticity of substitution among varieties, the profit-maximizing price set by a firm in location $i$ exhibits the constant-markup property:
\begin{equation}
\label{eq:markup}
    p_i^N=\frac{\sigma}{\sigma-1} \frac{w_{iN}^{1-\gamma^N} Q_i^{\gamma^N}}{b_i^N},
\end{equation}
where we omit the firm index $\iota$ because there is no longer heterogeneity across firms in the same location. Thus, the zero profit condition implies that the equilibrium output for each firm in location $i$ is $x_i^N=(\sigma-1)f_i^Nb_i^N$. Plugging this into (\ref{eq:costN}), the total costs incurred by a firm in location $i$ can be calculated as $\sigma f_i^N w_{iN}^{1-\gamma^N} Q_i^{\gamma^N}$. Because each variety is produced by a single firm, the total mass of firms producing non-tradable services in location $i$ is represented by $n_i$. The aggregate payments to labor and floor space are $w_{iN}L_{iN}=\sigma(1-\gamma^N)f_i^N n_i w_{iN}^{1-\gamma^N} Q_i^{\gamma^N}$ and $Q_iH_{iN}=\sigma\gamma^N f_i^N n_i w_{iN}^{1-\gamma^N} Q_i^{\gamma^N}$, respectively. From these two equations, the equilibrium mass of firms in the non-tradable service sector in location $i$ can be expressed as follows:
\begin{equation}
\label{eq:variety}
    n_i=\frac{1}{\sigma f_i^N} \left( \frac{L_{iN}}{1-\gamma^N} \right)^{1-\gamma^N} \left( \frac{H_{iN}}{\gamma^N} \right)^{\gamma^N}.
\end{equation}

Substituting (\ref{eq:markup}) and (\ref{eq:variety}) into (\ref{eq:price_index_N_2}), the price index of non-tradable services in location $i$ can be written as
\begin{equation}
\label{eq:priceN}
    P_{i}^N=\frac{w_{iN}^{1-\gamma^N}Q_i^{\gamma^N}}{B_i^N},
\end{equation}
where $B_i^N$ is given by
\begin{equation}
\label{eq:productivityN}
    B_i^N=\widetilde{b}_i^N L_{iN}^{\frac{1-\gamma^N}{\sigma-1}} H_{iN}^{\frac{\gamma^N}{\sigma-1}},
\end{equation}
and $\widetilde{b}_i^N$ is defined as $\widetilde{b}_i^N \equiv (\sigma-1)\left[ \sigma^\sigma (1-\gamma^N)^{1-\gamma^N}(\gamma^N)^{\gamma^N} f_i^N \right]^{1/(1-\sigma)} b_i^N$. Note that $B_i^N$ increases in $L_{iN}$ because $(1-\gamma^N)/(\sigma-1)>0$ holds, implying that the non-tradable service sector also exhibits agglomeration economies. 
Finally, substituting equation (\ref{eq:priceN}) into equation (\ref{eq:price_index_N}) gives the price index for workers living in $i$ as follows:
\begin{equation}
\label{eq:price_index_N_3}
    \mathbb{P}_{i}^N=\left[
    \sum_j \xi_{ij} \left( \frac{w_{jN}^{1-\gamma^N}Q_j^{\gamma^N}}{B_j^N} \right)^{1-\sigma}
    \right]^{\frac{1}{1-\sigma}}.
\end{equation}

\paragraph{The US military base sector.}
Although the US military base sector does not produce any good transacted in the market, it needs to employ some Okinawan people to support the operation of the US army.\footnote{
We can consider that as in \cite{becker2021impact}, US bases produce global public goods (i.e., military forces), which is irrelevant for location choices within the Okinawa island.
}
Note the distinction between the Okinawan people working in the US bases and the US military personnel, who are Americans.
The mass of Okinawan workers demanded by the base in location $i$ is represented by $L_{iB}$, which is exogenously given to the model because it is determined based on military considerations.

We denote the set of locations hosting the US bases by $\mathcal{I}_B$ and let the number of such locations be $I_B$ (i.e., $|\mathcal{I}_B|=I_B$). For locations included in $\mathcal{I}_B$, the wage in the base in location $i$, $w_{iB}$, is endogenously determined such that the labor supply matches the demand in each location. For locations without bases, we let $w_{iB}=0$ hold because there is no demand for labor in the base sector in such locations.

\subsection{Floor space supply}

Floor space is supplied by a perfectly competitive construction sector that uses land and tradable goods as inputs. Following \cite{epple2010new} and \cite{combes2021production}, we assume that the production function for the floor space takes a Cobb-Douglas form and the corresponding cost function is given by $C_i(H_i)=\varkappa \left( P^T \right)^\mu \mathbb{Q}_i^{1-\mu} H_i$, where $\mu \in [0,1]$ is the cost share of tradable inputs and $\varkappa$ is a positive constant that governs the level of marginal cost.\footnote{
\cite{yoshida18} and \cite{KII2022} show evidence that the Cobb-Douglas housing production function is also a good approximation of reality in Japan.
}

Although the land endowment of location $i$ is $S_i$, a fraction $m_i$ is occupied by the US bases, indicating that the amount of land available for development is $(1-m_i)S_i$. Thus, Shephard's lemma yields
\begin{equation}
    (1-m_i)\mathbb{Q}_iS_i=(1-\mu)Q_iH_i.
    \label{eq:fp_lp}
\end{equation}
This equation allows us to compute the land rents, given the floor space rents and the total floor space supply.

Note also that perfect competition implies $Q_i=\varkappa \left( P^T \right)^\mu \mathbb{Q}_i^{1-\mu}$. Substituting (\ref{eq:fp_lp}) into this equation and arranging the terms, we obtain the following floor space supply function:
\begin{equation*}
    H_i=(1-m_i)S_i\left( \frac{Q_i}{P^T} \right)^\phi ,
\end{equation*}
where $\phi \equiv \mu/(1-\mu)$ represents the floor space supply elasticity.\footnote{
Here, without loss of generality, we simplify notations by setting the unit of floor space such that $\varkappa=(1-\mu)^{-(1-\mu)}$ holds.}


\subsection{Market equilibrium}

\paragraph{Labor market.}
Each worker chooses the residence--workplace--sector combination that maximizes utility in (\ref{eq:utility}). Since $\varepsilon_{ijs}(\nu)$ follows a Fr\'{e}chet distribution, the probability that a worker chooses to live in location $i$ and work in sector $s$ in location $j$ is
\begin{equation}
\label{eq:prob_live_work}
    \lambda_{ijs}=
    \frac{A_{is}\left( w_{js} \mathbb{P}_{i}^{-\alpha} Q_i^{-(1-\alpha)} \right)^\theta \tau_{ij}^{-\theta}}{\sum_k \sum_l \sum_r A_{kr}\left( w_{lr} \mathbb{P}_{k}^{-\alpha} Q_k^{-(1-\alpha)} \right)^\theta \tau_{kl}^{-\theta}}.
\end{equation}
Summing these probabilities across all workplaces for a given pair of residence $i$ and sector $s$, we obtain the probability that a worker chooses to live in location $i$ and work in sector $s$:
\begin{equation}
\label{eq:prob_live}
    \lambda_{is}^R \equiv \sum_j \lambda_{ijs}.
\end{equation}
Therefore, the probability that a worker commutes to location $j$ conditional on living in location $i$ and working in sector $s$ is
\begin{equation}
\label{eq:prob_commute}
    \lambda_{ij \mid is} \equiv \frac{\lambda_{ijs}}{\lambda_{is}^R}=
    \frac{w_{js}^\theta \tau_{ij}^{-\theta}}{\sum_l w_{ls}^\theta \tau_{il}^{-\theta}}.
\end{equation}

Let $R_{is}$ represent the mass of workers who live in location $i$ and work in sector $s$. By equation (\ref{eq:prob_live}), $R_{is}$ is determined as follows:
\begin{equation}
\label{eq:residents}
    R_{is}=\lambda_{is}^R L.
\end{equation}
Using equation (\ref{eq:prob_commute}), we obtain the mass of workers in location $i$ for each sector:
\begin{equation}
\label{eq:workers}
\begin{split}
    & L_{is}=\sum_k\lambda_{ki \mid ks}R_{ks}
    \quad \mbox{for all $i$ and $s \in \{ A,M,N \}$} ,\\
    &\mbox{and}\quad
    L_{iB}=\sum_k\lambda_{ki \mid kB}R_{kB}
    \quad \mbox{for $i \in \mathcal{I}_B$}.
\end{split}
\end{equation}
Note that the mass of US-base workers in location $i$, $L_{iB}$, is exogenously given and locations without bases exhibit $L_{iB}=0$. The masses of workers living and employed in location $i$ are, respectively, expressed as
\begin{equation}
\label{eq:total_pop}
    R_i=\sum_s R_{is} \quad \mbox{and} \quad L_i=\sum_s L_{is}.
\end{equation}

We now turn to describing wages.
We can express private-sector wages by using (\ref{eq:priceA}), (\ref{eq:priceM}), (\ref{eq:productivityM}), (\ref{eq:priceN}), (\ref{eq:productivityN}) as follows:
\begin{equation}
\label{eq:wageA}
    w_{iA}=(\widetilde{b}_i^A)^{\frac{1}{1-\gamma^A}} Q_i^{-\frac{\gamma^A}{1-\gamma^A}},
\end{equation}
\begin{equation}
\label{eq:wageM}
    w_{iM}=(\widetilde{b}_i^M)^{\frac{1}{1-\gamma^M}} Q_i^{-\frac{\gamma^M}{1-\gamma^M}} L_{iM}^{\frac{\eta}{1-\gamma^M}},
\end{equation}
\begin{equation}
\label{eq:wageN}
    w_{iN}=(\widetilde{b}_i^N)^{\frac{1}{1-\gamma^N}} (P_i^N)^{\frac{1}{1-\gamma^N}} Q_i^{-\frac{\gamma^N}{1-\gamma^N}} L_{iN}^{\frac{1}{\sigma-1}} H_{iN}^{\frac{\gamma^N}{(\sigma-1)(1-\gamma^N)}},
\end{equation}
where $\widetilde{b}_i^A \equiv P^Ab_i^A$ and $\widetilde{b}_i^M \equiv P^Mb_i^M$ are the price-adjusted productivities of agricultural and manufacturing goods in location $i$, respectively.\footnote{
We cannot identify $P^A$ and $b_i^A$ (resp. $P^M$ and $b_i^M$) separately in our calibration process. Instead, we recover the value of $\widetilde{b}_i^A$ (resp. $\widetilde{b}_i^M$) such that it rationalizes the observed equilibrium.
}
In locations with US bases, the wage rate of US-base workers is determined so that it solves the second equation of (\ref{eq:workers}).


\paragraph{Floor space market.}
Since equation (\ref{eq:utility}) implies that the share of expenditure on housing is $1-\alpha$, the total amount of floor space demanded by residents in location $i$, denoted by $H_{iR}$, is given by
\begin{equation}
\label{eq:land_house}
    H_{iR}=(1-\alpha)\sum_j \sum_s \frac{\lambda_{ij \mid is}w_{js}R_{is}}{Q_i}.
\end{equation}
Private uses of floor space include not only housing but also the production of both tradable goods and non-tradable services. Because each sector exhibits a Cobb-Douglas production technology, the amount of floor space used for production in sector $s \in \{ A,M,N \}$ can be written as follows:
\begin{equation}
\label{eq:land_production}
    H_{is}=\frac{\gamma^s w_{is}L_{is}}{(1-\gamma^s)Q_i}.
\end{equation}

In location $i$, the equilibrium of the floor space market is characterized by $H_i=H_{iR}+\sum_{s \in \{ A,M,N \}} H_{is}$. Combining this equation with (\ref{eq:land_house}) and (\ref{eq:land_production}), we obtain the following condition for floor space market clearing:
\begin{equation}
\label{eq:land_market}
    Q_iH_i=(1-\alpha)\sum_j\sum_{s \in \{ A,M,N,B \}} \lambda_{ij \mid is}w_{js}R_{is}+\sum_{s \in \{ A,M,N \}}\frac{\gamma^s w_{is}L_{is}}{1-\gamma^s},
\end{equation}
where the first and second terms on the right-hand side represent the aggregate rents from housing and production, respectively.
%\footnote{Note that the Okinawan residents working for US bases consume floor space for their housing, which is the reason why the set for the summation index $s$ contains $B$ in the first term of the right-hand side of (\ref{eq:land_market}), while it does not in the second term.} 
%Note that we can rewrite equation (\ref{eq:land_market}) as a function of land rents, $\mathbb{Q}_i$, by using equation (\ref{eq:fp_lp}).

\paragraph{Non-tradable service market.}
In each location, the total revenue of firms producing non-tradable services must be equal to the total expenditure on non-tradable services spent by workers living or working there, local landowners, and the US military personnel in that location.

Workers living in location $i$ spend a fraction $\xi_{ij} \chi_{i}^N (P_{j}^N/\mathbb{P}_{i}^N)^{1-\sigma}$ of their income on non-tradable services served in location $j$. Multiplying this expenditure share by $1/\alpha$, we obtain the shares of expenditure spent by local landowners and the US military personnel in location $i$ on non-tradable services produced in location $j$. Thus, the market clearing condition for non-tradable services in location $i$ can be expressed as follows: 
\begin{equation}
\label{eq:goods_market_pre}
\begin{split}
    \left( \frac{w_{iN}L_{iN}}{1-\gamma^N} \right)^{1-\gamma^N} \left( \frac{Q_i H_{iN}}{\gamma^N} \right)^{\gamma^N}=
    & \sum_k \sum_j \sum_s
    \xi_{ki}\chi_{k}^N\left( \frac{P_i^N}{\mathbb{P}_k^N} \right)^{1-\sigma} \lambda_{kj \mid ks}w_{js}R_{ks}\\
    & +\sum_k \frac{\xi_{ki}\chi_{k}^N}{\alpha}\left( \frac{P_i^N}{\mathbb{P}_k^N} \right)^{1-\sigma} \left( \Omega_k+e_{US}L_k^{US} \right),
\end{split}
\end{equation}
where the left-hand-side captures the total revenue of non-tradable firms and the first and second terms on the right-hand side represent the total expenditure on non-tradable services in location $i$ from all workers and that from local landowners and the US military personnel, respectively.\footnote{Note that the total costs incurred by all non-tradable firms in location $i$ are expressed as $\sigma f_i^N n_i w_{iN}^{1-\gamma^N} Q_i^{\gamma^N}$. As the present model features free entry, the total revenue equals the total costs. Then, using (\ref{eq:variety}), the total revenue of firms in location $i$ can be written as $\left[ w_{iN}L_{iN}/(1-\gamma^N) \right]^{1-\gamma^N} (Q_iH_{iN}/\gamma^N)^{\gamma^N}$.

}

From (\ref{eq:landowner_income}), (\ref{eq:fp_lp}), (\ref{eq:land_production}) and (\ref{eq:land_market}), the market clearing condition (\ref{eq:goods_market_pre}) can be transformed into
\begin{equation}
\label{eq:goods_market}
\begin{split}
    \frac{w_{iN}L_{iN}}{1-\gamma^N}
    =&\frac{1-\mu+\alpha\mu}{\alpha}\sum_k \sum_j \sum_{s \in \{ A,M,N,B \}} \xi_{ki}\chi_{k}^N\left( \frac{P_i^N}{\mathbb{P}_k^N} \right)^{1-\sigma} \lambda_{kj \mid ks}w_{js}R_{ks}\\
    &+\frac{1-\mu}{\alpha}\sum_k \sum_{s \in \{ A,M,N \}} \xi_{ki}\chi_{k}^N\left( \frac{P_i^N}{\mathbb{P}_k^N} \right)^{1-\sigma} \frac{\gamma^sw_{ks}L_{ks}}{1-\gamma^s}
    +\frac{1}{\alpha}\sum_k \xi_{ki}\chi_{k}^N\left( \frac{P_i^N}{\mathbb{P}_k^N} \right)^{1-\sigma} \Phi_k,
\end{split}
\end{equation}
where $\Phi_k \equiv m_k\mathbb{Q}_k^{US}S_k+e_{US}L_k^{US}$ is the sum of the consumption of the US military personnel in Okinawa and the land rents of the US bases. Note that the land rents of the US bases are paid to local landowners, but they are assumed to have the same spending structure as the US military personnel. Therefore, ultimately, the land rent of the US bases is also considered as the consumption coming from the US bases. Thus, we interpret $\Phi_k$ as the US-base sourced consumption in municipality $k$.


\paragraph{Workers' welfare.}

Prior to the realization of the idiosyncratic preference shocks, expected utility from living in Okinawa is
\begin{equation}
\label{eq:welfare}
    \mathbb{W} = \Gamma\left( \frac{\theta-1}{\theta} \right)
    \left[ \sum_i \sum_j \sum_s A_{is}\left( w_{js} \mathbb{P}_{i}^{-\alpha} Q_i^{-(1-\alpha)} \right)^\theta \tau_{ij}^{-\theta} \right]^{\frac{1}{\theta}},
\end{equation}
where $\Gamma(z)=\int_0^\infty t^{z-1}\mathrm{e}^{-t} \mathrm{d}t$ is the Gamma function. In our baseline model, we assume a closed economy so that the total population $L$ is exogenously determined. In Section \ref{sec:counterfactual_external_validity}, we consider a small-open economy situation in which workers choose whether to live in 
 Okinawa, which gives the expected utility $\mathbb{W}$, or live outside Okinawa. 

\paragraph{Equilibrium conditions.}
The present model has $19 \times I+I_B$ endogenous variables: the spatial distribution of residents and employment, $\{ R_i,L_i \}_{i=1}^I$; the masses of residents and employment by sector, $\{ R_{is} \}_{i=1}^I$ for $s \in \{ A,M,N,B \}$ and $\{ L_{is} \}_{i=1}^I$ for $s \in \{ A,M,N \}$; the total amount of floor space used for housing, $\{ H_{iR} \}_{i=1}^I$; the total amount of floor space used for production in each sector, $\{ H_{is} \}_{i=1}^I$ for $s \in \{ A,M,N \}$; the wages in private sectors, $\{ w_{is} \}_{i=1}^I$ for $s \in \{ A,M,N \}$; the wages in the US base sector, $\{ w_{iB} \}_{i \in \mathcal{I}_B}$; the floor space rent, $\{ Q_i \}_{i=1}^I$;  the land rent, $\{ \mathbb{Q}_i \}_{i=1}^I$; and the price index of non-tradable services in location $i$, $\{ P_{i}^N \}_{i=1}^I$. These variables are determined in equations (\ref{eq:fp_lp}), (\ref{eq:residents}), (\ref{eq:workers}), (\ref{eq:total_pop}), (\ref{eq:wageA}), (\ref{eq:wageM}), (\ref{eq:wageN}), (\ref{eq:land_house}), (\ref{eq:land_production}), (\ref{eq:land_market}), and (\ref{eq:goods_market}).


\section{Calibration}\label{sec:calibration}

In this section, we calibrate the model parameters to match the Okinawan economy around 1970. 
We choose 1970 as our target of calibration in light of the substantial fraction of Okinawa's GNI in 1970 (see Figure \ref{fig:okinawa_GDPshare}) and rich data availability.
The calibration is done in steps, which we describe below. Appendix \ref{appsec:calibration_supplement} provides supplementary details on calibration procedure. For convenience, Table \ref{tab:para_value} summarizes calibrated parameter values and their sources.

\begin{table}[t]
    \centering
    \caption{Values of structural parameters in the model calibration}\label{tab:para_value}
    \adjustbox{width=\textwidth}{
    \begin{tabular}{llll}
        \toprule\toprule
        Parameter & Value & Description & Source \\
        \midrule
        $1-\alpha$ & 0.25 & Housing expenditure share & \citet{davis2011household} \\
        %$\beta$ & xxx & Strength of preference for non-tradable services & Own estimation \\
        $\sigma$ & 4.6 & Elasticity of substitution among non-tradable varieties & \citet{miyauchi2022economics} \\
        $\kappa$ & 1.87 & Elasticity of substitution between tradable and non-tradable goods & Own estimation \\
        $\theta$ & 2.19 & Fr\'{e}chet shape parameter & \cite{hayakawa2021highspeed} \\
        $\zeta$ & 0.845 & Commuting cost elasticity & Own estimation \\
        $\psi_S$ & 2.3 & Travel cost elasticity for shopping & Own estimation \\
        $\gamma^A$ & 0.3 & Cost share of floor space in the agricultural sector & \citet{hayashi2008japan} \\
        $\gamma^M$ & 0.2 & Cost share of floor space in the manufacturing sector & \cite{ahlfeldt2015economics} \\
        $\gamma^N$ & 0.2 & Cost share of floor space in the non-tradable service sector & \cite{ahlfeldt2015economics} \\
        $\eta$ & 0.05 & Agglomeration elasticity in the manufacturing sector & \citet{rosenthal2004agglomeration} \\
        $\phi$ & 2 & Floor space supply elasticity & \citet{heblich2020making} \\
        % $m_i$ & - & Land share of the US bases of each location & Okinawa Statistical Yearbook \\
        % $\Phi_i$ & - & US-base income of each location & Okinawa Statistical Yearbook \\
        % $A_i$& - & Exogenous residential amenities of each location & Own estimation \\
        % $\widetilde{b}_i^s$& - & Sector-specific adjusted productivity of each location & Own estimation \\
\bottomrule\bottomrule
    \end{tabular} 
}
\end{table}

\paragraph{Step 1: Calibrating parameters from external information.}

We first determine some parameter values based on external studies. The housing expenditure share ($1-\alpha$) is set at 0.25 based on \citet{davis2011household}.
The elasticity of substitution among non-tradable varieties ($\sigma$) is set at 4.6 based on \citet{miyauchi2022economics}, which is also consistent with \citet{hobijn2019sticker}. 
We set the Fr\'{e}chet shape parameter $\theta=2.19$, following \citeauthor{hayakawa2021highspeed}'s (\citeyear{hayakawa2021highspeed}) estimate using across-municipality Japanese data.\footnote{
We also experimented with estimating $\theta$ to match the observed dispersion of residential income across municipalities, analogously to \cite{ahlfeldt2015economics}. 
While this procedure has limitations in our context due to the relatively small number of municipalities, it yields a value close to $\theta=2.19$.
}
The cost share of floor space in the agricultural sector ($\gamma^A$) is set at 0.3 following \cite{hayashi2008japan}. The cost share of floor space in the manufacturing sector and the non-tradable service sector ($\gamma^M, \gamma^N$) is set at 0.2 following \citet{ahlfeldt2015economics}. 
The agglomeration elasticity in manufacturing productivity ($\eta$) in each municipality is set at 0.05, which is consistent with \cite{rosenthal2004agglomeration} and \cite{ahlfeldt2019economic}. Finally, the floor-space supply elasticity is set at 2 following \cite{heblich2020making}, which is consistent with Japanese floor-space supply technology in land-scarce environments like Okinawa (\citealt{yoshida18}). 


\paragraph{Step 2: Gravity equation estimation.}

We estimate the commuting cost parameter using commuting flows observed separately for workers in agriculture, manufacturing, and non-tradable services. We assume that the commuting cost is written as $\tau_{ij} = d_{ij}^\zeta $, where $d_{ij}$ is the distance between residence $i$ and workplace $j$ and $\zeta$ is the commuting cost elasticity with respect to commuting distance. Then, taking the log of equation (\ref{eq:prob_commute}) gives the sector-conditioned gravity equation
\begin{equation}
\label{eq:commuting_gravity_sector}
    \ln\lambda_{ij\mid is}
    =
    -\vartheta\ln d_{ij}
    +o_{is}^{R}+o_{js}^{L} ,
\end{equation}
where $\vartheta$ is defined as $\vartheta \equiv \zeta\theta$, and $o_{is}^{R}$ and $o_{js}^{L}$ are residence-by-sector and workplace-by-sector fixed effects, respectively.
%The former absorbs the sector-specific commuting-market-access denominator, while the latter absorbs the transformed final wage, $\ln\omega_{js}=\theta\ln w_{js}$, at workplace $j$ in sector $s$.

%We estimate the commuting cost parameter by estimating the commuting gravity equation.
%Summing (\ref{eq:prob_live_work}) across all sectors, we obtain the following commuting gravity equation:
%\begin{equation}
%\label{eq:gravity}
%    \lambda_{ij}=
%    \frac{\left( A_{i} \widetilde{w}_j \mathbb{P}_{i}^{-\alpha} Q_i^{-(1-\alpha)} \right)^\theta \tau_{ij}^{-\theta}}{\sum_k \sum_l \left( A_{k} \widetilde{w}_l \mathbb{P}_{k}^{-\alpha} Q_k^{-(1-\alpha)} \right)^\theta \tau_{kl}^{-\theta}},
%\end{equation}
%where $\lambda_{ij}$ is the share of workers living in $i$ and working in $j$. 
%We assume that the commuting cost is written as $\tau_{ij} = d_{ij}^\zeta $, where $d_{ij}$ is the commuting distance from residence $i$ to workplace $j$ and $\zeta$ is the commuting cost elasticity with respect to commuting distance. 
%Then, the gravity equation implied by (\ref{eq:gravity}) is rewritten as
%\begin{equation}
%\label{eq:log_gravity}
%    \ln{\lambda_{ij}}=
%    -\zeta \theta \ln d_{ij}
%    +o_i^R+o_j^L+\epsilon_{ij},
%\end{equation}
%where $o_i^R$ is the residence fixed effect, $o_j^L$ is the workplace fixed effect, and $\epsilon_{ij}$ is the measurement error. Using the data on intermunicipal commuting flows, we estimate the gravity equation (\ref{eq:log_gravity}) and obtain the estimate of $\zeta$.

%For estimation, let $F_{ijs}=R_{is}\lambda_{ij\mid is}$ denote the corresponding sector-specific commuting flow. Multiplication by $R_{is}$ only adds $\ln R_{is}$ to the residence-by-sector fixed effect, leaving the distance coefficient unchanged. 
Using trip survey data, we observe $\lambda_{ij|is}$ and estimate the equation \eqref{eq:commuting_gravity_sector} by PPML and OLS. The estimated distance coefficient identifies $\vartheta=\zeta\theta$; given the externally calibrated value of $\theta$, we recover the commuting cost elasticity as $\zeta=\vartheta/\theta$. Table \ref{tab:gravity_commutingDT} presents estimation results, which yield our preferred estimate of 1.85. Since we have set $\theta=2.19$, this implies $\zeta=1.85/2.19\simeq0.845$.

We next turn to estimating the travel cost for shopping using a gravity specification. The distance-based shopping-trip shares used in the model are given by
\begin{equation}
\label{eq:shoppingprob}
    \xi_{ij}^{\mathrm{trip}} = \frac{d_{ij}^{-\psi_S}}{ \sum_k d_{ik}^{-\psi_S}} ,
\end{equation}
where $d_{ij}$ is the travel distance between municipalities $i$ and $j$, and $\psi_S$ is the parameter governing travel costs.
Appendix \ref{appsec:gravity} reports shopping gravity estimates using destination shares, and we set $\psi_S=2.3$ based on Column 1 of Table \ref{tab:gravity_shopping}. We use this estimate as a calibration benchmark and construct the model's shopping weights $\xi_{ij}$ from the distance-based shares in equation (\ref{eq:shoppingprob}).

\paragraph{Step 3: Wages and rents.}

By combining (\ref{eq:prob_commute}) and (\ref{eq:workers}), we obtain the following labor market clearing condition for each sector:
\begin{equation}
\label{eq:LMC}
    L_{is}=\sum_k \frac{\omega_{is} d_{ki}^{-\vartheta}}{\sum_l \omega_{ls} d_{kl}^{-\vartheta}} R_{ks},
\end{equation}
where $\omega_{is} \equiv w_{is}^\theta$ is the transformed wage of sector $s$ in location $i$. Note also that the probability that a worker resides in location $i$ conditional on working in sector $s$ in location $j$ is given by $\lambda_{ijs}/\sum_k \lambda_{kjs}=\widetilde{A}_{is}d_{ij}^{-\vartheta}/\sum_k \widetilde{A}_{ks}d_{kj}^{-\vartheta}$, where $\widetilde{A}_{is} \equiv A_{is}\left(\mathbb{P}_{i}^{\alpha}Q_i^{1-\alpha}\right)^{-\theta}$ captures the composite attractiveness of living in $i$ and working in sector $s$, inclusive of amenities, goods prices, and housing rents. Then, $\{R_{is}\}_{i=1}^I$ and $\{L_{is}\}_{i=1}^I$ satisfy the following system of equations:
\begin{equation}
\label{eq:LMC_2}
    R_{is}=\sum_j \frac{\widetilde{A}_{is}d_{ij}^{-\vartheta}}{\sum_l \widetilde{A}_{ls}d_{lj}^{-\vartheta}}L_{js} .
\end{equation}

Given the estimate of $\vartheta$ obtained in the previous step and the values of $L_{is}$ and $R_{is}$ observed in the data, we solve equation (\ref{eq:LMC}) for $\omega_{is}$ and equation  (\ref{eq:LMC_2}) for $\widetilde{A}_{is}$. Note that equations (\ref{eq:LMC}) and (\ref{eq:LMC_2}) are homogeneous of degree zero in $\omega_{is}$ and $\widetilde{A}_{is}$, respectively, implying that we need to impose normalization conditions to determine the absolute levels of both variables. For $\widetilde{A}_{is}$ ($s \in \{ A,M,N \}$), we normalize the geometric means of $\widetilde{A}_{is}$ across all locations to one, which implicitly assumes that the scale of the amenities $A_{is}$ is common across all private sectors $s \in \{ A,M,N \}$. For the US-base sector, we allow the geometric mean of $\widetilde{A}_{iB}$ to differ from one and denote it by $a_B$, reflecting that the US-base sector may have specific job (dis)amenities compared with the private sectors.\footnote{
For example, military-base workers were subject to a distinctive and restrictive labor-law regime, such as the restrictions on union activities  (\citealt{Toya2019LabourRelationsCommission}).  \cite{Fukuyama2023OkinawanMaids} also documents the language and cultural barriers faced by the Okinawan maids employed by US military personnel.
}
The introduction of the US-base job amenity $a_B$ allows us to match the observed GNI share of the US-base sourced labor income.  
For the values of $\omega_{is}$, we normalize the geometric mean of $\omega_{iA}$ to one by suitably choosing the unit of labor. Then, the scales of $\omega_{iM}$, $\omega_{iN}$, and $\omega_{iB}$ are chosen to match the employment share of each sector (see Appendix \ref{appsec:calibration_supplement}).

Based on the normalization described above, we first back out $\widetilde{A}_{is}$ for $s \in \{ A,M,N \}$ and then $\omega_{is}$ for $s \in \{ A,M,N \}$. Given the value of $\theta$, we can readily recover the wages in sectors $A$, $M$, and $N$ by using $w_{is} = \omega_{is}^{1/\theta}$.
For the US-base sector, we first back out  $\widetilde{A}_{iB}^\prime \equiv \widetilde{A}_{iB}/a_B$, which is the amenity level renormalized to have the geometric mean of one by netting out the US-base job amenity. We then calibrate $a_B$ so that the model-implied GNI share of wages earned by workers in the US-base sector matches the observed value, which we describe in more detail in Appendix \ref{appsec:calibration_GNI}. 
The knowledge of $a_B$, together with the observed employment share of the US bases, also pins down the scale of $\omega_{iB}$ and hence $w_{iB}$. 

% For the US-base sector, we cannot directly back out the values of $\omega_{iB}$ and $\widetilde{A}_{iB}$ because of the additional unknown parameter $a_B$. Instead, we define the auxiliary variables $\omega_{iB}^\prime \equiv a_B\omega_{iB}$ and $\widetilde{A}_{iB}^\prime \equiv \widetilde{A}_{iB}/a_B$ and recover these variables. This transformation allows us to compute the values of $\omega_{iB}^\prime$ and $\widetilde{A}_{iB}^\prime$ in the same manner as for the other sectors, because they also satisfy equations (\ref{eq:LMC}) and (\ref{eq:LMC_2}) and the geometric mean of $\widetilde{A}_{iB}^\prime$ is one.\footnote{
% Because $\widetilde{A}_{iB}\omega_{iB}=\widetilde{A}_{iB}^\prime\omega_{iB}^\prime$ holds, we can also determine the scale of $\omega_{iB}^\prime$ analogously to the other sectors.
% }
% Then, we calibrate $a_B$ so that the model-implied GNI share of wages earned by workers in the US-base sector matches the observed value, which we describe in more detail in Appendix \ref{appsec:calibration_GNI}. Once $\omega_{iB}^\prime$, $\widetilde{A}_{iB}^\prime$, and $a_B$ have been obtained, we recover $w_{iB}$ and $\widetilde{A}_{iB}$ using $w_{iB}=(\omega_{iB}^\prime/a_B)^{1/\theta}$ and $\widetilde{A}_{iB}=a_B\widetilde{A}_{iB}^\prime$.

%Using the estimate of $\vartheta$ obtained in the previous step and the values of $L_{is}$ and $R_{is}$ observed from the data, we can solve (\ref{eq:LMC}) for $\omega_{is}$, up to scale.
%Without loss of generality, we normalize the geometric mean of $\omega_{iA}$ to one by suitably choosing the unit of labor. The scale of $\omega_{iM}, \omega_{iN}, \omega_{iB}$ is chosen to match the employment share of each sector.\footnote{More specifically, we determine the scale of $\omega_{is}$ ($s = M, N, B$) to satisfy $\lambda_s / \lambda_A = L_s / L_A$, that is, 
%    \begin{equation}
%    \label{eq:emp_ratio_AM}
%        \sum_i \sum_j \omega_{js}d_{ij}^{-\vartheta}=\frac{L_s}{L_A}\sum_i \sum_j \omega_{jA}d_{ij}^{-\vartheta} ,
%    \end{equation}
%}
%Finally, we recover the local wage rate in each sector, $w_{is}$, by $w_{is} = \omega_{is}^{1/\theta}$.

Finally, given the recovered wages $w_{is}$, we back out the floor space rents $Q_i$ from the floor space market clearing condition (\ref{eq:land_market}). We also back out the land rents $\mathbb{Q}_i$ from equation (\ref{eq:fp_lp}).

% \aiba{
% Maybe (\ref{eq:emp_ratio_AM}) is wrong and the correct equation is
% \begin{equation*}
%     \sum_i \sum_j \left( A_i\mathbb{P}_i^{-\alpha}Q_i^{-(1-\alpha)} \right)^\theta \omega_{js}d_{ij}^{-\vartheta}=\frac{L_s}{L_A}\sum_i \sum_j \left( A_i\mathbb{P}_i^{-\alpha}Q_i^{-(1-\alpha)} \right)^\theta \omega_{jA}d_{ij}^{-\vartheta} ,
% \end{equation*}
% which can be derived from (\ref{eq:prob_live_work}) by calculating $\lambda_s=\sum_i\sum_j\lambda_{ijs}$. But we can't use this equation for calibration because $A_i$, $\mathbb{P}_i$, and $Q_i$ are not yet calibrated at this stage.

% Instead, I propose the following procedure.
% \begin{itemize}
%     \item From (\ref{eq:prob_live}), the share of workers who choose to live in location $i$ and work for sector $s$ is
%     \begin{equation*}
%         \lambda_{is}^R=\sum_j\lambda_{ijs}=
%         \frac{\left( A_{i}\mathbb{P}_{i}^{-\alpha} Q_i^{-(1-\alpha)} \right)^\theta W_{is}}{\sum_k \sum_l \left( A_{k} \widetilde{w}_l \mathbb{P}_{k}^{-\alpha} Q_k^{-(1-\alpha)} \right)^\theta \tau_{kl}^{-\theta}} ,
%     \end{equation*}
%     where $W_{is} \equiv \sum_j\omega_{js}d_{ij}^{-\vartheta}$ measures the sector-$s$ commuting market access of location $i$.
%     \item Taking the ratio between sector $A$ and $s$ ($\neq A$), we have
%     \begin{equation*}
%         \frac{\lambda_{is}^R}{\lambda_{iA}^R}=\frac{W_{is}}{W_{iA}} ,
%     \end{equation*}
%     or equivalently,
%     \begin{equation*}
%         \frac{R_{is}}{R_{iA}}=\frac{W_{is}}{W_{iA}} .
%     \end{equation*}
%     \item We can observe $R_{is}$ from the data. In addition, $W_{iA}$ is readily calculated because we have already obtained $\omega_{iA}$. Thus, the commuting market accesses of location $i$ for sectors $M$, $N$, and $B$ are recovered as
%     \begin{equation*}
%         W_{is}=\frac{R_{is}}{R_{iA}} \times W_{iA} \quad \text{for} \quad s \in \{M,N,B\} .
%     \end{equation*}
%     \item From (\ref{eq:LMC}), $\omega_{is}$ for $s \in \{M,N,B\}$ is calibrated as
%     \begin{equation*}
%         \omega_{is}=\frac{L_{is}}{\sum_k \frac{d_{ki}^{-\vartheta}R_{ks}}{W_{ks}}} ,
%     \end{equation*}
%     or equivalently,
%     \begin{equation*}
%         \omega_{is}=\frac{L_{is}}{\sum_k \frac{d_{ki}^{-\vartheta}R_{kA}}{W_{kA}}} .
%     \end{equation*}
% \end{itemize}
% }

%\paragraph{Step 4: Price-adjusted residential amenity.}
%By (\ref{eq:prob_live}), the share of workers who reside in location $i$ can be written as
%\begin{equation}
%\label{eq:prob_live_calibration}
%    \lambda_i^R=\frac{\mathbb{A}_i W_i Q_i^{-(1-\alpha)\theta}}{\sum_k \mathbb{A}_k W_k Q_k^{-(1-\alpha)\theta}},
%\end{equation}
%where $W_i \equiv \sum_j \sum_s \omega_{is}d_{ij}^{-\vartheta}$ measures location $i$'s commuting market access and $\mathbb{A}_i \equiv (A_i\mathbb{P}_i^{-\alpha})^\theta$ is the price-adjusted residential amenity in location $i$. Because $W_i$ is immediately calculated using the values of $\vartheta$ and $\omega_{is}$, (\ref{eq:prob_live_calibration}) can be solved for $\mathbb{A}_i$, given the observed values of $\lambda_i^R$ and $Q_i$.
        
\paragraph{Step 4: Price index of non-tradable services.}
Given the values of $w_{is}$ for any pair of $i$ and $s$ and $\xi_{ij}$ for any pair of $i$ and $j$, we can back out the equilibrium values of $P_i^N$ that rationalize the observed patterns of shopping and commuting trips. Specifically, we solve the goods market clearing condition (\ref{eq:goods_market}) with respect to $P_i^N$.
In solving equations (\ref{eq:goods_market}), we need the value of US-base-sourced consumption $\Phi_i$, which is defined as $\Phi_i \equiv m_i\mathbb{Q}_i^{US}S_i+e_{US}L_i^{US}$. We allocate US military personnel across municipalities in proportion to US-base land area and calibrate $\mathbb{Q}_i^{US}$ and $e_{US}$ so that US-base land rents and US personnel's purchases of Okinawan goods and services separately match their observed shares of GNI. More details are provided in Appendix \ref{appsec:calibration_GNI}. Together with the base-sector wage level calibrated in the previous step, this procedure matches the three separately reported base-related GNI shares: labor income from base employment, base land rents, and US personnel's purchases of goods and services.

In solving the goods market clearing condition (\ref{eq:goods_market}) with respect to $P_i^N$, we also need to determine the value of $\kappa$, governing the substitutability between tradable and non-tradable goods. After solving the equation (\ref{eq:goods_market}) under many values of $\kappa$, we set $\kappa$ to match the relative spending share for restaurants in the largest urban place (Naha) and the entire Okinawan economy, which is 1.16 according to the Okinawa Household Expenditure Survey in 1970. We find $\kappa=1.87$, which is well within the 95\% confidence interval reported by \cite{hobijn2019sticker}.\footnote{
Note that in our model, housing consumption is excluded from the non-tradable sector in defining the substitutability between tradable and non-tradable sectors. \citeauthor{hobijn2019sticker}'s (\citeyear{hobijn2019sticker}) estimate is conceptually in line with ours in distinguishing housing consumption from other non-tradable goods. Studies that include housing as a non-tradable good (e.g., \citealt{duarte2010structural}; \citealt{fajgelbaum2022argentina}; \citealt{buera2022structural}) tend to find a smaller elasticity of substitution between tradable and non-tradable goods. 
} 

\paragraph{Step 5: Locational fundamentals.}
We back out the exogenous amenities and productivity fundamentals. First, 
equation (\ref{eq:priceA}) enables us to identify the price-adjusted productivity of agricultural goods in location $i$, defined as $\widetilde{b}_i^A \equiv P^A b_i^A$, because the values of $w_{iA}$ and $Q_i$ are already known.
Similarly, we define the price-adjusted productivity of manufacturing goods in location $i$ as $\widetilde{b}_i^M \equiv P^M b_i^M$, and combine equations (\ref{eq:priceM}) and (\ref{eq:productivityM}) to obtain the following expression: $\widetilde{b}_i^M=w_{iM}^{1-\gamma^M}Q_i^{\gamma^M}L_{iM}^{-\eta}$. Substituting the values of $w_{iM}$, $Q_i$, and $L_{iM}$, we can recover the value of $\widetilde{b}_i^M$.
For the non-tradable service sector, we can solve equation (\ref{eq:price_index_N_3}) with respect to $B_i^N$, given the values of $w_{iN}$, $Q_i$, and $\mathbb{P}_i^N$. After obtaining the value of $B_i^N$, we can readily back out $\widetilde{b}_i^N$ from equation (\ref{eq:productivityN}).
Finally, we back out the amenity associated with living in $i$ and working in sector $s$. Now that $\widetilde{A}_{is}$, $\mathbb{P}_i$ and $Q_i$ are known, $A_{is}$ can be obtained by using $A_{is}=\widetilde{A}_{is}\left(\mathbb{P}_{i}^{\alpha}Q_i^{1-\alpha}\right)^\theta$.
%Finally, we recover the values of $A_i$ that rationalize the observed distribution of residents. By equation (\ref{eq:prob_live}), the share of workers who reside in location $i$ can be written as
%\begin{equation}
%\label{eq:prob_live_calibration}
%    \lambda_i^R=\frac{A_i W_i \left( \mathbb{P}_i^{-\alpha} Q_i^{-(1-\alpha)} \right)^\theta}{\sum_k A_k W_k \left( \mathbb{P}_k^{-\alpha} Q_k^{-(1-\alpha)} \right)^\theta},
%\end{equation}
%where $W_i \equiv \sum_j \sum_s \omega_{is}d_{ij}^{-\vartheta}$ measures location $i$'s commuting market access. Because $W_i$ is immediately calculated using the values of $\vartheta$ and $\omega_{is}$, Equation (\ref{eq:prob_live_calibration}) can be solved for $A_i$, given the observed values of $\lambda_i^R$ and the calibrated values of $\mathbb{P}_i$ and $Q_i$. As the right-hand side of (\ref{eq:prob_live_calibration}) is homogeneous of degree zero with respect to $A_i$, we normalize the scale of $A_i$ such that their geometric mean equals one.

\paragraph{Model fit.} As we have used information on the population and employment of each sector in each municipality in our calibration, our model prediction perfectly matches the observed values. To assess predictions for variables not used in the calibration, we compare model-predicted bilateral commuting probabilities with observed probabilities in 1967, and model-implied floor-space rents with observed municipal house prices.
Figure \ref{fig:modelfit} reports correlations of 0.958 for commuting probabilities and 0.873 between model-implied floor-space rents and observed house prices. Overall, the model successfully predicts these variables not used for our calibration.

%%%%%%%%%%%%%%%%%%%%%%%%%%%%%%%
%%% Counterfactuals
%%%%%%%%%%%%%%%%%%%%%%%%%%%%%%%

\section{Counterfactual Analysis}\label{sec:counterfactuals}

We quantify the effects of US-base-sourced expenditure and employment by comparing the calibrated 1970 equilibrium with a series of counterfactual equilibria. Section \ref{sec:benchmark_cf} jointly removes US-base-sourced consumption, base-rent income, and employment while keeping base land unavailable for private use, holding other exogenous fundamentals at their calibrated levels, and fixing Okinawa's aggregate worker population. This benchmark case allows us to isolate the effect of consumption demand and employment at the US bases. Section \ref{sec:counterfactual_migration} then examines the role of sectoral and spatial mobility. Section \ref{sec:counterfactual_land} adds the return of base land to private use, while Section \ref{sec:counterfactual_aggregate} summarizes the Okinawa-wide outcomes across counterfactuals. 
Section \ref{sec:elasticity} then summarizes the direct and spillover effects in elasticity form under the benchmark specification.
Finally, Section \ref{sec:counterfactual_external_validity} assesses external validity of our results by examining how the benchmark counterfactual results depend on external mobility and manufacturing import supply. 
Appendix \ref{appsec:counterfactual} describes further details.

\subsection{Counterfactual effects of US-base-related flows and land use}\label{sec:counterfactual_main}

\subsubsection{Benchmark: Removing the US-base sourced consumption and employment}\label{sec:benchmark_cf}


We first ask how jointly removing US-base sourced consumption and employment changes the spatial distribution of economic activities. 
More specifically, we set to zero the sum of the consumption of the US military personnel and land-rent income (i.e., $\Phi_i=0$) and employment of Okinawans at the bases (i.e., $L_{iB}=0$ for every location $i$), while keeping base area unavailable for private use.
Therefore, the counterfactual considered in this section is different from the ``US-base closure,'' in which the base acreage is returned to Okinawans for private use in addition to eliminating the US-base sourced consumption and employment, which we analyze in Section \ref{sec:counterfactual_land}.


\paragraph{Geographical and sectoral reallocation.}

We begin with the population maps in Figure \ref{fig:cf_pop_change_map}, which show both the total number of residents and the numbers of residents working in agriculture, manufacturing, and non-tradable services.

\begin{figure}[p]
    \centering
    \includegraphics[width=\linewidth]{map/cf_all_total_AMN_Rcf_over_R_maps_en.png}
    \vspace{0.15in}
    \caption{Municipal worker-resident responses, overall and by civilian sector, in the benchmark counterfactual}
    \vspace{-0.2in}
    \floatfoot{\textbf{Note}: The maps report $100\times(\text{counterfactual}/\text{observed}-1)$ for total residents and for residents working in agriculture (A), manufacturing (M), and non-tradable services (N). Red denotes an increase and blue a decrease. Base employment is included in total residents but is not included in the three sector panels. The benchmark sets $\Phi_i=L_{iB}=0$ while leaving base acreage unavailable for private use. Municipal boundaries are those of 1970; labels report changes for selected municipalities. For the labeled municipalities, the shares of 1970 municipal area occupied by US bases are Nago 11.2\%, Kin 66.1\%, Ishikawa 17.4\%, Koza 65.3\%, Ginowan 38.9\%, Naha 22.7\%, and Itoman 1.2\%.}
    \label{fig:cf_pop_change_map}
\end{figure}

The top-left map shows considerable geographical mobility away from the US bases in response to the elimination of US-base-sourced consumption income and employment. More specifically, the number of worker residents falls by 26.57\% in Chatan, 18.51\% in Ishikawa, 20.45\% in Kin, 16.88\% in Koza, and 14.11\% in Ginowan, all of which have relatively large US-base area shares. In contrast, Naha, Nago, and Itoman, which have relatively low US-base area shares, gain worker residents by 5.62\%, 11.01\%, and 10.16\%. 

In addition to such large-scale geographical mobility, the sector maps in the rest of Figure \ref{fig:cf_pop_change_map} reveal considerable sectoral mobility. The number working in civilian non-tradable services falls in every municipality, with an aggregate-level decline of 9.63\%.
By contrast, the number of residents working in agriculture and manufacturing rises in every municipality, with the aggregate-level increases of 33.73\% and 47.87\%. Moreover, Figure \ref{fig:cf_pop_change_map} also reveals that these sectoral changes are stronger in areas with a higher share of the US bases, highlighting that these sectoral changes are driven by the evaporation of the US-base sourced consumption and employment.

Intuitively, the difference in the tradability of the outputs is the key factor in explaining these patterns. The loss of US-base sourced consumption and employment induces the decrease in the local non-tradable service demand, which reduces the non-tradable goods prices in Okinawa and reduces its employment share in Okinawa.  
In contrast, although the loss of US-base sourced consumption and employment also reduces the demand for agricultural and manufacturing goods, these goods are tradable so that it does not reduce the output prices. Therefore, relative to the non-tradable prices, the price of these tradable goods becomes higher, especiallly in areas around the US bases, and this makes the employment share in these sectors expand. 

To quantify how this geographical and sectoral mobility varies with initial base exposure, Figure \ref{fig:cf_result_scatter} summarizes the municipal responses in a scatterplot. Panel \ref{fig:cf_pop} plots the counterfactual-to-observed population ratio, while Panel \ref{fig:cf_empshare} reports counterfactual-minus-observed changes in broad sector-of-employment shares among residents: primary (agriculture), secondary (manufacturing), and tertiary (civilian non-tradables plus base employment).
The population ratio has a fitted slope of $-0.477$ (standard error $0.059$). Thus, a municipality with a base-area share 10 percentage points higher is predicted to lose population by 4.77\%. The broad sector composition also reverses most strongly in base-intensive municipalities: A municipality with 10 percentage-point higher base-area share is predicted to experiene 2.16 percentage points increase in the tertiary employment share, while their manufacturing- and agricultural-share are predicted to decrease by 1.44 and 0.72 percentage points, respectively. 

Note that these numbers in Figure \ref{fig:cf_result_scatter} are comparable to the stylized facts in Figure \ref{fig:reduced_crosssection1970} (Section \ref{sec:reducedform}) because we report the association between the US bases and the change in the relevant outcomes in the same form.\footnote{To ensure comparability, we include the US-base sector as tertiary in Figure \ref{fig:cf_result_scatter},  while Figure \ref{fig:cf_pop_change_map} does not include the US-base sector as tertiary for transparency.}
Generally, although they qualitatively match, the association is stronger in the stylized facts than in Figure \ref{fig:cf_result_scatter}. A possible reason for this discrepancy is that unlike our counterfactual scenario, the stylized facts include the effect of not only the consumption and employment of the US bases but also other effects such as infrastructure investment around the bases. We return to this issue in Section \ref{sec:counterfactual_fundamentals}.

\begin{figure}[t]
    \centering
    \subfloat[Residents\label{fig:cf_pop}]{%
      \includegraphics[width=0.475\textwidth]{figure/cf_baseline_residents.pdf}
    }
    \hfill
    \subfloat[Broad sector-of-employment shares\label{fig:cf_empshare}]{%
      \includegraphics[width=0.475\textwidth]{figure/cf_baseline_sector_share.pdf}
    }
    \vspace{0.2in}
    \caption{Exposure gradients in population and sector-of-employment shares}
    \vspace{-0.2in}
    \floatfoot{\textbf{Note}: Each point is a municipality, and the horizontal axis is the share of municipal area occupied by US bases in 1970. Panel \ref{fig:cf_pop} plots the counterfactual-to-observed ratio of residents. Panel \ref{fig:cf_empshare} plots counterfactual-minus-observed changes in sector-of-employment shares among residents. Red circles denote the tertiary share, which combines civilian non-tradable services and base employment; blue squares and green triangles denote manufacturing and agriculture. The three share changes therefore sum to zero within each municipality. Lines are linear fits, and the boxes report slopes and standard errors. The benchmark sets $\Phi_i=L_{iB}=0$ but does not return base land to private use.}
    \label{fig:cf_result_scatter}
\end{figure}



\paragraph{The effects on wages, rents, and non-tradable prices and varieties.}

Non-tradable wages fall in every municipality, and service varieties decline in 32 of 35 municipalities. These outcomes reflect the loss of base-related spending and the departure of residents from more exposed municipalities, which reduce demand for local services and therefore local labor demand in the non-tradable sector. 
%In the remaining three municipalities, population inflows sustain enough service demand for varieties to increase.
Floor-space rents tend to fall around the bases, which is natural given the population outflows.
These rent changes affect wages because firms use floor space in production. With tradable output prices fixed, lower floor-space costs allow firms to pay higher wages in the tradable sector. In addition to this, manufacturing wages rise because of agglomeration economies: the expansion of manufacturing employment around the bases raises the manufacturing productivity.
Regarding the service price index (equation \ref{eq:price_index_N_2}), lower non-tradable wages and, in the most exposed municipalities, lower rents decrease marginal costs (equation \ref{eq:markup}). These cost reductions outweigh the loss of non-tradable varieties, so the non-tradable price index falls in every municipality despite the contraction in service provision. Appendix \ref{appsec:counterfactual_prices} reports the corresponding maps and exposure gradients in Figures \ref{fig:cf_economic_map} and \ref{fig:cf_result_scatter_prices}.

\paragraph{Separating the consumption-side and employment channels.}\label{sec:decompose_bases}

The benchmark removes US-base-sourced consumption demand and employment jointly. To distinguish these two operating channels, we compare a consumption-only counterfactual, which sets $\Phi_i=0$ and therefore removes both US personnel's purchases and consumption coming from the US-base rents while retaining local base employment, with an employment-only counterfactual, which sets $L_{iB}=0$ while retaining $\Phi_i$. Both experiments leave base land occupied to focus on the role of US-base sourced consumption and employment. Appendix \ref{appsec:counterfactual_channels} reports the corresponding maps and fitted exposure gradients, including the side-by-side comparison in Figure \ref{fig:cf_channel_scatter}.

The two shocks have qualitative similarities: both shift residents away from more base-exposed municipalities, increase the numbers of residents working in agriculture and manufacturing in every municipality, and lower non-tradable wages and prices throughout Okinawa. Their magnitudes and service-sector quantity responses, however, differ sharply. The fitted resident gradient is $-0.190$ when $\Phi_i$ is set to zero but $-0.275$ when base employment is removed, indicating that base jobs play the larger role in attracting population around the US bases. The consumption-side shock reduces non-tradable employment in every municipality and by 14.20\% islandwide, and reduces service varieties everywhere. Removing base employment instead raises the non-tradable employment in every municipality and by 5.64\% islandwide, while varieties rise in 34 of 35 municipalities as former base workers move into the private non-tradable sector.
Overall, the comparison highlights that US personnel's purchases and base-rent income play a larger role in sustaining non-tradable demand and market size, whereas base employment plays a larger role in attracting more population around the US bases. We return to this distinction  in the elasticity analysis in Section~\ref{sec:elasticity}.

\subsubsection{Migration as a spatial adjustment margin}\label{sec:counterfactual_migration}

The benchmark allows residents to adjust their municipality of residence as well as their sector and workplace, implying that both the ``sectoral reallocation'' and ``geographical reallocation'' are feasible in response to the removal of the US-base sourced consumption and employment.

To separate the role of these two adjustment margins, we re-solve the benchmark while imposing $R_i^{cf}=R_i$ in every municipality.
This restriction fixes municipal resident totals but leaves sector and workplace choices, commuting and shopping patterns, and local wages, rents, and prices endogenous.\footnote{Therefore, precisely speaking, this exercise still accommodates changes in the commuting pattern as a geographical response.}

Figures \ref{fig:cf_fixR_pop_map} and \ref{fig:cf_fixR_economic_map} in Appendix \ref{appsec:counterfactual_fixed_residence} report the corresponding municipal maps, while Figure \ref{fig:cf_fixed_residence_scatter} summarizes the fitted gradients.
Under fixed residence, the fitted gradients in the tertiary, manufacturing, and agricultural sector-of-employment shares are $-0.191$, $+0.141$, and $+0.050$, compared with $-0.216$, $+0.144$, and $+0.072$ in the benchmark case. Thus, the structural reallocation does not arise because of workers' reallocation toward tradable-intensive municipalities, but sectoral reallocation \textit{within} the same municipality plays the key role.

Residential reallocation matters in determining where and how the local adjustment occurs. 
Even though residence is fixed, the employment gradients are $-0.238$ for non-tradables, $+0.882$ for manufacturing, and $+0.732$ for agriculture, compared with $-0.606$, $+0.345$, and $+0.371$ in the benchmark. This implies that commuting responses could generate a substantial reallocation of employment even if workers' residences are fixed. 
At the same time, the floor-space-rent gradient flattens from $-0.167$ to $-0.030$, and the variety gradient falls in magnitude from $-0.610$ to $-0.273$. These attenuated gradients are consistent with residential sorting amplifying local housing and non-tradable demand responses.


\subsubsection{Returning base land to private use}\label{sec:counterfactual_land}

The benchmark and fixed-residence exercises keep base acreage unavailable for private use. We now ask how the results change when that land is released for housing and private production. By increasing the amount of available land for production or residence, this additional land-supply channel may offset some of the economic losses caused by removing US-base-sourced expenditure and employment. The value of returned base land is also central to policy discussions in assessing the economic impact of base returns: some studies argue that the land-occupying effect is large (e.g., \citealt{okinawa07}), while others find a smaller effect (e.g., \citealt{kusuyama12}).

Specifically, we repeat the benchmark analysis in Section \ref{sec:benchmark_cf}, but the only difference is that we additionally set $m_i=0$ for every municipality: all base acreage becomes available for private use.
Note that we assume that the returned land has the same municipality-specific technology as other private land in the same municipality. We therefore abstract from any additional productivity associated with releasing a large contiguous site (\citealt{hornbeck2017creative}). Consistent with this abstraction, \citet{yamasaki2021samurai} documents no lot-size premium in Japan in 1970, when high-rise construction was not yet widely feasible, although large sites could become especially valuable if skyscraper technology were available.


Figure \ref{fig:cf_result_scatter_return} in Appendix \ref{appsec:counterfactual_landreturn} shows that returning the land makes the residential reallocation much more muted. In Panel \ref{fig:cf_pop_return}, the resident-ratio slope is $-0.034$ (standard error $0.061$), compared with $-0.477$ in the benchmark, implying that the US base share and the change in population are no longer correlated. 
Intuitively, more availability of land around the US bases directly increases utility through inducing lower housing costs and indirectly through increasing labor productivity due to more land input in production, offsetting the negative effect due to the loss of the local non-tradable demand from the base and the US-base employment.
By contrast, Panel \ref{fig:cf_empshare_return} shows a broad tertiary sector-of-employment-share slope of $-0.220$, close to the benchmark value of $-0.216$, while agriculture and manufacturing gain share. Overall, at least in our context, land availability therefore strongly dampens residential reallocation without undoing the sectoral reallocation.


\subsubsection{Aggregate outcomes across counterfactuals}\label{sec:counterfactual_aggregate}

The preceding subsections focused on the geography and sectoral composition of adjustment. We now turn to the aggregate economic consequences for Okinawa as a whole. Table \ref{tab:macro-counterfactual-summary} summarizes the effects on welfare, GNI, and economy-wide employment shares. The first three rows report the benchmark (Section \ref{sec:counterfactual_main}), the fixed-residence variant (Section \ref{sec:counterfactual_migration}), and the counterfactual that also returns base land to private use (Section \ref{sec:counterfactual_land}). As in the last part of Section \ref{sec:counterfactual_main}, the final two rows isolate the consumption-side and employment channels: one sets $\Phi_i=0$ while retaining local base employment, and the other sets $L_{iB}=0$ while retaining $\Phi_i$. 

{\floatsetup[table]{capposition=top}
% Table maintained directly in draft.tex.
\begin{table}[!t]
  \centering
  \caption{Okinawa-wide effects of counterfactual scenarios}
  \label{tab:macro-counterfactual-summary}
  \begin{threeparttable}
  \small
  \sisetup{
    input-signs = +-,
    retain-explicit-plus = true,
    table-number-alignment = center
  }
  \begin{tabularx}{\textwidth}{@{}>{\raggedright\arraybackslash}X *{6}{S[table-format=+2.2]}@{}}
    \toprule
    & \multicolumn{2}{c}{Aggregate effects (\%)} & \multicolumn{4}{c}{Employment-share changes (pp)} \\
    \cmidrule(lr){2-3} \cmidrule(l){4-7}
    Scenario & \multicolumn{1}{c}{Workers' Welfare} & \multicolumn{1}{c}{GNI} & \multicolumn{1}{c}{Agric.} & \multicolumn{1}{c}{Manuf.} & \multicolumn{1}{c}{Non-trad.} & \multicolumn{1}{c}{U.S. base} \\
    \midrule
    \multicolumn{7}{@{}l}{\textit{Benchmark and model variants}} \\
    Benchmark counterfactual & -7.25 & -11.99 & +5.58 & +10.26 & -4.92 & -10.93 \\
    Fixed residence; sector mobility allowed & -7.10 & -14.04 & +4.85 & +10.39 & -4.31 & -10.93 \\
    Base land also returned & -1.95 & -11.23 & +5.35 & +10.19 & -4.61 & -10.93 \\
    \addlinespace[0.35em]
    \multicolumn{7}{@{}l}{\textit{Partial-channel counterfactuals}} \\
    US-base sourced consumption removed & -2.42 & -11.60 & +2.57 & +4.68 & -7.25 & 0.00 \\
    Base employment removed & -5.22 & -1.87 & +2.92 & +5.12 & +2.88 & -10.93 \\
    \bottomrule
  \end{tabularx}
  \begin{tablenotes}[flushleft]
    \footnotesize
    \item \textit{Notes:} Welfare and GNI entries are percentage changes from the initial equilibrium, $100(\text{counterfactual}/\text{initial}-1)$. Employment entries are changes in Okinawa-wide sectoral employment shares in percentage points (pp). Initial employment shares are 16.55, 21.44, 51.08, and 10.93 percent for agriculture, manufacturing, non-tradable services, and US-base employment, respectively. 
  \end{tablenotes}
  \end{threeparttable}
\end{table}
}

In the first row of Table \ref{tab:macro-counterfactual-summary}, the benchmark counterfactual predicts both a substantial sectoral reallocation and sizable aggregate losses: Welfare falls by 7.25\% and GNI by 11.99\%. 
While this is large in absolute value, it is relatively small given that the US-base sourced consumption and employment constituted around 30\% of the GNI in the observed Okinawan economy in 1970.
The employment shares of agriculture and manufacturing rise by 5.58 and 10.26 percentage points, while the shares of civilian non-tradable services and base employment fall by 4.92 and 10.93 percentage points. Thus, the islandwide counterpart of the spatial patterns documented above is a shift from services and base employment toward tradable production, accompanied by lower income and welfare.

In this sense, this result is broadly consistent with the ``Dutch disease'' hypothesis that a large facility, which accompanies large consumption demand and employment, does not contribute to economic activities by preventing efficient allocation of resources  (\citealt{gollin2016urbanization}; \citealt{allcott2018dutch}). Our result on Okinawa suggests that even if the military base's demand and employment disappear, sectoral and geographical mobility can substantially mitigate the negative economic impact and restore a large part of the lost income and welfare. That said, the overall economic impact of the US bases in 1970 is still estimated to be positive.

In the second row of Table \ref{tab:macro-counterfactual-summary}, restricting residential mobility leaves this sectoral reallocation broadly unchanged compared with the benchmark. Under fixed municipal residence, the agricultural, manufacturing, and non-tradable employment shares change by $+4.85$, $+10.39$, and $-4.31$ percentage points. GNI also falls by 14.04\%, compared with 11.99\% in the benchmark. This comparison is consistent with internal migration somewhat cushioning the aggregate income loss by reallocating residents toward locations with higher productivity.\footnote{The reported welfare in the fixed residence case hardly differs from the benchmark value. However, note that the reported entry evaluates the same free-residence inclusive-value index used in the benchmark at the fixed-residence equilibrium, which is computable but does not have the standard expected welfare interpretation because the location choice has been fixed.}
However, the aggregate impacts with or without residential reallocation are similar, implying that sectoral adjustments play the key role in mitigating the negative economic impact of the loss of the US-base consumption and employment. 

In the third row of Table \ref{tab:macro-counterfactual-summary}, returning base land affects a different aggregate margin: It reduces the welfare loss from 7.25\% to 1.95\%, a recovery of 5.30 percentage points, but reduces the GNI loss by only 0.76 percentage points, from 11.99\% to 11.23\%. Its civilian-sector employment-share changes---$+5.35$, $+10.19$, and $-4.61$ percentage points in agriculture, manufacturing, and non-tradable services---also remain close to the benchmark. Additional land therefore offsets much of the welfare loss without reversing either the income decline or the service-to-tradable reallocation.

In the fourth row, removing the US-base sourced consumption lowers welfare by 2.42\% and GNI by 11.60\%; the non-tradable employment share falls by 7.25 percentage points, which is a larger response than in the benchmark case. 
In the fifth row, removing the US-base employment lowers welfare by 5.22\% but GNI by only 1.87\%; displaced base workers enter all three civilian sectors, and the non-tradable share rises by 2.88 percentage points. The consumption-side counterfactual therefore produces the larger GNI decline and service-sector contraction, whereas the employment-only counterfactual produces the larger welfare decline. 

Overall, the aggregate effects are economically large and broadly align with the local effects discussed in Sections \ref{sec:counterfactual_main}--\ref{sec:counterfactual_land}, but
the aggregate economic loss in neither the benchmark welfare decline nor its GNI decline approaches contemporary claims that an immediate US withdrawal would reduce Okinawan living standards by more than half \citep{toriyama09}.
%The aggregate manufacturing crowding-out effect depends on the trade environment, as we show in Section \ref{sec:counterfactual_external_validity}. 

\subsection{Direct and Spillover Elasticities of the US-base Consumption and Employment}
\label{sec:elasticity}

Under the benchmark counterfactual scenario, we examine separate 1\% increases in FSI and base employment in one municipality, holding other exogenous inputs fixed. This allows us to express the responses as elasticities, providing normalized measures of the effects of US-base-sourced consumption and employment that facilitate comparisons of effect sizes with existing studies. 
The increase in only one municipality also allows us to how consumption and employment shock has a local effect on its own and a spillover effect to the rest of the economy.

Table \ref{tab:local_base_elasticities} reports the average elasticity across all municipalities.
The first column reports the local elasticity in the municipality receiving the shock; the second reports the elasticity of the Okinawa-wide aggregate with respect to FSI or base employment in that municipality; and the third reports the percentage of the Okinawa-wide impact accounted for by the local effect in the source municipality. Subtracting this percentage from 100\% gives the share attributable to effects in all other municipalities, that is, the net spillover contribution. Appendix \ref{appsec:local_elasticities} provides the calculation details.

{\floatsetup[table]{capposition=top}
\begin{table}[tbp]
\centering
\caption{Local and aggregate responses to a 1\% increase in one municipality}
\label{tab:local_base_elasticities}
\begin{threeparttable}
\small
\begin{tabular*}{\textwidth}{@{\extracolsep{\fill}}lrrr@{}}
\toprule\toprule
\multicolumn{4}{l}{\textit{Panel A. FSI: a 1\% increase in one municipality}} \\
Outcome & \makecell[r]{Local response\\(\%)} & \makecell[r]{Okinawa-wide\\response (\%)} & \makecell[r]{Local share of\\aggregate impact (\%)} \\
\midrule
Agricultural employment & $-0.0187$ & $-0.00482$ & $11.7$ \\
Manufacturing employment & $-0.0176$ & $-0.00656$ & $13.2$ \\
Non-tradable employment & $+0.1074$ & $+0.00432$ & $47.3$ \\
Non-base employment & $+0.0394$ & $0.00000$ & $\textnormal{NA}$ \\
Resident workers & $+0.0126$ & $0.00000$ & $\textnormal{NA}$ \\
\addlinespace
Agricultural wage & $-0.0041$ & $-0.00017$ & $80.8$ \\
Manufacturing wage & $-0.0035$ & $-0.00070$ & $21.2$ \\
Non-tradable wage & $+0.0537$ & $+0.00405$ & $24.6$ \\
Floor-space rent & $+0.0096$ & $+0.00107$ & $33.3$ \\
Local non-tradable price index & $+0.0126$ & $+0.00236$ & $18.7$ \\
\midrule
\addlinespace
\multicolumn{4}{l}{\textit{Panel B. Base employment: a 1\% increase in one municipality}} \\
Outcome & \makecell[r]{Local response\\(\%)} & \makecell[r]{Okinawa-wide\\response (\%)} & \makecell[r]{Local share of\\aggregate impact (\%)} \\
\midrule
Agricultural employment & $-0.0233$ & $-0.00713$ & $6.3$ \\
Manufacturing employment & $-0.0216$ & $-0.00939$ & $14.9$ \\
Non-tradable employment & $+0.0023$ & $-0.00112$ & $-2.9$ \\
Non-base employment & $-0.0091$ & $-0.00423$ & $9.5$ \\
Resident workers & $+0.0409$ & $0.00000$ & $\textnormal{NA}$ \\
\addlinespace
Agricultural wage & $-0.0053$ & $+0.00009$ & $-72.1$ \\
Manufacturing wage & $-0.0044$ & $-0.00065$ & $27.3$ \\
Non-tradable wage & $+0.0068$ & $+0.00308$ & $18.6$ \\
Floor-space rent & $+0.0123$ & $+0.00030$ & $91.7$ \\
Local non-tradable price index & $+0.0075$ & $+0.00272$ & $18.0$ \\
\bottomrule
\end{tabular*}
\begin{tablenotes}[flushleft]
\footnotesize
\item \textbf{Note}: Each experiment raises FSI or base employment by 1\% in one municipality, holding the other channel and occupied base land fixed. FSI comprises US personnel expenditure and base-rent income, excluding Okinawan base workers' payroll. The first two columns average local and aggregate responses across the same 29 sources; all 35 municipalities enter the aggregate. The third column gives the local share of the net aggregate change; negative shares indicate offsetting local changes and NA denotes a zero aggregate response due to the assumption of fixed total population of Okinawa. Non-base employment comprises agriculture, manufacturing, and non-tradables.  Appendix \ref{appsec:local_elasticities} defines the weights and calculations.
\end{tablenotes}
\end{threeparttable}
\end{table}}

\paragraph{Elasticities with respect to US-base-sourced consumption.}
A 1\% increase in FSI raises local non-tradable employment by 0.1074\% and non-base employment by 0.0394\%. Each additional local service job is accompanied by losses of 0.112 agricultural and 0.088 manufacturing jobs in the same municipality. Local tradable-sector contraction thus offsets only one fifth of the service employment gain. Intuitively, additional spending supports service production, and the labor market competition reallocates some workers from tradable sectors to the non-tradable sector, with the overall positive employment impact.

In terms of the magnitude of the spillovers, service expansion is more concentrated in the source municipality than manufacturing contraction. The source municipality accounts for 47.3\% of the aggregate service employment gain but only 13.2\% of the manufacturing loss. 
Shopping costs concentrate additional service demand near the income source, directly supporting nearby service providers. The resulting service expansion, however, draws workers from tradable sectors across municipalities through commuting and relocation, spreading manufacturing employment losses beyond the source municipality.
Our results show that manufacturing accounts for a large part of the job losses, with a smaller contribution from agriculture (Appendix \ref{appsec:local_elasticities}). 

Both commuting and relocation yield spillover effects onto other municipalities. 
The importance of commuting can be seen in the contrast between local employment and population responses: local workplace employment rises by 0.0310\% and population increases by 0.0126\%. While population increase implies the role of residential relocation, workers can also exploit the additional FSI by commuting to the source municipality from elsewhere, yielding a larger response of employment.

Local service wages, prices, and floor-space rents rise by 0.0537\%, 0.0126\%, and 0.0096\%, respectively, with most of their aggregate adjustment occurring elsewhere. Service demand supports higher wages and prices, whereas higher rents lower agricultural and manufacturing wages at fixed tradable output prices. Manufacturing contraction also reduces manufacturing wages by weakening agglomeration economies. 

The estimated impact of the US-base sourced consumption on employment and wages is broadly consistent with the existing literature in different settings.
After normalization by the induced change in nominal consumer expenditure, our result implies a local non-tradable employment response of 0.412\% per 1\% increase in expenditure, close to the 0.48\% inferred by \citet{mian2014explains}.
The impact of the US-base sourced income on the employment and wage is also broadly consistent with the cash-transfer evidence of \citet{gerard2026cash} and \citet{mendes2026macroeconomic}, although employment and wage adjustment differ across settings.\footnote{\citet{faber2019tourism} also find positive local employment effects of external demand. Appendix \ref{appsec:local_elasticities} discusses differences in the shocks and outcomes across these studies.} \citet{gerard2026cash} report increased labor force participation and little wage adjustment holding worker composition fixed. Their estimates imply a local private formal employment response of about 0.15\% per 1\% increase in transfers, compared with our local non-base employment response of 0.0394\% per 1\% increase in FSI with local service wages rising by 0.0537\%. In our benchmark, Okinawa's workforce is fixed, so service expansion draws on existing workers through sectoral and spatial reallocation, which induces a wage increase rather than the employment adjustment. In contrast, \citet{mendes2026macroeconomic} find higher formal-sector real wages without a statistically significant increase in total employment, which is likely if labor supply elasticity is smaller than in our setting. 
Taken together, our results are consistent with the available relevant estimates in the literature, although the impact seems to depend on labor supply elasticity that differs by context. 


\paragraph{Elasticities with respect to base employment.}
A 1\% increase in Okinawan employment at the bases, holding FSI and occupied base land fixed, raises total local workplace employment by 0.1201\% and reduces non-base employment by 0.0091\%. Converting elasticities into numbers of jobs, we find that, on average, each additional base job displaces 0.044 non-base jobs locally, leaving a net increase of 0.956 local jobs. 
This suggests a limited crowding-out of local private jobs.
However, limited local displacement coexists with substantial displacement elsewhere: Only 9.5\% of the aggregate decline occurs in the source municipality. Commuting and relocation allow local hiring to draw on workers from other municipalities, leading to the spillover of the private-sector employment contraction effect. Indeed, resident workers increase by 0.0409\%, or 0.473 workers per base job, with the remaining 0.483 local jobs accounted for by increased net inward commuting. 

Note that the elasticity of the local service employment is substantially smaller than in the case of the FSI. 
The US-base employment is similar to the FSI in that it also creates the spending expansion through the payroll of Okinawan base workers. Base hiring, however, also competes directly with civilian employers for Okinawa's fixed workforce, while the inflow of residents increases housing demand and floor-space rents. These demand and cost pressures leave local service employment only 0.0023\% higher on average. It also induces increases in nominal service wages, service prices, and floor-space rents of 0.0068\%, 0.0075\%, and 0.0123\%, respectively. The small local service employment gain is more than offset by service job losses elsewhere, yielding an Okinawa-wide decline of 0.00112\%. 

Our total local employment multiplier of 0.956 is broadly consistent with the mostly below-one base-closure estimates of \citet{hooker2001measuring} and the near-unit employment multiplier from prison openings in \citet{chirakijja2024local}. Our small non-base response of $-0.044$ is also comparable in magnitude to the statistically insignificant estimate of 0.085 private jobs per additional public-sector job in \citet{faggio2014effect}. \citet{zou2018local} further highlights population adjustment accompanying civilian job losses after military personnel reductions. While our results are in line with these results, the comparisons remain suggestive because estimands are different: \citet{hooker2001measuring} and \citet{zou2018local} capture concurrent changes in personnel and spending, \citet{chirakijja2024local} examines prison openings, and \citet{faggio2014effect} examine broader changes in public-sector employment. Our results shed new light on this discussion by distinguishing the effects of consumption of the military personnel and base employment.

\subsection{External validity: international mobility and trade}\label{sec:counterfactual_external_validity}

Our baseline counterfactual analysis so far has assumed fixed population of Okinawa and no trade cost for tradable goods, which is justified given Okinawa's island geography with restrictions on movement to mainland Japan and active external trade. We relax these assumptions separately to assess how the benchmark counterfactual results on welfare, aggregate income, and sectoral reallocation depend on external mobility and manufacturing import supply.\footnote{Note that altering these assumptions affect the counterfactuals, not the calibration of the observed equilibrium.} Appendix \ref{appsec:counterfactual_smallopen} describes the population adjustment, and Appendix \ref{appsec:counterfactual_external_maps} reports the extended specifications and full results in Table \ref{tab:macro-counterfactual-summary-additional}.

\paragraph{External migration.}
We accommodate migration in and out of Okinawa, which was prohibited in our baseline analysis.
As described in Appendix \ref{appsec:counterfactual_smallopen}, the migration elasticity parameter $\theta_{\mathrm{out}}$ governs how strongly workers' choice between Okinawa and the outside economy responds to relative expected welfare. We use $\theta_{\mathrm{out}}=1.1$ as an intermediate sensitivity case, informed by \citet{ortega2013effect}, whose preferred income elasticity of international migration flows is $0.76$ and whose estimate for migration within Europe is $1.90$. 

Table \ref{tab:macro-counterfactual-summary-additional} shows that our main findings on workers' welfare and sectoral reallocation are robust to relaxing the fixed-population assumption. Across external-migration elasticities of 0.5--2.19, welfare losses remain between 7.26\% and 7.28\%, almost identical to the benchmark's 7.25\%. Sectoral employment-share changes also remain close to the benchmark: manufacturing gains 9.96--10.19 percentage points, compared with 10.26 points, and civilian non-tradable services lose 5.01--5.32 points, compared with 4.92 points. Thus, the welfare loss and the shift from services and base employment toward tradable production do not depend on ruling out migration between Okinawa and the outside economy.
With the external migration elasticity set to $\theta_{\mathrm{out}}=1.1$, the aggregate population falls by 7.91\% and GNI by 18.93\%, compared with the benchmark GNI decline of 11.99\%. These larger aggregate income losses coexist with stable welfare and sectoral-composition results. The main spatial pattern of population is also preserved: a municipality with 10 percentage-points increase in the base-area share is predicted to lose population by 4.45\%, very close to 4.77\% in our benchmark prediction.

\paragraph{Manufacturing imports and the trade environment.}
We accommodate the possibility that manufacturing imports face frictions.
To capture such frictions, we denote by $\varepsilon_M$ the price elasticity of manufacturing import supply, which measures the percentage response of the quantity supplied to Okinawa by external producers to a 1\% increase in the manufacturing-goods price. The benchmark assumes perfectly elastic import supply ($\varepsilon_M\rightarrow\infty$): imports can freely adjust at a fixed manufacturing price, so the loss of base-related demand lowers service prices relative to tradable prices and shifts workers toward tradable production. With finite $\varepsilon_M$, the decline in import demand due to eliminating US-base sourced consumption and employment also lowers the manufacturing price. The resulting fall in manufacturing returns relative to agriculture weakens the expansion of manufacturing and redirects more workers toward agriculture. At the same time, cheaper manufactured goods reduce consumption costs and the cost of inputs used to construct floor space, cushioning the decline in workers' real welfare.

Table \ref{tab:macro-counterfactual-summary-additional} shows that these price adjustments affect welfare, aggregate income, and sectoral composition differently under different values of $\varepsilon_M$. A useful benchmark is $\varepsilon_M=10$, guided by \citet{nicita2018cooperation}, who estimate that the median elasticity of export supply faced by World Trade Organization member importers is approximately 10.
In this case, welfare falls by 6.25\% and GNI by 18.64\%, compared with 7.25\% and 11.99\% in the benchmark. Manufacturing's employment share rises by 7.13 percentage points rather than 10.26, while agriculture's gain increases from 5.58 to 9.41 points. 
In this case, eliminating the US-base sourced consumption and employment decreases the manufacturing prices, which dampens reallocation to manufacturing from the US-base sector. While this substantially lowers the nominal income, the welfare impact is cushioned because a lower manufacturing price also benefits workers as consumers. That said, qualitative implications remain similar to our benchmark case.\footnote{However, the aggregate effect can even reverse in an extreme scenario: when import quantities are fixed ($\varepsilon_M=0$), removing US-base-sourced consumption and employment reduces manufacturing's islandwide employment share by 1.77 percentage points (Table \ref{tab:macro-counterfactual-summary-additional}), implying that their presence supports aggregate manufacturing employment in that case.} However, even under elastic import supply, removing base-related flows induces a stronger shift toward manufacturing from services or base employment in more base-exposed municipalities. In this sense, the elastic import elasticity does not change the spatial tendency for bases to favor services over manufacturing.

\section{Role of Past Investment in Infrastructure and Durable Capital around the Bases}\label{sec:discussion_extensions}\label{sec:counterfactual_fundamentals}

The counterfactuals in Section \ref{sec:counterfactuals} quantify the effects of the ongoing consumption demand and direct employment emphasized in the policy discussions introduced at the outset of the paper. They remove base-related expenditure, land-rent income, and local base employment while holding other local fundamentals fixed at their calibrated 1970 levels. We now ask how far these ongoing flows can account for the broader historical association between the US bases and local development.

Our counterfactual analysis suggests that these flows alone may not fully account for the observed association.
To illustrate this point, we compare the changes in the tertiary employment share in the stylized facts and the benchmark counterfactual. In Figure \ref{fig:serv_crosssection1970}, a 10-percentage-point higher base-area share is associated with an approximately 5.2-percentage-point larger increase in the tertiary employment share in 1970, relative to the period without US bases.
In contrast, in the benchmark counterfactual of eliminating US-base sourced consumption demand and employment, the same difference in exposure is associated with a 2.16-percentage-point larger decline in this share (Figure \ref{fig:cf_empshare}).
Therefore, the effects of US-base-sourced consumption demand and employment qualitatively match the observed pattern of service-sector development, but the association might be weaker than in the observed data.
Note that our counterfactual analysis holds local fundamentals fixed at their calibrated 1970 levels and therefore does not evaluate how postwar investment shaped these fundamentals.\footnote{Likewise, the potential impact of other base-related externalities, such as noise and pollution, is also absent in our baseline counterfactual analysis.}
By contrast, the observed data in Figure \ref{fig:serv_crosssection1970} cover changes from the pre-war period (1938), which may reflect both ongoing US-base-sourced consumption demand and employment and changes in local amenities and productivity during the postwar period.

This comparison motivates examining past investment in infrastructure and durable capital as a possible source of the local advantages inherited in 1970. Base development may have encouraged both public infrastructure provision and private investment in durable capital, such as buildings and local human capital. The US bases could thereby have improved subsequent living conditions and production opportunities, consistent with the role of historical agglomeration in shaping later amenities and productivity emphasized by \citet{allen2020persistence}. Related evidence documents lasting local effects of early land-use regulation (\citealt{twinam2018longrun}), transport infrastructure (\citealt{jedwab2016permanent}), and publicly financed industrial investment (\citealt{garin2025long}). In our static model, its calibrated fundamentals, which are held constant in our counterfactual analysis, can embody the legacy of such investments made before 1970.

We examine three municipal measures that capture physical infrastructure and the early commercial opportunities that may have encouraged durable private investment. The first two directly measure infrastructure provision: the share of settlements with piped water around 1963 and a road quality index combining the z-scores of municipal-road improvement, paving, and road density per unit of non-base land. These measures capture infrastructure in place before 1970, the calibration year.
The third measure identifies municipalities containing road corridors designated to host authorized markets serving US personnel in 1949.\footnote{Military Government Directive No.\ 19, issued on September 6, 1949, restricted occupation personnel's purchases of Ryukyuan services and goods not designated as essential commodities to authorized markets along specified highways and at designated places.
% https://www3.archives.pref.okinawa.jp/RDA/USCAR1/0000002175/index.html?page=39
} By concentrating early access to military customers, designation may have encouraged business entry and commercial agglomeration, raising the returns to investment in shops, commercial premises, and related equipment. As in \cite{allen2020persistence}, such investment could have left a durable local advantage beyond the contemporaneous spending that initially attracted businesses. We therefore use designation as an indicator of historical conditions conducive to private capital formation, not as a direct measure of infrastructure or capital stocks. Together, the three measures allow us to examine public infrastructure provision and a potential source of durable private investment within the same historical mechanism.


Table \ref{tab:base_channels_1970_four_channels} examines the first link in this argument: whether these infrastructure and early commercial conditions were more prevalent in municipalities with greater US-base exposure. It reports separate cross-sectional regressions of the three measures on the 1970 base-area share for 35 municipalities defined by their 1970 boundaries. We include region fixed effects and standardized elevation and slope measures to control for topographic conditions.
We find positive and statistically significant associations for all three measures. For a 10-percentage-point higher base-area share, the estimated associations correspond to 0.17, 0.25, and 0.16 outcome standard deviations for piped-water provision, road quality, and corridor designation, respectively. Thus, greater base exposure is associated with better infrastructure and earlier commercial access. 


{\floatsetup[table]{capposition=top}
% Table maintained directly in draft.tex.
\begin{table}[!htbp]
\centering
\caption{U.S. Military-Base Land Share, Infrastructure, and Early Market Access}
\label{tab:base_channels_1970_four_channels}
\footnotesize
\setlength{\tabcolsep}{5pt}
\begin{tabular}{@{}lccc@{}}
\toprule
 & \shortstack{Water-installed\\settlement share} & \shortstack{Road quality\\index} & \shortstack{1949 designated\\road corridor} \\
 & (1) & (2) & (3) \\
\midrule
1970 base share (+10 pp) & 5.686\textsuperscript{**} & 0.565\textsuperscript{**} & 8.063\textsuperscript{**} \\
HC1 robust standard error & (2.210) & (0.209) & (3.502) \\
\addlinespace
Outcome mean & 62.075 & 0.000 & 54.286 \\
Outcome standard deviation & 33.013 & 2.287 & 50.543 \\
Observations & 35 & 35 & 35 \\
\(R^{2}\) & 0.423 & 0.500 & 0.548 \\
Region fixed effects & Yes & Yes & Yes \\
Terrain controls & Yes & Yes & Yes \\
\bottomrule
\end{tabular}
\par\smallskip
\parbox{0.98\linewidth}{\scriptsize \textit{Notes:} The unit of observation is the municipality as of 1970. Each column is estimated separately by OLS with North and South indicators (Central omitted), standardized $\log(1+\text{mean elevation})$, standardized mean slope, and HC1 standard errors. Coefficients report the change associated with a 10-percentage-point increase in 1970 U.S. military-base land share. Coefficients, standard errors, and outcome standard deviations are in percentage points in columns (1) and (3) and index points in column (2); outcome means are in percent in columns (1) and (3) and index points in column (2). The road quality index is the unweighted sum of full-35 sample z-scores for the municipal-road improvement rate, paving rate, and road density per historical non-base km$^2$. The 1949 outcome is a dummy variable for containing roads designated to host authorized markets.  \textsuperscript{***}\(p<0.01\), \textsuperscript{**}\(p<0.05\), \textsuperscript{*}\(p<0.10\).}
\end{table}
}

The next question is whether these observed conditions are also associated with the local fundamentals that we backed out in our calibration. 
For this purpose, we regress log calibrated amenities and productivity on water supply, road quality, and commercial-corridor designation, conditioning on the base-area share and the same region and terrain controls as in Table \ref{tab:base_channels_1970_four_channels}. The regressions ask whether municipalities with similar base exposure and observed geography, but different infrastructure or early commercial access, exhibit advantages in fundamentals in 1970.\footnote{
Methodologically, this approach is closely related to \citet{ang2026amenities}, who relate amenities recovered from quantitative spatial models to observable local characteristics. We apply this logic to both amenities and productivity to investigate whether past infrastructure and durable investment improved both amenities and productivity.
} If these are positively correlated with the calibrated fundamental amenities and productivity, combined with Table \ref{tab:base_channels_1970_four_channels}, we can establish that US bases are associated with better fundamentals through these channels.

{\floatsetup[table]{capposition=top}
% Generated by analysis/amenity_infrastructure_1970/build_paper_main_table.py.
\begin{table}[!htbp]
\centering
\caption{Infrastructure, Early Commercial Access, and Local Fundamentals in 1970}
\label{tab:fundamentals_infrastructure_1970}
\footnotesize
\setlength{\tabcolsep}{2pt}
\begin{tabular*}{\linewidth}{@{\extracolsep{\fill}}lcccccc@{}}
\toprule
 & Agriculture & Manufacturing & \shortstack{Non-\\tradables} & US-base & \shortstack{Private\\mean} & \shortstack{All-four\\mean} \\
 & (1) & (2) & (3) & (4) & (5) & (6) \\
\midrule
\multicolumn{7}{@{}l}{\textit{Panel A: Log amenity}} \\
\addlinespace[2pt]
Base share (+10 pp) & -0.087 & -0.026 & 0.067 & 0.276 & -0.015 & 0.058 \\
 & (0.104) & (0.127) & (0.105) & (0.183) & (0.108) & (0.114) \\
Water share (+10 pp) & 0.065 & 0.155\textsuperscript{**} & 0.143\textsuperscript{**} & 0.049 & 0.121\textsuperscript{**} & 0.103\textsuperscript{*} \\
 & (0.046) & (0.059) & (0.056) & (0.093) & (0.051) & (0.058) \\
Road quality index & -0.021 & -0.059 & -0.018 & -0.235\textsuperscript{**} & -0.033 & -0.083 \\
 & (0.053) & (0.081) & (0.063) & (0.110) & (0.063) & (0.068) \\
Commercial road (1949) & -0.285 & 0.066 & 0.416 & 0.160 & 0.066 & 0.089 \\
 & (0.267) & (0.366) & (0.306) & (0.518) & (0.290) & (0.319) \\
\addlinespace[2pt]
Observations & 35 & 35 & 35 & 35 & 35 & 35 \\
\(R^{2}\) & 0.562 & 0.505 & 0.551 & 0.388 & 0.472 & 0.413 \\
\midrule
\multicolumn{7}{@{}l}{\textit{Panel B: Log productivity}} \\
\addlinespace[2pt]
Base share (+10 pp) & -0.067 & -0.144\textsuperscript{**} & -0.254\textsuperscript{***} & --- & -0.155\textsuperscript{***} & --- \\
 & (0.040) & (0.055) & (0.079) &  & (0.053) &  \\
Water share (+10 pp) & 0.030 & 0.039 & 0.090\textsuperscript{*} & --- & 0.053 & --- \\
 & (0.021) & (0.035) & (0.051) &  & (0.034) &  \\
Road quality index & 0.024 & 0.107\textsuperscript{**} & 0.152\textsuperscript{**} & --- & 0.095\textsuperscript{*} & --- \\
 & (0.028) & (0.044) & (0.073) &  & (0.047) &  \\
Commercial road (1949) & 0.116 & 0.549\textsuperscript{**} & 0.822\textsuperscript{**} & --- & 0.496\textsuperscript{**} & --- \\
 & (0.134) & (0.258) & (0.369) &  & (0.240) &  \\
\addlinespace[2pt]
Observations & 35 & 35 & 35 & --- & 35 & --- \\
\(R^{2}\) & 0.232 & 0.378 & 0.377 & --- & 0.339 & --- \\
\bottomrule
\end{tabular*}
\par\smallskip
\parbox{0.98\linewidth}{\fontsize{8}{9.5}\selectfont \textit{Notes:} All regressions include an intercept, North and South indicators (Central omitted), standardized $\log(1+\text{mean elevation})$, and standardized mean slope. Outcomes are natural logs of calibrated sectoral fundamentals. Private and all-four means are equal-weight geometric means over agriculture, manufacturing, and non-tradables, and over these sectors plus US-base employment, respectively. US-base productivity is not a calibrated outcome, so the corresponding productivity columns are unavailable. Base share and water share coefficients are per 10-percentage-point increase. Water share is the fraction of settlements with water installations in June 1963. Road quality is the sum of full-sample z-scores of municipal-road improvement, paving, and density per historical non-base km$^2$. Commercial road is a dummy variable of containing 1949 road centerlines designated to host authorized markets. HC1 robust standard errors appear in parentheses. \textsuperscript{***}$p<0.01$, \textsuperscript{**}$p<0.05$, \textsuperscript{*}$p<0.10$.}
\end{table}
}

Table \ref{tab:fundamentals_infrastructure_1970} shows that all three channels are positively and significantly associated with at least one dimension of calibrated fundamentals amenities and productivity. Water supply is positively associated with amenities for manufacturing and non-tradable service workers, as well as with the geometric mean across the three private sectors. This is consistent with \citet{devoto2012happiness}, who show that piped-water connections improve well-being by freeing time and reducing conflicts over water. Road quality is positively associated with productivity in manufacturing and non-tradable services. Early commercial-corridor designation is also positively associated with productivity in both sectors and with the private-sector geometric mean. These patterns have intuitive interpretations: better roads can facilitate production and exchange, while early access to military customers may have encouraged investment in durable capital such as buildings and human capital. The persistence of such local advantages is consistent with the role of historical agglomeration emphasized by \citet{allen2020persistence}. Together with Table \ref{tab:base_channels_1970_four_channels}, these results suggest that base development was accompanied by improvements in civilian living and production conditions.\footnote{Conditional on these channels and geographic controls, however, base share is negatively and significantly associated with productivity in manufacturing, non-tradable services, and the geometric mean across the three private sectors. This residual association may reflect restrictions on civilian business expansion and investment near bases beyond the aggregate land-supply constraint captured by the model.}

Taken together, our results suggest that attributing all observed urbanization and service-sector development around the US bases to their consumption demand and employment may overstate the effects of these channels. We provide suggestive evidence that the accompanying investment in infrastructure and durable capital, driven by public investment and early zoning regulation, may also have contributed to local development. This distinction is important for evaluating base expansion or closure in that the gains from additional consumption demand and employment from the bases may differ from what a simple observation of Okinawa in Section \ref{sec:reducedform} suggests: The base expansion might not bring a large economic benefit if it does not accompany infrastructure investment, and the base closure might not have an immediate devastating impact because the associated infrastructure and durable capital do not immediately disappear. 
Indeed, we can find examples of recent policy discussions that recognize the potential benefits of infrastructure investment accompanying base development. For example, in March 2026, Canadian Prime Minister Mark Carney linked new Arctic military hubs and upgrades to existing facilities to investment in energy, telecommunications, and wastewater infrastructure, stating that ``as we build these military operations, we can transform and benefit the surrounding communities.''\footnote{\href{https://www.pm.gc.ca/en/videos/2026/03/12/prime-minister-carney-announces-new-plan-defend-build-and-transform-north}{Prime Minister of Canada, speech transcript, March 12, 2026}, last accessed on October 1, 2026.} In addition to the role of base-related consumption demand and employment, our findings also underscore the potential importance of such infrastructure investment when evaluating the local economic effects of base development.


\section{Conclusion}\label{sec:conclusion}

This paper examines how US-base-related expenditure and employment shaped local economic development in Okinawa. Using newly digitized municipal data, we document ``urbanization without industrialization'' around the bases: greater base exposure is associated with growth in population density on non-base land and in the broad tertiary employment share, including base jobs, but has little association with changes in the manufacturing employment share. Combining base-related GNI components with a quantitative spatial model calibrated to 1970, we quantify the effects of consumption demand and employment. 

Removing base-related consumption, land-rent income, and employment, while holding local fundamentals and Okinawa's workforce fixed and leaving base land unavailable for private use, shifts residents away from more base-exposed municipalities and reallocates employment from tertiary activities toward agriculture and manufacturing. Consumption has larger effects on civilian service employment, whereas base employment has larger effects on population concentration.
In the aggregate, the benchmark counterfactual reduces workers' welfare by 7.25\% and GNI by 11.99\%. These losses are large in absolute value but substantially smaller than the roughly 30\% initial base-related GNI share, with sectoral reallocation cushioning the shock alongside sizable within-island migration. Returning base land dampens population outflows around the bases and reduces the welfare loss to 1.95\%, with little change in sectoral composition. Base-related consumption generates relatively localized service employment gains, while manufacturing losses spread more widely across municipalities. Allowing external migration leaves welfare and sectoral responses close to the benchmark but increases GNI losses through population outflows. Less elastic manufacturing import supply weakens aggregate manufacturing crowding out and can reverse it when import quantities are fixed.
Finally, the consumption and employment channels predict a smaller tertiary-share response than the observed postwar association. Evidence on infrastructure and early commercial access suggests that past investment in infrastructure and durable capital may also have contributed to development around the bases.

As geopolitical tensions renew attention to military basing decisions, our findings highlight the importance of assessing their consequences for host economies. 
Our analysis on Okinawa suggests that consumption spending and job creation around the bases, which many policymakers expect to stimulate the local economy, does induce economic development that accompany a large-scale in-migration and service-sector development. Despite the positive impact, there are two important qualifications. First, the sectoral reallocation away from manufacturing and toward the base employment or service-sector jobs may partially dampen its economic benefits. Second, the impact of consumption spending and employment could be smaller than a naive observation of Okinawa suggests: our analysis also highlights the potential importance of investment in infrastructure and durable capital. Incorporating the impact of base development on public and private investment in a dynamic spatial model, which requires more detailed data on infrastructure and durable capital, remains for future research.



{
%\footnotesize
\bibliographystyle{aer}
\bibliography{reference.bib}
}

\newpage 


\appendix

\pagenumbering{arabic}% resets `page` counter to 1
\renewcommand*{\thepage}{A\arabic{page}}
\def\thesection{\Alph{section}}

\numberwithin{equation}{section}
\numberwithin{figure}{section}
\numberwithin{footnote}{section}
\numberwithin{table}{section}
\setcounter{prop}{0}
% \renewcommand{\thefigure}{A\arabic{figure}}
% \renewcommand{\thepostfigure}{A\arabic{postfigure}}
\setcounter{figure}{0}


\begin{center}
    {\Large \textbf{Appendix to ``Bases and the City: Spatial Economics of US Military Bases in Okinawa''  (Not for Publication)}}
\end{center}


\startcontents[section]
\printcontents[section]{ }{0}{\setcounter{tocdepth}{3}\contentsmargin{2em}}

\newpage 


\section{Details on the data}\label{appsec:datadetails}


\paragraph{Okinawan workers at the US bases.}

Municipal counts of Okinawan base workers by residence and workplace are available in a special 1970 Census tabulation preserved by the \href{https://www2.archives.pref.okinawa.jp/opa/OPA600_RESULT_BUNSYO.aspx?cont_cd=R00008312B}{Okinawa Prefectural Archives}. For 1980 onward, we obtain these counts through \href{https://www.e-stat.go.jp/microdata/jisseki/order/30020020240003}{custom-made Census tabulations}. Both sources are in Japanese (last accessed on May 1, 2025).

\paragraph{US base areas.}

The US base area in 1956 is obtained from \href{https://www2.archives.pref.okinawa.jp/opa/OPA600_RESULT_BUNSYO.aspx?cont_cd=T00017797B&src_keyword=&keyword_hit=&lang=jp}{\textit{Okinawa Nenkan}}, while that during the occupation period is obtained from the Ryukyus (Okinawa) Statistical Yearbook, and that after Okinawa's return to Japan is obtained from \textit{Okinawa-no beigun kichi} and \textit{Okinawa-no beigun oyobi jieitai kichi}.
We also use the shape file of the US bases constructed by the Okinawa Prefecture government, as illustrated in Figure \ref{fig:okinawa_map}. The data are available online at \url{http://gis.pref.okinawa.jp/OpenData/} (in Japanese, last accessed on May 2 2025).


\paragraph{Addressing municipal mergers and divisions.}

We construct municipal crosswalks between years $t_1$ and $t_2$ and aggregate municipalities to common boundaries. For postwar comparisons, these generally follow the later boundaries because of municipal mergers. For prewar--postwar comparisons, we also combine postwar municipalities that belonged to the same prewar municipality.
Two prewar municipalities, \textit{misato-son} and \textit{oozato-son}, were divided around 1950 and subsequently merged into different municipalities. We allocate their 1938 variables, such as population, using the corresponding shares after the division, then aggregate to the boundaries used in each comparison.

\paragraph{Geographic scope and treatment of GNI.}
The GNI statistics cover Okinawa Prefecture. We use them as an approximation for the main island, which accounts for roughly 90\% of the prefecture's population today. Table \ref{tab:data_sources} describes the component shares and their interpolation. Foreign aid from the United States and Japan was about one-fourth of US-base-sourced income around 1970. We exclude this aid from base-related demand and assume that it is spent uniformly across Okinawa.

\paragraph{Summary of the data sources}

For convenience, Table \ref{tab:data_sources} summarizes the dataset used in this paper. 



{
\small
\setlength{\tabcolsep}{3pt}
\renewcommand{\arraystretch}{1.15}

\begin{longtable}{
  @{}
  >{\raggedright\arraybackslash}p{0.14\textwidth}
  >{\raggedright\arraybackslash}p{0.20\textwidth}
  >{\raggedright\arraybackslash}p{0.28\textwidth}
  >{\raggedright\arraybackslash}p{0.34\textwidth}
  @{}
}
\caption{Variables and Data Sources}
\label{tab:data_sources}\\

\toprule
Category
  & Variable or measure
  & Data source(s)
  & Coverage and construction \\
\midrule
\endfirsthead

\multicolumn{4}{c}{
  \tablename\ \thetable\ --- Continued from previous page
}\\
\toprule
Category
  & Variable or measure
  & Data source(s)
  & Coverage and construction \\
\midrule
\endhead

\midrule
\multicolumn{4}{r}{Continued on next page}\\
\endfoot

\bottomrule
\endlastfoot

Population
  & Municipality-level resident population
  & Population Census
  & Available at five-year intervals from 1950 onward.
    Tabulations for the US occupation period were newly digitized
    by the authors. The data cover the Japanese population and
    exclude US citizens. \\

\addlinespace

Employment (residence-based)
  & Residents employed by industry
  & Population Census; \textit{Okinawa Prefectural Statistical
    Yearbook} (\textit{Okinawa-ken toukei sho}) for 1934 and 1938
  & Used to construct residence-based sectoral employment shares.
    The reduced-form analysis groups industries into agriculture,
    manufacturing (including construction), and services (including
    base employment). The model treats base employment separately by utulizing the separate base-employment data below. \\

\addlinespace

Employment (workplace-based)
  & Number of workers by industry and
    municipality of workplace
  & 1970 Population Census
  & This is a workplace-based measure and may differ from
    residence-based employment because of commuting across
    municipalities. These data provide the sectoral employment
    distribution used in the 1970 calibration. \\

\addlinespace

US-base employment
  & Okinawan workers employed at US military bases, by
    municipality of residence and workplace
  & Population Census; a special 1970 Census tabulation preserved
    by the Okinawa Prefectural Archives; custom-made Census
    tabulations provided by the Ministry of Internal Affairs and
    Communications from 1980 onward
  & These data identify US-base workers separately from other
    sectors, by both residence and workplace. \\

\addlinespace

US-base size
  & US-base area and the share of municipal land occupied by
    US bases
  & \textit{Okinawa Nenkan} for 1956; \textit{Ryukyus Statistical
    Yearbook} and \textit{Okinawa Statistical Yearbook} for the
    occupation period; \textit{Okinawa-no beigun kichi} and
    \textit{Okinawa-no beigun oyobi jieitai kichi} for the
    post-reversion period
  & The municipal base-area share is the main measure of exposure
    to US-base-sourced consumption and employment. \\

\addlinespace

US-base geography
  & Geographic boundaries of US military bases
  & GIS shapefiles provided by the Okinawa Prefectural Government
  & Used to identify and map the locations of bases. For Figure \ref{fig:okinawa_map},
    the land boundary of Okinawa Island is taken from the Digital
    National Land Information, using the 1975 geography as an
    approximation for 1970. \\

\addlinespace

Alternative measures of base importance
  & Number of Okinawan workers and US military personnel at
    each base
  & Komeito (1969), \textit{Okinawa beigun kichi no jittai chosa}
    [Survey on the Actual Conditions of US Military Bases in
    Okinawa]
  & Used to validate the base-area measure and in supplementary
    robustness analyses. Base-level counts are aggregated to
    municipalities and divided by total municipal land area,
    including land occupied by US bases. \\

\addlinespace

Aggregate income
  & Gross National Income of Okinawa
  & \textit{Ryukyus Statistical Yearbook}
    (\textit{Ryukyus Toukei Nenkan}) and
    \textit{Okinawa Statistical Yearbook}
    (\textit{Okinawa Toukei Nenkan})
  & Available annually from 1955 onward. %The statistics refer to Okinawa Prefecture rather than Okinawa Island alone. 
    \\

\addlinespace

US-base-sourced income
  & Share of Okinawa's GNI derived from US bases, separately for purchases of goods and services; salaries of Okinawan base
    workers; and rents paid for base land
  & \textit{Ryukyus Statistical Yearbook} and
    \textit{Okinawa Statistical Yearbook}
  & Includes base-related payments financed by the Japanese
    government. Component shares for 1966--1971 are linearly
    interpolated. Economic aid and intergovernmental transfers,
    including base-related grants to local governments, are excluded.\\

\addlinespace

Household expenditure
  & Restaurant expenditure share in Naha relative to Okinawa as a whole
  & 1970 Okinawa Household Expenditure Survey
    (National Diet Library Digital Collections:
    \url{https://dl.ndl.go.jp/pid/12101674})
  & The observed relative expenditure share, 1.16, is used as the
    calibration target for the elasticity of substitution between
    tradable and non-tradable goods. \\

\addlinespace

Model validation
  & Bilateral commuting probabilities in 1967 and municipality-level house prices
  & 1967 commuting origin--destination data; 1970 edition of
    \textit{Shichoson koufuzei seido kaisetsu}
  & These variables are not calibration targets. Commuting probabilities
    are compared across aggregated origin--destination pairs. House prices
    and model-implied floor-space rents are standardized before comparison. \\

\addlinespace

Agricultural trade
  & Okinawa's agricultural-goods trade balance around 1970
  & \textit{Okinawa Statistical Yearbook}
  & Used in the external-trade extension to discipline the manufacturing
    share in the tradable-goods composite. \\

\addlinespace

Travel behavior
  & Trip origins, destinations, purposes, and workers' industries at the
    individual-trip level
  & 1977 Okinawa Main Island Central and Southern Urban Area
    Person Trip Survey, provided by the Okinawa Prefectural
    Government
  & Trip purposes allow commuting and shopping trips to be
    identified, and worker industries allow commuting flows to be
    constructed separately by sector. The 1977 survey does not cover northern Okinawa
    Island. The corresponding 1989 and 2006 surveys are used for
    robustness checks. \\

\addlinespace

Commuting and shopping flows
  & Bilateral commuting matrices for agriculture, manufacturing, and
    non-tradable services, and a bilateral shopping matrix between municipalities
  & Authors' construction using the 1977 Person Trip Survey
  & The sectoral commuting matrices are pooled to estimate the
    sector-conditioned commuting gravity equation. Shopping and leisure trips are used in
    the calibration of shopping travel costs. \\

\addlinespace

Travel distance
  & Bilateral road distance between municipalities
  & Newly digitized 1973 road network based on maps from the
    Geospatial Information Authority of Japan. 
    1930 and 2005 road networks are also digitized for robustness.
  & Distances are measured along the road network between
    municipal centroids and are used in the gravity-equation
    estimation. Bilateral distances change little when the 1973 network is replaced by the 1930 or 2005 network. \\

\addlinespace

Water infrastructure
  & Share of settlements with installed water facilities
  & \textit{Ryukyus Statistical Yearbook}, 8th ed. (1963),
    Table 323
  & Measured as of June 30, 1963 and constructed as the number of
    settlements with installed water facilities divided by the total
    number of settlements. Predecessor municipalities are aggregated to
    the fixed 1970 municipal units. \\

\addlinespace

Road infrastructure
  & Municipal-road improvement rate, paving rate, road density outside
    base land, and the road quality index
  & Okinawa Prefectural Archives, \textit{Doro ni kansuru shorui}
    [Documents Concerning Roads] series; \textit{Kensetsu unyu yoran
    1965} for aggregate cross-checks
  & Municipal-road ledgers refer mainly to 1963--64 and exclude military
    and Ryukyu Government roads. The first two measures divide improved
    or paved length by total municipal-road length; density divides total
    road length by 1954 municipal area net of March 1956 base land. The
    index is the unweighted sum of the full-sample z-scores of these three
    measures. \\

\addlinespace

Early market access
  & Indicator for a 1949 road corridor designated to host markets
    authorized to transact with US occupation personnel
  & US Military Government Directives Nos. 38 (1948), 19 (1949), and
    21 (1949); US Navy, \textit{Okinawa Shima Road Map} (1945);
    authors' road-to-municipality crosswalk
  & The indicator equals one when the reconstructed centerline of a
    designated road crosses the interior of a fixed 1970 municipality
    (or for the separately designated Naha establishment). It represents
    access to authorized markets, not a statutory land-use zone. \\

\addlinespace

Terrain controls
  & Municipality-level mean elevation and mean slope
  & Ministry of Land, Infrastructure, Transport and Tourism,
    Digital National Land Information: administrative-area data
    (N03, 1950 and 1975) and elevation and slope fourth-mesh data
    (G04-c, 2011 release)
  & The approximately 500-meter grid cells are intersected with the
    reconstructed fixed 1970 municipal polygons and aggregated using
    intersection-area weights. In the infrastructure and fundamentals
    regressions, $\log(1+\text{mean elevation})$ and mean slope are
    standardized across the 35 municipalities. \\



\end{longtable}
}



\clearpage

\section{Supplementary results on stylized facts}\label{appsec:reducedform}

\subsection{Overview of supplementary results on stylized facts}\label{sec:reduced_additional}

\paragraph{Levels in 1970.} Figure \ref{fig:reduced_crosssection1970_level} reproduces Facts \ref{result1} and \ref{result2} using 1970 levels instead of changes from 1938.

\paragraph{Endogeneity of base location.} Table \ref{tab:base-location-ols-iv} shows that the patterns persist after controlling for terrain and instrumenting base exposure with early occupation near the US landing point.

\paragraph{Prewar trends.} Figure \ref{fig:reduced_crosssection1938} examines changes from 1934 to 1938. The population-density coefficient is positive ($p=0.057$), with a cumulative point estimate about 3\% of the postwar coefficient, but the magnitude is much smaller than the observed association. Sectoral employment-share coefficients are not statistically significant.

\paragraph{Base returns.} Using the 1956 base-area share as an instrument yields the same qualitative conclusions (Figure \ref{fig:reduced_iv}).

\paragraph{Alternative exposure measures.} Replacing base-area shares with the densities of Okinawan base workers or US military personnel preserves the positive pre-reversion associations with population density and services (Figures \ref{fig:reduced_japanese} and \ref{fig:reduced_military}).

\paragraph{Patterns over time.} Figure \ref{fig:reduced_persistence} shows similar patterns across years, with smaller coefficients after reversion in 1972 as the bases' GNI share declined.

\paragraph{Service subsectors.} Table \ref{tab:detailsector1970} shows that wholesale and retail and intangible services have the strongest positive associations with base exposure.

\clearpage

\subsection{Stylized Facts for the 1970 level variables}\label{sec:facts_1970level}

\begin{figure}[H]
    \subfloat[Population density\label{fig:popdens_crosssection1970_level}]{%
      \includegraphics[width=0.34\textwidth, page=5]{figure/postlog3870_scatter_areaUS.pdf}
    }
    \hfill
    \subfloat[Population level\label{fig:pop_crosssection1970_level}]{%
      \includegraphics[width=0.34\textwidth, page=4]{figure/postlog3870_scatter_areaUS.pdf}
    } \\  
    \subfloat[Service share\label{fig:serv_crosssection1970_level}]{%
      \includegraphics[width=0.34\textwidth, page=1]{figure/postlog3870_scatter_areaUS.pdf}
    }
    \hfill
    \subfloat[Manufacturing share\label{fig:manu_crosssection1970_level}]{%
      \includegraphics[width=0.34\textwidth, page=2]{figure/postlog3870_scatter_areaUS.pdf}
    }
    \hfill \\
    \subfloat[Agriculture share\label{fig:agri_crosssection1970_level}]{%
      \includegraphics[width=0.34\textwidth, page=3]{figure/postlog3870_scatter_areaUS.pdf}
    }
    \vspace{0.2in}
    \caption{Relationship between US bases, Urbanization, and Industry Structure (1970, without differencing)}
    \vspace{-0.2in}
    \floatfoot{\textbf{Note}: Each point represents a municipality. The horizontal axis is the 1970 US-base area share. The outcomes are log population density, log population, and the employment shares of services, manufacturing, and agriculture. Population density uses municipal land area outside US bases. Lines are OLS fits.
    }
    \label{fig:reduced_crosssection1970_level}
\end{figure}


\clearpage

% BEGIN added base-location endogeneity analysis
% Additional material only; included at the end of the supplementary reduced-form results.
\subsection{Endogeneity of US base location}\label{appsec:base_location_endogeneity}

\begin{figure}[!htbp]
  \centering
  \includegraphics[width=0.96\linewidth]{figure/iv30km_early_occupation_map.pdf}
  \caption{US landing areas and the early-occupation instrument}
  \label{fig:base_location_iv}
  \floatfoot{\textbf{Note}: The star marks the US landing reference point; purple lines show alternative landing sectors. Dark teal denotes early-occupation municipalities with a mean distance below 30 km from the landing point, and light green denotes the remaining early-occupation municipalities. Boundaries are harmonized between 1938 and 1970; only the main island is shown. Sources: project GIS data and the historical accounts cited below. Place names are from the \href{https://doi.org/10.20676/00000448}{CODH historical place-name dataset} (CC BY 4.0) and the Okinawa Prefectural Archives.}
\end{figure}
\clearpage

Figure \ref{fig:base_location_iv} shows that the Hagushi beaches, where US forces landed on April 1, 1945, were not the only landing areas considered. If the west-coast landing proved impracticable, the alternative plan called for Marines to land between Minatoga and Chinen Point, followed by Army divisions between Kuba and Yonabaru \citeapp[pp.~32--34]{ApplemanEtAl1948Okinawa_app}. These were successive components of the same alternative plan. In assessing this alternative, \citeapp{ApplemanEtAl1948Okinawa_app} write (p.~34):
\begin{quote}
\small
``Although the alternate plan met most of the requirements for a successful landing operation, it was distinctly a second choice because it would allow the enemy reserves to offer maximum opposition to the second landings and would require a prolonged assault against all the enemy forces on the island to complete the first phase of the mission.''
\end{quote}
The alternative was less attractive because of anticipated Japanese resistance. The Marine Corps history also describes the east-coast alternative and the weather concerns that prompted it \citeapp[pp.~23--25]{NicholsShaw1955Okinawa_app}. These accounts emphasize military and logistical considerations, including airfield access, unloading capacity, and the expected enemy response.

The colored areas in Figure \ref{fig:base_location_iv} were occupied within about a week of the landing, whereas resistance delayed occupation of southern Okinawa for several months. The US began building bases in the early-occupation area during the battle, partly to support air raids on mainland Japan. Many remained after the war, making early occupation a predictor of subsequent base exposure.

To implement this idea, we construct an instrumental variable as follows. 
We combine the early-occupation classification shown in Figure \ref{fig:base_location_iv} with proximity to the actual landing point:
\begin{equation}
  Z_i = \mathrm{Early Occupation}_i\,\mathbf{1}\{\bar d_i<30\text{ km}\}.
  \label{eq:base_location_instrument}
\end{equation}
Here, $\bar d_i$ is the distance from the landing reference point to the municipal representative point, averaged across constituent 1970 municipalities when harmonizing boundaries. $\mathrm{Early Occupation}_i$ equals one if municipality $i$ was occupied during the early stage of the battle. The instrument combines proximity to the landing point with occupation timing. We control for municipal mean altitude and slope in both stages.

\clearpage
% Numerical entries copied verbatim from the Stata-generated six-column table.
\begin{table}[H]
\centering
\begin{threeparttable}
\caption{Base exposure, population density, and sectoral employment: OLS and IV}
\label{tab:base-location-ols-iv}
\small
\setlength{\tabcolsep}{3pt}
\def\sym#1{\ifmmode^{#1}\else\(^{#1}\)\fi}
\begin{tabular*}{\linewidth}{@{\extracolsep{\fill}}l*{6}{c}@{}}
\toprule
&\multicolumn{3}{c}{OLS}&\multicolumn{3}{c}{IV (2SLS)}\\
\cmidrule(lr){2-4}\cmidrule(lr){5-7}
& (1) & (2) & (3) & (4) & (5) & (6)\\
& \shortstack{Population\\density} & \shortstack{Service\\share} & \shortstack{Manufacturing\\share}
& \shortstack{Population\\density} & \shortstack{Service\\share} & \shortstack{Manufacturing\\share}\\
\midrule
Base area share (1970)    &        2.4584\sym{***}&        0.4050\sym{***}&       -0.0500         &        2.3998\sym{***}&        0.6598\sym{***}&        0.0204         \\
                          &      (0.4322)         &      (0.0760)         &      (0.0429)         &      (0.4483)         &      (0.1849)         &      (0.0872)         \\
Mean altitude (m)         &        0.0034         &        0.0004         &       -0.0008         &        0.0036         &       -0.0006         &       -0.0011\sym{*}  \\
                          &      (0.0041)         &      (0.0008)         &      (0.0005)         &      (0.0039)         &      (0.0009)         &      (0.0006)         \\
Mean slope                &       -0.1316\sym{**} &       -0.0249\sym{*}  &        0.0096         &       -0.1356\sym{**} &       -0.0076         &        0.0144         \\
                          &      (0.0550)         &      (0.0126)         &      (0.0067)         &      (0.0540)         &      (0.0160)         &      (0.0092)         \\
Constant                  &        0.7066\sym{***}&        0.4658\sym{***}&        0.1838\sym{***}&        0.7270\sym{***}&        0.3771\sym{***}&        0.1592\sym{***}\\
                          &      (0.1640)         &      (0.0596)         &      (0.0356)         &      (0.1935)         &      (0.0834)         &      (0.0482)         \\
\midrule
Observations              &            29         &            29         &            29         &            29         &            29         &            29         \\
First-stage KP F          &                       &                       &                       &         17.27         &         17.27         &         17.27         \\
\bottomrule
\end{tabular*}
\begin{tablenotes}[flushleft]
\footnotesize
\item \textit{Notes:} Outcomes are 1938--1970 changes in log population density and the service (broad tertiary) and manufacturing (broad secondary) employment shares. All columns use 29 harmonized municipalities and control for mean altitude and slope. Columns (4)--(6) instrument the 1970 base-area share with early occupation interacted with mean distance below 30 km from the landing point; both stages include the terrain controls. KP F is the Kleibergen--Paap first-stage F statistic. Robust standard errors are in parentheses. $^{*}p<0.10$, $^{**}p<0.05$, $^{***}p<0.01$.
\end{tablenotes}
\end{threeparttable}
\end{table}

Table \ref{tab:base-location-ols-iv} reports the results for population density growth and changes in the service and manufacturing employment shares. Columns (1)--(3) report the OLS results, while columns (4)--(6) report two-stage least squares estimates using the variable \eqref{eq:base_location_instrument} as an instrument.
The OLS estimates are 2.4584 for population density growth and 0.4050 for the change in the service employment share. The corresponding IV estimates are 2.3998 and 0.6598, implying that the population density estimate is almost unchanged under IV, while the service-share estimate is larger.
For the manufacturing employment share, the OLS and IV estimates are $-0.0500$ and $0.0204$, respectively, and neither is statistically significant. The first-stage Kleibergen--Paap F statistic is 17.27 in all three IV columns.
These results indicate that the positive association of US bases with urbanization and service-sector expansion survives, while neither specification provides evidence of an increase in the manufacturing employment share.


\clearpage

\subsection{Testing the pre-trends}
Figure \ref{fig:reduced_crosssection1938} relates changes from 1934 to 1938 to the 1956 US-base area share. The coefficient for log population-density growth is $0.096$ (standard error $0.049$, $p=0.057$). Because municipal area is constant across the prewar years, log population and population-density changes are identical.

The cumulative prewar coefficient is about 3\% of the postwar population-density coefficient of $3.045$ in Figure \ref{fig:popdens_crosssection1970}. A 10-percentage-point difference in base share is associated with an approximately 1\% higher ratio of end-period to initial density in 1934--1938, compared with 36\% in 1938--1970. Postwar density changes reflect both population change and the reduction in civilian land caused by base occupation.\footnote{The periods span four and 32 years. Annualized log-change coefficients are $0.024$ and $0.095$, respectively, so the prewar coefficient is about 25\% of the postwar coefficient. The prewar regressions use 39 municipalities and the 1956 base-area share; the postwar regressions use 29 aggregated municipalities and the 1970 share.}

The prewar coefficients for service, manufacturing, and agricultural employment shares are $0.029$ (standard error $0.032$), $0.004$ ($0.008$), and $-0.034$ ($0.031$), respectively. None is statistically significant at the 10\% level.

\begin{figure}[p]
    \subfloat[Population density\label{fig:popdens_crosssection1938}]{%
      \includegraphics[width=0.38\textwidth, page=5]{figure/loggrowth3438_scatter_areaUS.pdf}
    }
    \hfill
    \subfloat[Service share\label{fig:serv_crosssection1938}]{%
      \includegraphics[width=0.38\textwidth, page=1]{figure/sharediff3438_scatter_areaUS.pdf}
    } \\
    \subfloat[Manufacturing share\label{fig:manu_crosssection1938}]{%
      \includegraphics[width=0.38\textwidth, page=2]{figure/sharediff3438_scatter_areaUS.pdf}
    }
    \hfill
    \subfloat[Agriculture share\label{fig:agri_crosssection1938}]{%
      \includegraphics[width=0.38\textwidth, page=3]{figure/sharediff3438_scatter_areaUS.pdf}
    }
    \vspace{0.15in}
    \caption{Prewar trends, 1934--1938}
    \vspace{-0.1in}
    \floatfoot{\textbf{Note}: Each point represents one of 39 municipalities. The horizontal axis is the 1956 US-base area share. Outcomes are 1934--1938 changes in log population density and sectoral employment shares; a share change of 0.01 equals one percentage point. Lines are OLS fits, and boxes report slopes and conventional standard errors.}
    \label{fig:reduced_crosssection1938}
\end{figure}



\clearpage


\subsection{Addressing endogeneity of US base returns} We instrument each year's base-area share with its 1956 value to address selective base returns. Figure \ref{fig:reduced_iv} repeats Figure \ref{fig:reduced_persistence} using two-stage least squares. The estimates preserve Facts \ref{result1} and \ref{result2}.\footnote{First-stage F statistics range from 27.63 in 2015 to 226.16 in 1965.}

\begin{figure}[p]
    \subfloat[Population density\label{fig:popdens_reduced_iv}]{%
      \includegraphics[width=0.375\textwidth, page=5]{figure/coef_plots_iv.pdf}
    }
    \hfill
    \subfloat[Population level\label{fig:pop_reduced_iv}]{%
      \includegraphics[width=0.375\textwidth, page=4]{figure/coef_plots_iv.pdf}
    } \\
    \subfloat[Service share\label{fig:serv_reduced_iv}]{%
      \includegraphics[width=0.375\textwidth, page=1]{figure/coef_plots_iv.pdf}
    }
    \hfill
    \subfloat[Manufacturing share\label{fig:manu_reduced_iv}]{%
      \includegraphics[width=0.375\textwidth, page=2]{figure/coef_plots_iv.pdf}
    } 
    \hfill \\ 
    \subfloat[Agriculture share\label{fig:agri_reduced_iv}]{%
      \includegraphics[width=0.375\textwidth, page=3]{figure/coef_plots_iv.pdf}
    }
    \vspace{0.2in}
    \caption{Urbanization without industrialization around the US bases (IV estimation)}
    \vspace{-0.2in}
    \floatfoot{\textbf{Note}: Each panel reports year-specific two-stage least squares coefficients on the contemporaneous US-base area share, instrumented by its 1956 value. Outcomes and confidence intervals follow Figure \ref{fig:reduced_persistence}.
    }
    \label{fig:reduced_iv}
\end{figure}


\clearpage



\subsection{Using alternative measures of base importance} 
Instead of the share of land occupied by US bases in a municipality, we use alternative measures of base importance: the densities of Okinawan base workers and US military personnel.
\citeapp{Komeito69_app} reports the number of Okinawan workers and US military personnel at each base. We aggregate these counts to municipalities and divide by total municipal land area, including base land. Figures \ref{fig:reduced_japanese} and \ref{fig:reduced_military} use the resulting densities. Both measures are positively associated with population density and service employment shares before reversion in 1972.

\begin{figure}[p]
    \subfloat[Population density\label{fig:popdens_reduced_japanese}]{%
      \includegraphics[width=0.375\textwidth, page=10]{figure/coef_plots.pdf}
    }
    \hfill
    \subfloat[Population level\label{fig:pop_reduced_japanese}]{%
      \includegraphics[width=0.375\textwidth, page=9]{figure/coef_plots.pdf}
    } \\
    \subfloat[Service share\label{fig:serv_reduced_japanese}]{%
      \includegraphics[width=0.375\textwidth, page=6]{figure/coef_plots.pdf}
    }
    \hfill
    \subfloat[Manufacturing share\label{fig:manu_reduced_japanese}]{%
      \includegraphics[width=0.375\textwidth, page=7]{figure/coef_plots.pdf}
    } \\
    \subfloat[Agriculture share\label{fig:agri_reduced_japanese}]{%
      \includegraphics[width=0.375\textwidth, page=8]{figure/coef_plots.pdf}
    }
    \vspace{0.2in}
    \caption{Urbanization without industrialization and the US bases (Okinawan base-worker density)}
    \vspace{-0.2in}
    \floatfoot{\textbf{Note}: The explanatory variable is the number of Okinawan base workers per square kilometer of total municipal land area, using counts from \citeapp{Komeito69_app}. Outcomes and confidence intervals follow Figure \ref{fig:reduced_persistence}.}
    \label{fig:reduced_japanese}
\end{figure}



\begin{figure}[p]
    \subfloat[Population density\label{fig:popdens_reduced_military}]{%
      \includegraphics[width=0.375\textwidth, page=15]{figure/coef_plots.pdf}
    }
    \hfill
    \subfloat[Population level\label{fig:pop_reduced_military}]{%
      \includegraphics[width=0.375\textwidth, page=14]{figure/coef_plots.pdf}
    } \\
    \subfloat[Service share\label{fig:serv_reduced_military}]{%
      \includegraphics[width=0.375\textwidth, page=11]{figure/coef_plots.pdf}
    }
    \hfill
    \subfloat[Manufacturing share\label{fig:manu_reduced_military}]{%
      \includegraphics[width=0.375\textwidth, page=12]{figure/coef_plots.pdf}
    } \\
    \subfloat[Agriculture share\label{fig:agri_reduced_military}]{%
      \includegraphics[width=0.375\textwidth, page=13]{figure/coef_plots.pdf}
    }
    \vspace{0.2in}
    \caption{Urbanization without industrialization and the US bases (US military personnel density)}
    \vspace{-0.2in}
    \floatfoot{\textbf{Note}: The explanatory variable is the number of US military personnel per square kilometer of total municipal land area, using counts from \citeapp{Komeito69_app}. Outcomes and confidence intervals follow Figure \ref{fig:reduced_persistence}.}
    \label{fig:reduced_military}
\end{figure}



\clearpage

\subsection{The association with US bases in different years}

Figure \ref{fig:reduced_persistence} reports OLS estimates of $\beta_t$ from
\[
  y_{it}-y_{i,1938}=\beta_t\text{USBaseShare}_{it}+\text{constant}+\text{error},
\]
for $t=1960,1965,\ldots,2015$. Here, $\text{USBaseShare}_{it}$ is the municipal base-area share and $y_{it}$ is the outcome of interest.
Population, population density, and service employment shares are positively associated with base exposure in all years. The coefficients become smaller after reversion in 1972, when base-related income fell as a share of GNI (Figure \ref{fig:okinawa_GDPshare}). The service-share pattern mainly reflects growth in municipalities away from the bases. Manufacturing shares show little association with base exposure, while agricultural shares are negatively associated with it.

\begin{figure}[p]
    \subfloat[Population density\label{fig:popdens_persistence}]{%
      \includegraphics[width=0.375\textwidth, page=5]{figure/coef_plots.pdf}
    }
    \hfill
    \subfloat[Population level\label{fig:pop_persistence}]{%
      \includegraphics[width=0.375\textwidth, page=4]{figure/coef_plots.pdf}
    } \\  
    \subfloat[Service share\label{fig:serv_persistence}]{%
      \includegraphics[width=0.375\textwidth, page=1]{figure/coef_plots.pdf}
    } 
    \hfill 
    \subfloat[Manufacturing share\label{fig:manu_persistence}]{%
      \includegraphics[width=0.375\textwidth, page=2]{figure/coef_plots.pdf}
    } 
    \hfill  \\
    \subfloat[Agriculture share\label{fig:agri_persistence}]{%
      \includegraphics[width=0.375\textwidth, page=3]{figure/coef_plots.pdf}
    }
    \vspace{0.2in}
    \caption{Urbanization without industrialization around the US bases: Long-run Perspective}
    \vspace{-0.2in}
    \floatfoot{\textbf{Note}: Each panel reports year-specific OLS coefficients on the contemporaneous US-base area share. Outcomes are changes from 1938 in log population density, log population, and sectoral employment shares. Postwar population density uses land outside US bases; prewar density uses total municipal land area. Black bars are 95\% confidence intervals.
    }
    \label{fig:reduced_persistence}
\end{figure}

\clearpage

\subsection{Disaggregating service-sector employment}

We examine 1970 employment shares in wholesale and retail, finance and insurance, real estate, transportation and communications, utilities, intangible services, public administration, and unclassified jobs. Since comparable prewar categories are unavailable, we estimate a cross-sectional regression for each sector $j$:

\begin{align}\label{eq:regeq_diff_detailsector}
  y_{ij1970} =  \beta_j \text{USBaseShare}_{i1970} + \text{constant} + \text{error term}.
\end{align}
Here, $y_{ij1970}$ is employment in sector $j$ divided by total employment across all sectors in municipality $i$ in 1970, and $\beta_j$ is the sector-specific coefficient.

Table \ref{tab:detailsector1970} shows the strongest positive associations for wholesale and retail and intangible services, consistent with demand for consumer services around bases.\footnote{\citeapp{fan2023india_app} distinguish consumer and business services in their analysis of Indian growth.} Finance and insurance, real estate, and utilities also have positive, statistically significant coefficients. Transportation and communications and public administration do not.

\begin{table}[H] \centering 
\resizebox{\textwidth}{!}{ 
\begin{tabular}{@{\extracolsep{5pt}}lcccccccc} 
\\[-1.8ex]\hline 
\hline \\[-1.8ex] 
 & \multicolumn{8}{c}{Employment share in total municipal employment} \\ 
\cline{2-9} 
 & \shortstack{Wholesale\\and retail} & \shortstack{Finance\\and insurance} & \shortstack{Real\\estate} & \shortstack{Transportation\\and communications} & Utilities & \shortstack{Intangible\\service} & \shortstack{Public\\administration} & Unclassified \\ 
 & (1) & (2) & (3) & (4) & (5) & (6) & (7) & (8)\\ 
\hline \\[-1.8ex] 
US base area share & 0.214$^{***}$ & 0.009$^{*}$ & 0.004$^{***}$ & $-$0.035 & 0.013$^{***}$ & 0.339$^{***}$ & $-$0.007 & 0.0001 \\ 
  & (0.047) & (0.005) & (0.001) & (0.023) & (0.003) & (0.054) & (0.008) & (0.0001) \\ 
  & & & & & & & & \\ 
\hline \\[-1.8ex] 
Average employment share & 0.162 & 0.009 & 0.002 & 0.056 & 0.006 & 0.220 & 0.038 & 0.000 \\ 
Observations & 29 & 29 & 29 & 29 & 29 & 29 & 29 & 29 \\ 
R$^{2}$ & 0.434 & 0.111 & 0.318 & 0.082 & 0.442 & 0.593 & 0.027 & 0.061 \\ 
\hline 
\hline \\[-1.8ex] 
 \multicolumn{9}{l}{$^{*}$p$<$0.1; $^{**}$p$<$0.05; $^{***}$p$<$0.01} \\ 
\end{tabular} 
}
  \caption{Service employment share by sector and the US base area share (1970)} 
  \label{tab:detailsector1970} 
    \floatfoot{\textit{Notes:} The table reports OLS estimates of equation (\ref{eq:regeq_diff_detailsector}) for 29 municipalities harmonized between 1938 and 1970. Outcomes are sectoral employment shares in 1970, measured relative to total municipal employment. Conventional standard errors are in parentheses.}
\end{table} 


\clearpage
% END added base-location endogeneity analysis

\section{Calibration details}\label{appsec:calibration}

\subsection{Supplementary calibration details}
\label{appsec:calibration_supplement}

\paragraph{Demand weights.}
A more general version of the price index in equation (\ref{eq:price_index}) is
\begin{equation*}
    \mathbb{P}_{i}=
    \left[
        \beta_T(P_i^T)^{1-\kappa}
        +\beta_N(\mathbb{P}_{i}^N)^{1-\kappa}
    \right]^{\frac{1}{1-\kappa}},
\end{equation*}
where $\beta_T,\beta_N>0$ are preference weights for tradable goods and non-tradable services.
These weights are not separately identified from sectoral productivity in our calibration.
Following \citeapp{faber2019tourism_app}, we set $\beta_T=\beta_N=1$; the calibrated productivity terms therefore also absorb the demand weights.

\paragraph{Wage scales.}
Let $L_s\equiv\sum_i L_{is}$ and $\lambda_s\equiv\sum_i\sum_j\lambda_{ijs}$.
Given the private-sector amenity normalizations in Step 3 of Section \ref{sec:calibration}, we normalize the geometric mean of $\omega_{iA}$ to one and match the manufacturing and service wage scales to observed employment shares.
Using equation (\ref{eq:prob_live_work}), the restriction $\lambda_s/\lambda_A=L_s/L_A$ becomes
\begin{equation}
\label{eq:emp_ratio_AM}
    \frac{L_s}{L_A}
    =\frac{\sum_i\sum_j\widetilde{A}_{is}\omega_{js}d_{ij}^{-\vartheta}}
    {\sum_i\sum_j\widetilde{A}_{iA}\omega_{jA}d_{ij}^{-\vartheta}},
    \qquad s\in\{M,N\}.
\end{equation}
Appendix \ref{appsec:calibration_GNI} applies the analogous normalization to the US-base sector using its primed variables.

\paragraph{Service-price recovery.}
For a given $\kappa$, the service-price vector $\{P_i^N\}_{i=1}^I$ determines all the price indexes $\mathbb{P}_k^N$ through equation (\ref{eq:price_index_N}).
Under the benchmark normalization $P^T=1$, equations (\ref{eq:price_index}) and (\ref{eq:share_N}) also imply
\begin{equation*}
    \chi_k^N=
    \frac{\alpha\left[\sum_l\xi_{kl}(P_l^N)^{1-\sigma}\right]^{(1-\kappa)/(1-\sigma)}}
    {1+\left[\sum_l\xi_{kl}(P_l^N)^{1-\sigma}\right]^{(1-\kappa)/(1-\sigma)}}.
\end{equation*}
Substituting these functions into the service-market clearing condition (\ref{eq:goods_market}) leaves the service-price vector to be recovered, conditional on the other calibrated inputs.

\paragraph{Housing expenditure share.}
The housing expenditure share $1-\alpha=0.25$ used in Step 1 of Section \ref{sec:calibration} is also consistent with the \href{https://www.e-stat.go.jp/stat-search/files?page=1&layout=datalist&toukei=00200564&tstat=000000640001&cycle=0&tclass1=000000640041&tclass2=000000640134&tclass3=000000640135&iroha=9&tclass4val=0}{1999 National Survey of Family Income and Expenditure for Okinawa Prefecture} (in Japanese; accessed May 16, 2025).

\clearpage

\subsection{Gravity equation estimation}\label{appsec:gravity}

We estimate commuting and shopping gravity equations using the 1977 Okinawa person trip survey and intermunicipal distances from the digitized 1973 road network. For within-municipality trips, we use the expected distance between two points drawn independently and uniformly from a disk with municipality $i$'s area, $Area_i$: $d_{ii}=\frac{128}{45\pi}\sqrt{\frac{Area_i}{\pi}}$. Commuting costs take the form $\tau_{ij}=d_{ij}^{\zeta}$.

\paragraph{Commuting gravity equation.}

We classify workers into agriculture, manufacturing, and non-tradable services using their reported industry and pool the three municipality-to-municipality commuting flow matrices. We estimate the multiplicative counterpart of equation (\ref{eq:commuting_gravity_sector}), with origin-by-industry and destination-by-industry fixed effects. Table \ref{tab:gravity_commutingDT} reports distance coefficients of $-1.975$ under PPML and $-1.714$ under log OLS, both significant at the 1 percent level. We set $\vartheta=1.85$, approximately the midpoint of their absolute values. Given $\theta=2.19$, the commuting cost elasticity is $\zeta=\vartheta/\theta\simeq0.845$.

\paragraph{Shopping gravity equation.}

Using home-originating shopping and leisure trips from the 1977 person trip survey, we estimate
\begin{equation}
\label{eq:shoppingprob_gravity}
    \ln \xi_{ij}^{\mathrm{trip}} = -\psi_S \ln d_{ij} + \Psi_i^R + \Psi_j^D ,
\end{equation}
where $\Psi_i^R$ and $\Psi_j^D$ are origin and destination fixed effects. We also estimate the specification without these fixed effects.

Table \ref{tab:gravity_shopping} reports PPML estimates using destination shares, including zeros, and log-OLS estimates using positive shares. We set $\psi_S=2.3$ from the PPML estimate without fixed effects in Column 1 ($-2.2978$) and construct the model's shopping weights using equation (\ref{eq:shoppingprob}).

\clearpage

\begin{table}[H]
    \centering
\def\sym#1{\ifmmode^{#1}\else\(^{#1}\)\fi}
\begin{tabular}{l*{2}{c}}
\hline\hline
            &\multicolumn{1}{c}{(1)}&\multicolumn{1}{c}{(2)}\\
\hline
Log distance ($-\zeta\theta$) &     -1.9750\sym{***}&     -1.7138\sym{***}\\
            &    (0.0329)         &    (0.0422)         \\
\hline
Method & PPML & OLS \\
Origin $\times$ industry fixed effects & Yes & Yes\\
Destination $\times$ industry fixed effects & Yes & Yes\\
\(N\)       &       1,728         &         672         \\
\(R^{2}\)   &                     &       0.850         \\
\hline\hline
\multicolumn{3}{l}{\footnotesize Heteroskedasticity-robust standard errors in parentheses}\\
\multicolumn{3}{l}{\footnotesize \sym{*} \(p<0.1\), \sym{**} \(p<0.05\), \sym{***} \(p<0.01\)}\\
\end{tabular}
\caption{Pooled industry-conditioned commuting gravity estimates}\label{tab:gravity_commutingDT}
    \floatfoot{\textit{Notes:} The table reports gravity estimates for commuting flows from the 1977 Okinawa person trip survey. Agriculture, manufacturing, and non-tradable services are pooled with a common distance coefficient and origin-by-industry and destination-by-industry fixed effects. PPML includes all $24\times24\times3$ origin--destination--industry cells, including zeros; OLS uses log flows for positive cells. Within-municipality distances are imputed as described in the text.
    }
\end{table}


\begin{table}[H]
    \centering
\small
\setlength{\tabcolsep}{4pt}
\def\sym#1{\ifmmode^{#1}\else\(^{#1}\)\fi}
\begin{tabular}{l*{4}{c}}
\hline\hline
            &\multicolumn{1}{c}{(1)}&\multicolumn{1}{c}{(2)}&\multicolumn{1}{c}{(3)}&\multicolumn{1}{c}{(4)}\\
\hline
Log distance&     -2.2978\sym{***}&     -1.9928\sym{***}&     -3.6199\sym{***}&     -2.4406\sym{***}\\
            &    (0.1506)         &    (0.1400)         &    (0.1475)         &    (0.1369)         \\
\hline
Method & PPML &  OLS & PPML &  OLS  \\ 
Origin fixed effects & No & No & Yes & Yes\\
Destination fixed effects & No & No & Yes & Yes\\
\(N\)       &         576         &         146         &         576         &         146         \\
\(R^{2}\)   &                     &       0.495         &                     &       0.885         \\
\hline\hline
\multicolumn{5}{l}{\footnotesize Heteroskedasticity-robust standard errors in parentheses}\\
\multicolumn{5}{l}{\footnotesize \sym{*} \(p<0.1\), \sym{**} \(p<0.05\), \sym{***} \(p<0.01\)}\\
\end{tabular}
\caption{Shopping gravity estimates with and without origin and destination fixed effects}\label{tab:gravity_shopping}
    \floatfoot{\textit{Notes:} The sample contains one home-originating shopping or leisure trip per respondent from the 1977 Okinawa person trip survey. Destination shares use survey expansion weights and are normalized within each origin. PPML includes all $24\times24$ origin--destination pairs, including zero shares; OLS uses log shares for positive pairs. Columns (1)--(2) omit origin and destination fixed effects; columns (3)--(4) include both. Distances use the digitized 1973 road network, with within-municipality distances imputed as described in the text.
    }
\end{table}

\clearpage

\subsection{Calibration using GNI component shares}\label{appsec:calibration_GNI}

We use the three base-related components reported separately in the Okinawan GNI statistics: labor income earned by Okinawan workers employed at US bases, land rents paid for US-base land, and US personnel's purchases of Okinawan goods and services. Let their observed shares of GNI be denoted by $s_E$, $s_R$, and $s_C$, respectively.

\paragraph{GNI in the equilibrium.}

Model-implied GNI is the total income of Okinawan workers and landowners:
\begin{equation*}
    \begin{split}
    Y &= \sum_i \sum_s w_{is}L_{is}+\sum_i \Omega_i\\
    &= \underbrace{\sum_i \sum_{s \in \{ A,M,N \}} \frac{(1-\mu\gamma^s)w_{is}L_{is}}{1-\gamma^s}+(1-\alpha)(1-\mu)\sum_i \sum_{s \in \{ A,M,N,B \}} w_{is}L_{is}}_{\equiv Y_P}\\
    &\quad +\underbrace{\sum_i w_{iB}L_{iB}}_{\equiv Y_E}+\underbrace{\sum_i m_i\mathbb{Q}_i^{US}S_i}_{\equiv Y_R} .
    \end{split}
\end{equation*}
Here, $Y_P$ combines private-sector labor income and private-land income from production and housing; $Y_E$ is labor income from US-base employment; and $Y_R$ is rent paid for US-base land.

\paragraph{Matching the GNI share of the US-base labor income.}

Step 3 in Section \ref{sec:calibration} recovers private-sector wages and composite amenities. For the US-base sector, define $\omega_{iB}^\prime\equiv a_B\omega_{iB}$ and $\widetilde{A}_{iB}^\prime\equiv\widetilde{A}_{iB}/a_B$. We recover these variables from equations (\ref{eq:LMC}) and (\ref{eq:LMC_2}), normalize the geometric mean of $\widetilde{A}_{iB}^\prime$ to one, and set the scale of $\omega_{iB}^\prime$ to match the base-sector employment share. Equation (\ref{eq:emp_ratio_AM}) applies with $s=B$ and the primed variables because $\widetilde{A}_{iB}\omega_{iB}=\widetilde{A}_{iB}^\prime\omega_{iB}^\prime$.

We choose $a_B$ so that $Y_E/Y_P=s_E/(1-s_E-s_R)$. Given private-sector wages and $\omega_{iB}^\prime$,
\begin{equation*}
    \begin{split}
        Y_P&=\sum_i \sum_{s \in \{ A,M,N \}} \frac{(1-\mu\gamma^s)w_{is}L_{is}}{1-\gamma^s}+(1-\alpha)(1-\mu)\sum_i \sum_{s \in \{ A,M,N \}} w_{is}L_{is}\\
        &\quad +(1-\alpha)(1-\mu)a_B^{-\frac{1}{\theta}} \sum_i (\omega_{iB}^\prime)^{\frac{1}{\theta}}L_{iB}
    \end{split}
\end{equation*}
and
\begin{equation*}
    Y_E=a_B^{-\frac{1}{\theta}} \sum_i (\omega_{iB}^\prime)^{\frac{1}{\theta}}L_{iB} .
\end{equation*}
Solving $Y_E/Y_P=s_E/(1-s_E-s_R)$ determines $a_B$. We then recover $w_{iB}=(\omega_{iB}^\prime/a_B)^{1/\theta}$ and $\widetilde{A}_{iB}=a_B\widetilde{A}_{iB}^\prime$.

\paragraph{Matching the GNI share of the land rents for the US bases.}

We set municipal US-base rent rates proportional to private-sector land rents and choose their scale to satisfy $Y_R/Y_E=s_R/s_E$:
\begin{equation*}
    \mathbb{Q}_i^{US}
    =\frac{s_R}{s_E}\frac{Y_E}{\sum_k m_k\mathbb{Q}_k S_k}\mathbb{Q}_i.
\end{equation*}
Thus, aggregate US-base land rents equal $(s_R/s_E)Y_E$.

\paragraph{Matching the GNI share of consumption by the US military personnel.}

We allocate the observed total number of US military personnel, $L^{US}$, across municipalities in proportion to US-base land area:
\begin{equation*}
    L_i^{US}=L^{US}\frac{m_iS_i}{\sum_k m_kS_k}.
\end{equation*}
Per-capita spending is $e_{US}=s_CY/L^{US}$, which matches US personnel's purchases to their observed share of GNI. We then obtain US-base-sourced consumption as $\Phi_i=m_i\mathbb{Q}_i^{US}S_i+e_{US}L_i^{US}$.

\clearpage

\subsection{Model fit}\label{appsec:modelfit}

\begin{figure}[H]
    \subfloat[Bilateral commuting probability\label{fig:modelfit_commuting}]{%
      \includegraphics[width=0.475\textwidth]{figure/modelfit_commuting1970.png}
    }
    \hfill
    \subfloat[House prices and floor-space rents\label{fig:modelfit_rent}]{%
      \includegraphics[width=0.475\textwidth]{figure/modelfit_rent_standard.pdf}
    }
    \vspace{0.2in}
    \caption{Comparison of the observed and calibrated values}
    \vspace{-0.2in}
    \floatfoot{\textbf{Note}: Panel \ref{fig:modelfit_commuting} compares 1967 commuting probabilities with model predictions for 1970. Each point represents one of 56 origin--destination pairs across eight regions; probabilities are normalized within origin regions. Panel \ref{fig:modelfit_rent} compares municipal house prices with model-implied floor-space rents. Each point represents a municipality, and both variables are standardized to mean zero and standard deviation one. Observed values are on the horizontal axes and model predictions on the vertical axes. Lines are OLS fits; correlations are 0.958 and 0.873, respectively. Sources: the 1967 Person Trip Survey (\href{https://www.library.pref.okinawa.jp/winj/opac/switch-detail.do?idx=1}{Okinawa Prefectural Library}) and the 1970 \href{https://ndlsearch.ndl.go.jp/books/R100000002-I000001206601}{\textit{shichoson koufuzei seido kaisetsu}}.
    }
    \label{fig:modelfit}
\end{figure}

\clearpage

\section{Counterfactual details}\label{appsec:counterfactual}

\subsection{Computational procedure}\label{appsec:counterfactual_comp}

At a given aggregate workforce $L$, we iterate between wages and sector-specific residence counts. The benchmark sets $\Phi_i=L_{iB}=0$ and holds occupied base land fixed. We initialize $\{R_{is}\}$ with $\sum_i\sum_s R_{is}=L$ and $R_{iB}=0$.

\paragraph{Inner loop.}
Holding $\{R_{is}\}$ fixed, we proceed as follows:
\begin{enumerate}
    \item Initialize $\{w_{iA},w_{iM},w_{iN}\}$ and set $w_{iB}=0$.
    \item Compute private-sector employment from (\ref{eq:workers}).
    \item Solve (\ref{eq:land_market}) for $Q_i$, using $H_i=(1-m_i)S_iQ_i^\phi$ under the normalization $P^T=1$. Obtain $H_{iN}$ from (\ref{eq:land_production}) and recover service prices from (\ref{eq:wageN}): $        P_i^N=(\widetilde b_i^N)^{-1}w_{iN}^{1-\gamma^N}Q_i^{\gamma^N}
        L_{iN}^{\frac{1-\gamma^N}{1-\sigma}}H_{iN}^{\frac{\gamma^N}{1-\sigma}}$.
    \item Update agricultural and manufacturing wages from (\ref{eq:wageA}) and (\ref{eq:wageM}), and service wages from (\ref{eq:goods_market}).
    \item Repeat steps 2--4 until wages converge.
\end{enumerate}

\paragraph{Main loop.}
Update $R_{is}=L\lambda_{is}^R$ using (\ref{eq:prob_live}) and (\ref{eq:residents}), then repeat the inner loop. Iterate until the residence distribution converges.

For experiments retaining base employment, each inner-loop iteration also updates base wages. Combining (\ref{eq:prob_live_work}) and (\ref{eq:workers}) gives
\begin{equation*}
    w_{iB}=\left[
        \frac{L_{iB}}{\displaystyle\sum_k
        \frac{A_{kB}R_k^P\tau_{ki}^{-\theta}}
        {\sum_{s\in\{A,M,N\}}A_{ks}\sum_l w_{ls}^\theta\tau_{kl}^{-\theta}}}
    \right]^{1/\theta},
    \qquad R_k^P=\sum_{s\in\{A,M,N\}}R_{ks}.
\end{equation*}
Base wages are zero where $L_{iB}=0$. The open-population model adds the outer loop in Appendix \ref{appsec:counterfactual_smallopen}.

\clearpage

\subsection{Local price and wage adjustment in the benchmark counterfactual}\label{appsec:counterfactual_prices}

Figure \ref{fig:cf_economic_map} maps benchmark changes in wages, floor-space rents, service prices, and varieties. Non-tradable wages fall in every municipality, and varieties fall in 32 of 35. Population inflows sustain service demand in the three municipalities where varieties increase.
Manufacturing wages rise everywhere as employment expansion raises productivity through agglomeration. Lower floor-space rents reinforce these gains around the bases; elsewhere, productivity gains outweigh higher rents. With fixed agricultural prices and productivity, agricultural wages rise where rents fall and decline where rents rise, as implied by equation \eqref{eq:wageA}.
Floor-space rents fall in Chatan ($-12.19\%$), Koza ($-10.77\%$), and Kin ($-10.96\%$), but rise in 19 municipalities, consistent with the residential reallocation in Figure \ref{fig:cf_pop_change_map}. Lower marginal costs outweigh the effect of fewer varieties on the CES price index, so service prices fall in every municipality.

\begin{figure}[p]
    \centering
    \includegraphics[width=0.96\linewidth]{map/cf_all_wages_rent_priceN_variety_maps_en.png}
    \vspace{0.1in}
    \caption{Municipal wage, rent, price, and variety responses in the benchmark counterfactual}
    \vspace{-0.2in}
    \floatfoot{\textbf{Note}: The maps report $100\times(\text{counterfactual}/\text{observed}-1)$ for sectoral wages, floor-space rents, the non-tradable price index, and non-tradable varieties. The benchmark sets $\Phi_i=L_{iB}=0$ and leaves base land unavailable for private use. Red denotes an increase and blue a decrease.}
    \label{fig:cf_economic_map}
\end{figure}

Figure \ref{fig:cf_result_scatter_prices} reports exposure gradients of $-0.610$ for non-tradable varieties, $-0.167$ for floor-space rents, and $-0.114$ for service prices. The wage gradients are $-0.317$ for non-tradables, $0.060$ for manufacturing, and $0.075$ for agriculture. These responses are stronger in municipalities with greater base exposure.

\begin{figure}[p]
    \centering
    \subfloat[Non-tradable varieties\label{fig:cf_variety}]{%
      \includegraphics[width=0.475\textwidth]{figure/cf_baseline_variety.pdf}
    }
    \hfill
    \subfloat[Non-tradable price index\label{fig:cf_priceN}]{%
      \includegraphics[width=0.475\textwidth]{figure/cf_baseline_priceN.pdf}
    } \\
    \subfloat[Floor-space rent\label{fig:cf_rent}]{%
      \includegraphics[width=0.475\textwidth]{figure/cf_baseline_rent.pdf}
    }
    \hfill
    \subfloat[Sectoral wages\label{fig:cf_wage}]{%
      \includegraphics[width=0.475\textwidth]{figure/cf_baseline_wage.pdf}
    }
    \vspace{0.2in}
    \caption{Exposure gradients in prices, rents, varieties, and wages in the benchmark counterfactual}
    \vspace{-0.2in}
    \floatfoot{\textbf{Note}: Each panel plots a counterfactual-to-observed ratio against the 1970 base-area share. In panel \ref{fig:cf_wage}, red circles, blue squares, and green triangles denote non-tradable, manufacturing, and agricultural wages. Lines are linear fits, and the boxes report slopes and standard errors. The benchmark sets $\Phi_i=L_{iB}=0$ while leaving base acreage unavailable for private use.}
    \label{fig:cf_result_scatter_prices}
\end{figure}

\clearpage

\subsection{Spatial results for the consumption and employment channels}\label{appsec:counterfactual_channels}

This subsection reports the maps and fitted gradients underlying the channel comparison in Section \ref{sec:decompose_bases}. Figures \ref{fig:cf_noFSI_pop_map} and \ref{fig:cf_noUSemploy_pop_map} show population responses when $\Phi_i=0$ and when local base employment is removed, respectively.

\begin{figure}[p]
    \centering
    \includegraphics[width=0.90\linewidth]{map/cf_noFSI_total_AMN_Rcf_over_R_maps_en.png}
    \vspace{0.1in}
    \caption{Municipal population responses when base-sourced expenditure and rent income are removed}
    \vspace{-0.2in}
    \floatfoot{\textbf{Note}: The maps report $100\times(\text{counterfactual}/\text{observed}-1)$ for total residents and residents working in agriculture (A), manufacturing (M), and civilian non-tradable services (N). The total includes base workers. The counterfactual sets $\Phi_i=0$, removing US personnel's purchases and base-rent income while retaining base employment and occupied base land. Red denotes an increase and blue a decrease.}
    \label{fig:cf_noFSI_pop_map}
\end{figure}

\begin{figure}[p]
    \centering
    \includegraphics[width=0.90\linewidth]{map/cf_noUSemploy_total_AMN_Rcf_over_R_maps_en.png}
    \vspace{0.1in}
    \caption{Municipal population responses when local base employment is removed}
    \vspace{-0.2in}
    \floatfoot{\textbf{Note}: The maps report $100\times(\text{counterfactual}/\text{observed}-1)$ for total residents and residents working in agriculture (A), manufacturing (M), and civilian non-tradable services (N). The observed total includes base workers. The counterfactual sets $L_{iB}=0$, retaining US personnel's purchases, base-rent income, and occupied base land. Red denotes an increase and blue a decrease.}
    \label{fig:cf_noUSemploy_pop_map}
\end{figure}

\begin{figure}[p]
    \centering
    \subfloat[Purchases and rents removed: residents\label{fig:cf_noFSI_pop}]{%
      \includegraphics[width=0.475\textwidth]{figure/cf_noFSI_residents.pdf}
    }
    \hfill
    \subfloat[Purchases and rents removed: sector shares\label{fig:cf_noFSI_empshare}]{%
      \includegraphics[width=0.475\textwidth]{figure/cf_noFSI_sector_share.pdf}
    } \\
    \subfloat[Base employment removed: residents\label{fig:cf_noUSemploy_pop}]{%
      \includegraphics[width=0.475\textwidth]{figure/cf_noUSemploy_residents.pdf}
    }
    \hfill
    \subfloat[Base employment removed: sector shares\label{fig:cf_noUSemploy_empshare}]{%
      \includegraphics[width=0.475\textwidth]{figure/cf_noUSemploy_sector_share.pdf}
    }
    \vspace{0.2in}
    \caption{Exposure gradients for the consumption and employment channels}
    \vspace{-0.2in}
    \floatfoot{\textbf{Note}: Each point is a municipality; the horizontal axis is the 1970 base-area share. Panels \ref{fig:cf_noFSI_pop} and \ref{fig:cf_noFSI_empshare} set $\Phi_i=0$; panels \ref{fig:cf_noUSemploy_pop} and \ref{fig:cf_noUSemploy_empshare} set $L_{iB}=0$. The other channel and occupied base land remain fixed. Resident panels plot counterfactual-to-observed ratios; share panels plot counterfactual-minus-observed sector-share changes among residents. Red circles denote tertiary (civilian non-tradables plus base employment), blue squares manufacturing, and green triangles agriculture. Lines are linear fits; boxes report slopes and standard errors.}
    \label{fig:cf_channel_scatter}
\end{figure}

Figures \ref{fig:cf_noFSI_economic_map} and \ref{fig:cf_noUSemploy_economic_map} show the corresponding local prices, wages, rents, and varieties. Their common color scales facilitate direct comparison of the two mechanisms.

\begin{figure}[p]
    \centering
    \includegraphics[width=0.96\linewidth]{map/cf_noFSI_wages_rent_priceN_variety_maps_en.png}
    \vspace{0.1in}
    \caption{Municipal economic outcomes when base-sourced expenditure and rent income are removed}
    \vspace{-0.2in}
    \floatfoot{\textbf{Note}: The maps report $100\times(\text{counterfactual}/\text{observed}-1)$ for sectoral wages, floor-space rents, the non-tradable price index, and non-tradable varieties. The counterfactual sets $\Phi_i=0$, removing US personnel's purchases and base-rent income while retaining base employment and occupied base land. Red denotes an increase and blue a decrease.}
    \label{fig:cf_noFSI_economic_map}
\end{figure}

\begin{figure}[p]
    \centering
    \includegraphics[width=0.96\linewidth]{map/cf_noUSemploy_wages_rent_priceN_variety_maps_en.png}
    \vspace{0.1in}
    \caption{Municipal economic outcomes when local base employment is removed}
    \vspace{-0.2in}
    \floatfoot{\textbf{Note}: The maps report $100\times(\text{counterfactual}/\text{observed}-1)$ for sectoral wages, floor-space rents, the non-tradable price index, and non-tradable varieties. The counterfactual sets $L_{iB}=0$, retaining US personnel's purchases, base-rent income, and occupied base land. Red denotes an increase and blue a decrease.}
    \label{fig:cf_noUSemploy_economic_map}
\end{figure}

\clearpage

\subsection{Fixed municipal residence}\label{appsec:counterfactual_fixed_residence}

The fixed-residence counterfactual in Section \ref{sec:counterfactual_migration} imposes $R_i^{cf}=R_i$ while allowing sector and workplace choices, commuting, and local prices to adjust. Figure \ref{fig:cf_fixR_pop_map} shows resident-sector reallocation, Figure \ref{fig:cf_fixed_residence_scatter} reports exposure gradients, and Figure \ref{fig:cf_fixR_economic_map} shows wages, rents, prices, and varieties. Sectoral and workplace reallocation remains strongly related to base exposure, while the rent and variety exposure gradients are flatter than in the benchmark.

\begin{figure}[p]
    \centering
    \includegraphics[width=0.90\linewidth]{map/cf_fixR_sector_total_AMN_Rcf_over_R_maps_en.png}
    \vspace{0.1in}
    \caption{Municipal resident-sector reallocation with municipal residence fixed}
    \vspace{-0.2in}
    \floatfoot{\textbf{Note}: The maps report $100\times(\text{counterfactual}/\text{observed}-1)$ for total residents and residents working in agriculture (A), manufacturing (M), and civilian non-tradable services (N). The observed total includes base workers. The counterfactual fixes $R_i^{cf}=R_i$, sets $\Phi_i=L_{iB}=0$, and leaves base land unavailable for private use. Red denotes an increase and blue a decrease. Municipal boundaries are those of 1970; labels report changes for selected municipalities.}
    \label{fig:cf_fixR_pop_map}
\end{figure}

\begin{figure}[p]
    \centering
    \subfloat[Broad sector-of-employment shares\label{fig:cf_fixR_empshare}]{%
      \includegraphics[width=0.475\textwidth]{figure/cf_womigration_sector_share.pdf}
    }
    \hfill
    \subfloat[Civilian workplace employment\label{fig:cf_fixR_employment}]{%
      \includegraphics[width=0.475\textwidth]{figure/cf_womigration_employment.pdf}
    } \\
    \subfloat[Floor-space rent\label{fig:cf_fixR_rent}]{%
      \includegraphics[width=0.475\textwidth]{figure/cf_womigration_rent.pdf}
    }
    \hfill
    \subfloat[Non-tradable varieties\label{fig:cf_fixR_variety}]{%
      \includegraphics[width=0.475\textwidth]{figure/cf_womigration_variety.pdf}
    }
    \vspace{0.2in}
    \caption{Exposure gradients with municipal residence fixed}
    \vspace{-0.2in}
    \floatfoot{\textbf{Note}: Each point is a municipality, and the horizontal axis is the 1970 base-area share. Panel \ref{fig:cf_fixR_empshare} plots counterfactual-minus-observed changes in resident sector-of-employment shares; the other panels plot counterfactual-to-observed ratios. In the sectoral panels, red circles, blue squares, and green triangles denote non-tradables (tertiary in panel \ref{fig:cf_fixR_empshare}), manufacturing, and agriculture. Tertiary combines civilian non-tradables and base employment. Lines are linear fits, and boxes report slopes and standard errors. The counterfactual fixes $R_i^{cf}=R_i$, sets $\Phi_i=L_{iB}=0$, and leaves base land unavailable for private use.}
    \label{fig:cf_fixed_residence_scatter}
\end{figure}

\begin{figure}[p]
    \centering
    \includegraphics[width=0.96\linewidth]{map/cf_fixR_sector_wages_rent_priceN_variety_maps_en.png}
    \vspace{0.1in}
    \caption{Municipal wage, rent, price, and variety responses with municipal residence fixed}
    \vspace{-0.2in}
    \floatfoot{\textbf{Note}: The maps report $100\times(\text{counterfactual}/\text{observed}-1)$ for sectoral wages, floor-space rents, the non-tradable price index, and non-tradable varieties. The counterfactual fixes $R_i^{cf}=R_i$, sets $\Phi_i=L_{iB}=0$, and leaves base land unavailable for private use. Red denotes an increase and blue a decrease.}
    \label{fig:cf_fixR_economic_map}
\end{figure}

\clearpage

\subsection{Returning base land to private use}\label{appsec:counterfactual_landreturn}

\begin{figure}[H]
    \centering
    \subfloat[Resident population\label{fig:cf_pop_return}]{%
      \includegraphics[width=0.475\textwidth]{figure/cf_return_residents.pdf}
    }
    \hfill
    \subfloat[Broad sector-of-employment shares\label{fig:cf_empshare_return}]{%
      \includegraphics[width=0.475\textwidth]{figure/cf_return_sector_share.pdf}
    }
    \vspace{0.2in}
    \caption{Spatial effects when US-base land is returned to private use}
    \vspace{-0.2in}
    \floatfoot{\textbf{Note}: The horizontal axis is the 1970 base-area share. The left panel plots the counterfactual-to-observed resident ratio. The right panel plots counterfactual-minus-observed changes in primary, secondary, and tertiary sector-of-employment shares among residents; tertiary combines civilian non-tradable services and base employment. The counterfactual sets $\Phi_i=L_{iB}=m_i=0$, so base land is available for private use. Symbols, fitted lines, slopes, and standard errors follow Figure \ref{fig:cf_result_scatter}.}
    \label{fig:cf_result_scatter_return}
\end{figure}

\clearpage


\begin{figure}[t]
    \centering
    \includegraphics[width=\textwidth]{map/cf_land_return_total_AMN_Rcf_over_R_maps_en.png}
    \vspace{0.15in}
    \caption{Municipal population changes when US-base land is returned to private use}
    \vspace{-0.2in}
    \floatfoot{\textbf{Note}: The maps report $100\times(\text{counterfactual}/\text{observed}-1)$ for total residents and for residents working in agriculture (A), manufacturing (M), and non-tradable services (N). Red denotes an increase and blue a decrease. The counterfactual sets $\Phi_i=L_{iB}=m_i=0$. Municipal boundaries are those of 1970; labels report changes for selected municipalities.}
    \label{fig:cf_pop_change_map_CF2}
\end{figure}


\begin{figure}[tp]
    \centering
    \subfloat[Non-tradable varieties\label{fig:cf_variety_return}]{%
      \includegraphics[width=0.475\textwidth]{figure/cf_return_variety.pdf}
    }
    \hfill
    \subfloat[Non-tradable price index\label{fig:cf_priceN_return}]{%
      \includegraphics[width=0.475\textwidth]{figure/cf_return_priceN.pdf}
    } \\
    \subfloat[Floor-space rent\label{fig:cf_rent_return}]{%
      \includegraphics[width=0.475\textwidth]{figure/cf_return_rent.pdf}
    }
    \hfill
    \subfloat[Sectoral wages\label{fig:cf_wage_return}]{%
      \includegraphics[width=0.475\textwidth]{figure/cf_return_wage.pdf}
    }
    \vspace{0.2in}
    \caption{Prices, rents, varieties, and wages when US-base land is returned}
    \vspace{-0.2in}
    \floatfoot{\textbf{Note}: Each panel plots a counterfactual-to-observed ratio against the 1970 base-area share. In the wage panel, red circles, blue squares, and green triangles denote non-tradable, manufacturing, and agricultural wages, respectively. Lines are linear fits, and the boxes report slopes and standard errors. The counterfactual sets $\Phi_i=L_{iB}=m_i=0$.}
    \label{fig:cf_result_scatter_prices_return}
\end{figure}

\clearpage

\subsection{Local and aggregate elasticity calculations}\label{appsec:local_elasticities}

\paragraph{Experiments and elasticities.}
Section \ref{sec:elasticity} considers separate 1\% increases in foreign-sourced income (FSI), $\Phi_j$, and Okinawan base employment, $L_{jB}$, in one municipality $j$ from the calibrated 1970 equilibrium. All other exogenous inputs, including fundamental amenities and productivity, occupied base land and Okinawa's total workforce, remain fixed.

For outcome $Z_i$ in response municipality $i$ and input $x_j$ in source municipality $j$, define
\begin{equation}
\label{eq:local_base_elasticities}
    E_{ij}^{Z,x}\equiv\frac{\partial\ln Z_i}{\partial\ln x_j},
    \qquad x_j\in\{\Phi_j,L_{jB}\}.
\end{equation}
We compute these elasticities by symmetric log differences, changing $x_j$ to $1.01x_j$ and $x_j/1.01$ and dividing the difference in log outcomes by $2\ln(1.01)$. The diagonal entry is the local response; off-diagonal entries are spillovers. For a 1\% input increase, $E_{ij}^{Z,x}$ gives the first-order percentage change in $Z_i$. Both panels of Table \ref{tab:local_base_elasticities} use the same 29 source municipalities with positive FSI and base employment, denoted by $\mathcal J$, and all 35 response municipalities. For comparability across panels, we exclude five municipalities with zero base employment and one with zero FSI, since the corresponding log elasticities are undefined at zero. This leads to 29 (out of 35) source municipalities.

\paragraph{Aggregation.}
From the matrix of elasticity of the response in municipality $i$ with respect to a shock in municipality $j$, we compute the aggregate response in level as follows.
Let $s_i^Z$ denote municipality $i$'s aggregation weight for outcome $Z$, with $\sum_i s_i^Z=1$; superscript $0$ denotes the initial equilibrium. Employment and resident-worker responses use shares of the corresponding initial worker counts as the weights. Non-base employment is $L_i^P=\sum_{s\in\{A,M,N\}}L_{is}$, and total workplace employment is $L_i=L_i^P+L_{iB}$. For sectoral wages, floor-space rents, and service prices, the weights are
\begin{equation*}
    s_i^{w_s}=\frac{w_{is}^0L_{is}^0}{\sum_k w_{ks}^0L_{ks}^0},
    \qquad
    s_i^Q=\frac{Q_i^0H_i^0}{\sum_k Q_k^0H_k^0},
    \qquad
    s_i^{P^N}=\frac{w_{iN}^0L_{iN}^0/(1-\gamma^N)}{\sum_k w_{kN}^0L_{kN}^0/(1-\gamma^N)}.
\end{equation*}
%Wage and rent indices hold initial employment and floor-space compositions fixed; the wage index is $\sum_i L_{is}^0w_{is}/\sum_i L_{is}^0$. Floor space includes residential and production uses. 
The island-level service-price index is defined as $\prod_i(P_i^N)^{s_i^{P^N}}$, weighted by initial service-revenue shares across production locations.

For a shock in $j$, the aggregate response is $\sum_i s_i^ZE_{ij}^{Z,x}$, comprising the local contribution $s_j^ZE_{jj}^{Z,x}$ and spillovers $\sum_{i\ne j}s_i^ZE_{ij}^{Z,x}$. Table \ref{tab:local_base_elasticities}'s three columns are
\begin{align*}
    \text{Local response} &=\frac{1}{|\mathcal J|}\sum_{j\in\mathcal J}E_{jj}^{Z,x},\\
    \text{Okinawa-wide response} &=\frac{1}{|\mathcal J|}\sum_{j\in\mathcal J}\sum_i s_i^ZE_{ij}^{Z,x},\\
    \text{Local share of aggregate impact} &=100\frac{\sum_{j\in\mathcal J}s_j^ZE_{jj}^{Z,x}}{\sum_{j\in\mathcal J}\sum_i s_i^ZE_{ij}^{Z,x}}.
\end{align*}
The third column is the local share of the net aggregate change; negative values indicate offsetting local responses. NA denotes a zero aggregate response: resident workers are fixed in both experiments, as is non-base employment under FSI.

\paragraph{Worker-count effects.}
% For employment or resident workers, a 1\% increase in $j$ implies the first-order count change
% \begin{equation*}
%     \Delta Z_i^{(j,x)}=0.01Z_i^0E_{ij}^{Z,x}.
% \end{equation*}
% Total workplace employment changes combine non-base and base employment changes, $\Delta L_i=\Delta L_i^P+\Delta L_{iB}$; base employment changes only in the source municipality of the hiring experiment. Converting these totals to percentage changes and averaging across sources gives the local workplace responses of 0.0310\% under FSI and 0.1201\% under base hiring.

% For an FSI shock, we divide each tradable sector's employment loss by the service employment gain induced by the same shock. Averaging the source-municipality ratios equally across $\mathcal J$ gives losses of 0.112 agricultural and 0.088 manufacturing jobs per additional local service job. For the islandwide ratios, define
% \begin{equation*}
%     \Delta L_s^{(j)}=\sum_i\Delta L_{is}^{(j,\Phi)},
%     \qquad
%     \bar r_s=\frac{1}{|\mathcal J|}\sum_{j\in\mathcal J}
%     \frac{-\Delta L_s^{(j)}}{\Delta L_N^{(j)}},
%     \quad s\in\{A,M\}.
% \end{equation*}
% This gives $\bar r_A=0.372$ and $\bar r_M=0.628$. Fixed total labor supply and unchanged base employment imply a one-for-one offset between service gains and tradable losses across Okinawa, with manufacturing accounting for most of the losses.

Here we express the employment response as the absolute number of jobs in response to a one-job increase in the US bases. 
For base hiring, divide the count change by the additional $0.01L_{jB}^0$ base jobs. The local effect per added base job is
\begin{equation*}
    m_j^Z=\frac{Z_j^0E_{jj}^{Z,L_B}}{L_{jB}^0},
    \qquad
    \bar m^Z=\frac{1}{|\mathcal J|}\sum_{j\in\mathcal J}m_j^Z.
\end{equation*}
The mean non-base response is $-0.044$, giving a total local employment gain of $1-0.044=0.956$ jobs per added base job. Net inward commuting equals workplace employment minus resident workers, $L_i-R_i$. The increase of 0.956 local jobs therefore comprises 0.473 additional resident workers and an increase of $0.956-0.473=0.483$ in net inward commuting.

%Each additional base job displaces one non-base job across Okinawa because total labor supply is fixed. For non-base employment, the table's local contribution share weights source-specific responses by initial base employment: a 1\% increase adds more jobs in municipalities with larger bases. This gives a weighted local non-base response of $-0.095$ jobs per added base job and a local share of 9.5\%, whereas the worker-count ratios in the text use equal source weights.

% \paragraph{Comparisons with demand shocks.}
% For the consumer-spending comparison, define nominal goods and services expenditure excluding housing as
% \begin{equation*}
%     C_i=\alpha W_i+(1-\mu)Q_iH_i+\Phi_i,
% \end{equation*}
% where $W_i$ is residents' wage income, including commuting earnings, and $(1-\mu)Q_iH_i$ is private land-rent income. A 1\% FSI increase raises local $C_i$ by 0.244\% on average. The mean of the municipality-specific ratios $E_{jj}^{L_N,\Phi}/E_{jj}^{C,\Phi}$ is 0.412. \citeapp{mian2014explains_app} similarly divide the non-tradable employment response to housing net worth by the spending response in \citeapp{mian2013household}, obtaining $0.37/0.77\simeq0.48$. Their employment measure covers retail and restaurants, while spending covers automobiles and selected retail and restaurant purchases; our aggregates are broader.

% The response of approximately 0.15 for \citeapp{gerard2026cash_app} is calculated from their average 2010--2011 log effects on private formal employment and transfers: $0.020/0.134\simeq0.15$ (Table 2, Panel A, column 1). Our FSI elasticity of 0.0394 covers all non-base employment. Their employment expansion includes increased labor force participation, with little wage adjustment holding worker composition fixed. In our benchmark, local employment grows through commuting and relocation within a fixed islandwide workforce. \citeapp{mendes2026macroeconomic_app} find higher formal-sector real wages without a statistically significant increase in total employment, alongside growth in informal non-tradable employment. \citeapp{faber2019tourism_app} find that external demand, measured by hotel revenues, raises local employment, including manufacturing employment.

% \paragraph{Comparisons with employment shocks.}
% The total employment multipliers in \citeapp{hooker2001measuring_app} relate county-level employment losses to direct military and civilian job losses at closing bases. The relevant comparison is therefore our total local multiplier of 0.956; the non-base response of $-0.044$ excludes the added base job. \citeapp{chirakijja2024local_app} finds an approximately one-for-one increase in total local employment from prison openings, with little spillover to private employment. For public-sector employment in English local authorities during 2003--2007, \citeapp{faggio2014effect_app} estimate 0.085 private jobs per additional public job (standard error 0.325; Table 4, column 2). Their sectoral estimates are 0.518 non-tradable jobs and $-0.441$ manufacturing jobs (Table 6, column 2). Their non-tradable category includes construction, which belongs to the secondary sector in our data.

% \citeapp{zou2018local_app} finds civilian employment losses after military personnel reductions, with population adjusting mainly through lower in-migration, little wage adjustment, and falling rents. The military-base studies combine personnel and associated spending changes. Our experiments vary Okinawan base hiring and FSI separately, holding occupied land fixed and allowing base-worker payroll to adjust.

\clearpage

\subsection{Computing the open-population equilibrium}\label{appsec:counterfactual_smallopen}

Let $\theta_{\mathrm{out}}$ govern workers' choice between Okinawa and the outside economy. We use $\theta_{\mathrm{out}}\in\{0.5,1.1,2.19\}$ in the open-population exercises, while retaining $\theta=2.19$ for residence, workplace, and sector choice within Okinawa. With Fr\'{e}chet preferences at the outer margin, Okinawa's share of Japan's workforce is
\begin{equation}
\label{okinawa_res_share}
    \Lambda=\frac{\mathbb{W}^{\theta_{\mathrm{out}}}}{\mathbb{W}^{\theta_{\mathrm{out}}}+\overline{u}^{\theta_{\mathrm{out}}}},
\end{equation}
where $\mathbb{W}$ is expected welfare in (\ref{eq:welfare}) and $\overline u$ is welfare outside Okinawa. The benchmark fixes Okinawa's workforce and omits this migration margin.

Holding Japan's total workforce and $\overline u$ fixed, let $\widehat{\mathbb{W}}=\mathbb{W}/\mathbb{W}^0$ denote welfare relative to the initial equilibrium. Then
\begin{equation}
\label{okinawa_res_share_change}
    \widehat{\Lambda}\equiv\frac{\Lambda}{\Lambda^0}
    =\frac{\widehat{\mathbb{W}}^{\theta_{\mathrm{out}}}}
    {\Lambda^0\widehat{\mathbb{W}}^{\theta_{\mathrm{out}}}+1-\Lambda^0}
    =\frac{L}{L^0}.
\end{equation}
We set $\Lambda^0=0.0077$, using the Okinawa main island's share of Japan's 1970 population as a proxy for its worker share.

We use the following ``outer loop" to solve the model with endogeneous total population:

\paragraph{Outer loop.}
\begin{enumerate}
    \item Solve the inner and main loops in Appendix \ref{appsec:counterfactual_comp} at a trial aggregate workforce $L$ and compute $\mathbb{W}$.
    \item Use (\ref{okinawa_res_share_change}) to obtain the target workforce $L^0\widehat{\Lambda}$, with $L^0$, $\Lambda^0$, and $\mathbb{W}^0$ fixed at their initial values.
    \item Rescale all $R_{is}$ by $L^0\widehat{\Lambda}/L$ and repeat until $L=L^0\widehat{\Lambda}$.
\end{enumerate}

\clearpage

\subsection{Alternative external-market assumptions: specification and results}\label{appsec:counterfactual_external_maps}

This subsection provides the specifications and detailed results for the external-validity exercises in Section \ref{sec:counterfactual_external_validity}. Table \ref{tab:macro-counterfactual-summary-additional} reports aggregate outcomes across all parameter values; the figures below report spatial outcomes for the middle cases.

{\floatsetup[table]{capposition=top}
% Table maintained directly in draft.tex.
\begin{table}[!htbp]
  \centering
  \caption{Okinawa-wide effects of counterfactual scenarios under extended models}
  \label{tab:macro-counterfactual-summary-additional}
  \begin{threeparttable}
  \footnotesize
  \setlength{\tabcolsep}{3pt}
  \sisetup{
    input-signs = +-,
    retain-explicit-plus = true,
    table-number-alignment = center
  }
  \begin{tabularx}{\textwidth}{@{}>{\raggedright\arraybackslash}X *{7}{S[table-format=+2.2]}@{}}
    \toprule
    & \multicolumn{3}{c}{Aggregate effects (\%)} & \multicolumn{4}{c}{Employment-share changes (pp)} \\
    \cmidrule(lr){2-4} \cmidrule(l){5-8}
    Scenario & \multicolumn{1}{c}{\shortstack{Workers'\\Welfare}} & \multicolumn{1}{c}{GNI} & \multicolumn{1}{c}{\shortstack{Worker\\population}} & \multicolumn{1}{c}{Agric.} & \multicolumn{1}{c}{Manuf.} & \multicolumn{1}{c}{Non-trad.} & \multicolumn{1}{c}{U.S. base} \\
    \midrule
    \multicolumn{8}{@{}l}{\textit{Benchmark}} \\
    Benchmark counterfactual & -7.25 & -11.99 & 0.00 & +5.58 & +10.26 & -4.92 & -10.93 \\
    \addlinespace[0.35em]
    \multicolumn{8}{@{}l}{\textit{Larger geographical mobility in and out of Okinawa}} \\
    Low migration elasticity\newline ($\theta_{\mathrm{out}}=0.5$) & -7.26 & -15.21 & -3.67 & +5.74 & +10.19 & -5.01 & -10.93 \\
    Middle migration elasticity\newline ($\theta_{\mathrm{out}}=1.1$) & -7.27 & -18.93 & -7.91 & +5.93 & +10.11 & -5.12 & -10.93 \\
    High migration elasticity\newline ($\theta_{\mathrm{out}}=2.19$) & -7.28 & -25.29 & -15.16 & +6.28 & +9.96 & -5.32 & -10.93 \\
    \addlinespace[0.35em]
    \multicolumn{8}{@{}l}{\textit{Endogenous manufacturing goods prices}} \\
    High import elasticity\newline ($\varepsilon_M=20$) & -6.68 & -16.19 & 0.00 & +7.89 & +8.39 & -5.35 & -10.93 \\
    Middle import elasticity\newline ($\varepsilon_M=10$) & -6.25 & -18.64 & 0.00 & +9.41 & +7.13 & -5.62 & -10.93 \\
    Low import elasticity\newline ($\varepsilon_M=5$) & -5.63 & -21.53 & 0.00 & +11.43 & +5.45 & -5.95 & -10.93 \\
    Fixed import quantity\newline ($\varepsilon_M=0$) & -2.10 & -30.04 & 0.00 & +19.74 & -1.77 & -7.04 & -10.93 \\
    \bottomrule
  \end{tabularx}
  \begin{tablenotes}[flushleft]
    \footnotesize
    \item \textit{Notes:} Welfare, GNI, and worker-population entries are percentage changes from the initial equilibrium, $100(\text{counterfactual}/\text{initial}-1)$. Worker population is $\sum_i R_i$ across the 35 municipalities and adjusts in the external-migration scenarios. Employment entries are changes in Okinawa-wide sectoral shares in percentage points (pp). Initial shares are 16.55\%, 21.44\%, 51.08\%, and 10.93\% for agriculture, manufacturing, non-tradable services, and US-base employment, respectively. All scenarios set $\Phi_i=L_{iB}=0$ while keeping base land unavailable for private use.
  \end{tablenotes}
  \end{threeparttable}
\end{table}
}

\paragraph{Population adjustment through external migration.}
The open-population specification is described in Appendix \ref{appsec:counterfactual_smallopen}. We use $\theta_{\mathrm{out}}=1.1$ as the middle case, informed by \citeapp{ortega2013effect_app}, who estimate income elasticities of international migration flows of $0.76$ overall and $1.90$ within Europe. The low case sets $\theta_{\mathrm{out}}=0.5$; the high case uses $2.19$, equal to the within-Okinawa migration elasticity based on \citeapp{hayakawa2021highspeed_app}.

As $\theta_{\mathrm{out}}$ rises from $0.5$ to $1.1$ and $2.19$, worker population falls by $3.67\%$, $7.91\%$, and $15.16\%$, and GNI falls by $15.21\%$, $18.93\%$, and $25.29\%$ (Table \ref{tab:macro-counterfactual-summary-additional}). Welfare losses remain between $7.26\%$ and $7.28\%$, close to the benchmark loss of $7.25\%$, and sectoral employment-share changes remain similar. External migration mainly amplifies the contraction in population and total GNI.

\paragraph{Endogenous manufacturing prices and finite external trade elasticity.}
We next endogenize Okinawa's manufacturing-goods price through an external import-supply schedule while retaining agriculture as the numeraire. Let $\mathcal{I}^M$ denote Okinawa's net manufacturing import value and let a subscript $0$ denote the observed equilibrium. We assume that the manufacturing price $P^M$ is determined by
\begin{equation}
    \frac{P^M}{P_0^M}
    =\left(\frac{\mathcal{I}^M}{\mathcal{I}_0^M}\right)^{\frac{1}{1+\varepsilon_M}},
    \label{eq:inverse_import_supply}
\end{equation}
where $\varepsilon_M$ is the price elasticity of manufacturing import supply.\footnote{Equation (\ref{eq:inverse_import_supply}) implies that the import quantity of manufacturing goods, denoted by $\mathcal{Q}^M$, is given by $\mathcal{Q}^M/\mathcal{Q}_0^M=\left( P^M/P_0^M \right)^{\varepsilon_M}$, where $\mathcal{Q}_0^M$ represents the import quantity in the observed equilibrium.}
We specify the homothetic aggregator of tradable prices as $\Upsilon(P^A,P^M)=\left( P^A \right)^{1-\beta}\left( P^M \right)^\beta$ and set the manufacturing share in the tradable composite to $\beta=0.8$ to match the observed trade balance of the agricultural goods in the calibrated economy.\footnote{According to the Okinawa Statistical Yearbook, Okinawa's trade balance for the agricultural goods around 1970 was roughly balanced. }

Net manufacturing import value equals expenditure on manufacturing goods less local manufacturing output:
\begin{equation}
\label{eq:manufacturing_import_accounting}
    \mathcal{I}^M=\sum_i\left[
        \frac{\chi_i^M}{\alpha}C_i
        +\beta\mu Q_iH_i
        -\frac{w_{iM}L_{iM}}{1-\gamma^M}
    \right].
\end{equation}
Here, $\chi_i^M=\alpha\beta(P^T/\mathbb{P}_i)^{1-\kappa}$, and $C_i=\alpha W_i+(1-\mu)Q_iH_i+\Phi_i$ is expenditure on goods and services excluding housing. Residents' wage income $W_i$ includes commuting earnings. The three terms are manufacturing consumption, manufacturing inputs used in floor-space construction, and local manufacturing output. 

We use $\varepsilon_M=10$ as the middle case, guided by \citeapp{nicita2018cooperation_app}, and $5$ and $20$ for sensitivity analysis. At $\varepsilon_M=0$, the physical import quantity is fixed; as $\varepsilon_M\rightarrow\infty$, $P^M=P_0^M$, recovering the benchmark.


Table \ref{tab:macro-counterfactual-summary-additional} shows that at $\varepsilon_M=10$, welfare falls by $6.25\%$ and GNI by $18.64\%$, compared with $7.25\%$ and $11.99\%$ in the benchmark. Agricultural and manufacturing employment shares rise by $9.41$ and $7.13$ percentage points, compared with $5.58$ and $10.26$ points. Lower import-supply elasticity produces a larger manufacturing-price decline, redirects more workers toward agriculture, and increases the GNI loss. Cheaper manufactured goods reduce consumption and construction costs, cushioning the welfare loss. With fixed import quantity, manufacturing's employment share falls by $1.77$ points.

\paragraph{Spatial results.}

For external migration with $\theta_{\mathrm{out}}=1.1$, Figures \ref{fig:cf_openecon_middle_pop_map} and \ref{fig:cf_openecon_middle_workplace_map} map residents and workplace employment, and Figure \ref{fig:cf_openecon_middle_economic_map} maps wages, rents, service prices, and varieties. Figure \ref{fig:cf_openecon_middle_spatial} plots the corresponding municipal outcomes against base-area shares. The spatial patterns remain close to the benchmark: more base-exposed municipalities experience larger population losses and stronger shifts from tertiary employment toward manufacturing and agriculture. The resident-exposure gradient is $-0.445$, compared with $-0.477$ in the benchmark, and the sectoral-share gradients are also nearly unchanged. External migration amplifies the islandwide population decline while preserving these spatial patterns.

For elastic manufacturing import supply with $\varepsilon_M=10$, Figures \ref{fig:cf_import_middle_pop_map} and \ref{fig:cf_import_middle_workplace_map} map residents and workplace employment, and Figure \ref{fig:cf_import_middle_economic_map} maps wages, rents, service prices, and varieties. Figure \ref{fig:cf_import_middle_spatial} reports the corresponding scatterplots against base-area shares. These spatial patterns also remain close to the benchmark: population losses and the shift from tertiary employment toward tradable production are larger in more base-exposed municipalities. The resident-exposure gradient is $-0.538$, compared with $-0.477$ in the benchmark. Manufacturing-price adjustment changes aggregate income and the balance between agriculture and manufacturing, but preserves the main spatial pattern of reallocation away from services and base employment.

\begin{figure}[p]
    \centering
    \includegraphics[width=0.92\linewidth]{map/cf_openecon_middle_total_AMN_Rcf_over_R_maps_en.png}
    \vspace{0.1in}
    \caption{Municipal worker-resident responses with external migration: middle elasticity}
    \vspace{-0.2in}
    \floatfoot{\textbf{Note}: The maps report $100\times(\text{counterfactual}/\text{observed}-1)$ for total residents and residents working in agriculture (A), manufacturing (M), and non-tradable services (N). The observed total includes base workers. Red denotes an increase and blue a decrease. The maps use 1970 municipal boundaries and common color scales across scenarios. The counterfactual sets $\Phi_i=L_{iB}=0$ and $\theta_{\mathrm{out}}=1.1$, with base land unavailable for private use.}
    \label{fig:cf_openecon_middle_pop_map}
\end{figure}

\begin{figure}[p]
    \centering
    \includegraphics[width=0.92\linewidth]{map/cf_openecon_middle_total_AMN_Lcf_over_L_maps_en.png}
    \vspace{0.1in}
    \caption{Municipal workplace-employment responses with external migration: middle elasticity}
    \vspace{-0.2in}
    \floatfoot{\textbf{Note}: The maps report $100\times(\text{counterfactual}/\text{observed}-1)$ for total workplace employment and employment in agriculture (A), manufacturing (M), and non-tradable services (N). The observed total includes base jobs. Red denotes an increase and blue a decrease. The maps use 1970 municipal boundaries and common color scales across scenarios. The counterfactual sets $\Phi_i=L_{iB}=0$ and $\theta_{\mathrm{out}}=1.1$, with base land unavailable for private use.}
    \label{fig:cf_openecon_middle_workplace_map}
\end{figure}

\begin{figure}[p]
    \centering
    \includegraphics[width=0.82\linewidth]{map/cf_openecon_middle_wages_rent_priceN_variety_maps_en.png}
    \vspace{0.1in}
    \caption{Municipal wage, rent, price, and variety responses with external migration: middle elasticity}
    \vspace{-0.2in}
    \floatfoot{\textbf{Note}: The maps report $100\times(\text{counterfactual}/\text{observed}-1)$ for sectoral wages, floor-space rents, the non-tradable price index, and non-tradable varieties. Red denotes an increase and blue a decrease. The maps use 1970 municipal boundaries and common color scales across scenarios. The counterfactual sets $\Phi_i=L_{iB}=0$ and $\theta_{\mathrm{out}}=1.1$, with base land unavailable for private use.}
    \label{fig:cf_openecon_middle_economic_map}
\end{figure}

\begin{figure}[p]
    \centering
    \subfloat[Residents\label{fig:cf_openecon_middle_residents}]{%
      \includegraphics[width=0.475\textwidth]{figure/cf_openecon_middle_residents.pdf}
    }
    \hfill
    \subfloat[Broad sector-of-employment shares\label{fig:cf_openecon_middle_sector_share}]{%
      \includegraphics[width=0.475\textwidth]{figure/cf_openecon_middle_sector_share.pdf}
    }
    \vspace{0.15in}

    \subfloat[Workplace employment\label{fig:cf_openecon_middle_employment}]{%
      \includegraphics[width=0.475\textwidth]{figure/cf_openecon_middle_employment.pdf}
    }
    \hfill
    \subfloat[Non-tradable varieties\label{fig:cf_openecon_middle_variety}]{%
      \includegraphics[width=0.475\textwidth]{figure/cf_openecon_middle_variety.pdf}
    }
    \vspace{0.15in}
    \caption{Exposure gradients with external migration: middle elasticity}
    \vspace{-0.2in}
    \floatfoot{\textbf{Note}: Each point is a municipality, and the horizontal axis is the 1970 base-area share. Panel \ref{fig:cf_openecon_middle_sector_share} plots counterfactual-minus-observed changes in resident sector-of-employment shares; the other panels plot counterfactual-to-observed ratios. In the sectoral panels, red circles denote non-tradable services (including base employment in the tertiary share); blue squares and green triangles denote manufacturing and agriculture. Lines are linear fits, and the boxes report slopes and standard errors. The counterfactual sets $\Phi_i=L_{iB}=0$ and $\theta_{\mathrm{out}}=1.1$, with base land unavailable for private use.}
    \label{fig:cf_openecon_middle_spatial}
\end{figure}

\begin{figure}[p]
    \centering
    \includegraphics[width=0.92\linewidth]{map/cf_import_epsMis10_total_AMN_Rcf_over_R_maps_en.png}
    \vspace{0.1in}
    \caption{Municipal worker-resident responses with an endogenous manufacturing price: middle import elasticity}
    \vspace{-0.2in}
    \floatfoot{\textbf{Note}: The maps report $100\times(\text{counterfactual}/\text{observed}-1)$ for total residents and residents working in agriculture (A), manufacturing (M), and non-tradable services (N). The observed total includes base workers. Red denotes an increase and blue a decrease. The maps use 1970 municipal boundaries and common color scales across scenarios. The counterfactual sets $\Phi_i=L_{iB}=0$ and $\varepsilon_M=10$, with base land unavailable for private use.}
    \label{fig:cf_import_middle_pop_map}
\end{figure}

\begin{figure}[p]
    \centering
    \includegraphics[width=0.92\linewidth]{map/cf_import_epsMis10_total_AMN_Lcf_over_L_maps_en.png}
    \vspace{0.1in}
    \caption{Municipal workplace-employment responses with an endogenous manufacturing price: middle import elasticity}
    \vspace{-0.2in}
    \floatfoot{\textbf{Note}: The maps report $100\times(\text{counterfactual}/\text{observed}-1)$ for total workplace employment and employment in agriculture (A), manufacturing (M), and non-tradable services (N). The observed total includes base jobs. Red denotes an increase and blue a decrease. The maps use 1970 municipal boundaries and common color scales across scenarios. The counterfactual sets $\Phi_i=L_{iB}=0$ and $\varepsilon_M=10$, with base land unavailable for private use.}
    \label{fig:cf_import_middle_workplace_map}
\end{figure}

\begin{figure}[p]
    \centering
    \includegraphics[width=0.82\linewidth]{map/cf_import_epsMis10_wages_rent_priceN_variety_maps_en.png}
    \vspace{0.1in}
    \caption{Municipal wage, rent, price, and variety responses with an endogenous manufacturing price: middle import elasticity}
    \vspace{-0.2in}
    \floatfoot{\textbf{Note}: The maps report $100\times(\text{counterfactual}/\text{observed}-1)$ for sectoral wages, floor-space rents, the non-tradable price index, and non-tradable varieties. Red denotes an increase and blue a decrease. The maps use 1970 municipal boundaries and common color scales across scenarios. The counterfactual sets $\Phi_i=L_{iB}=0$ and $\varepsilon_M=10$, with base land unavailable for private use.}
    \label{fig:cf_import_middle_economic_map}
\end{figure}

\begin{figure}[p]
    \centering
    \subfloat[Residents\label{fig:cf_import_middle_residents}]{%
      \includegraphics[width=0.475\textwidth]{figure/cf_import_epsMis10_residents.pdf}
    }
    \hfill
    \subfloat[Broad sector-of-employment shares\label{fig:cf_import_middle_sector_share}]{%
      \includegraphics[width=0.475\textwidth]{figure/cf_import_epsMis10_sector_share.pdf}
    }
    \vspace{0.15in}

    \subfloat[Workplace employment\label{fig:cf_import_middle_employment}]{%
      \includegraphics[width=0.475\textwidth]{figure/cf_import_epsMis10_employment.pdf}
    }
    \hfill
    \subfloat[Sectoral wages\label{fig:cf_import_middle_wage}]{%
      \includegraphics[width=0.475\textwidth]{figure/cf_import_epsMis10_wage.pdf}
    }
    \vspace{0.15in}
    \caption{Exposure gradients with an endogenous manufacturing price: middle import elasticity}
    \vspace{-0.2in}
    \floatfoot{\textbf{Note}: Each point is a municipality, and the horizontal axis is the 1970 base-area share. Panel \ref{fig:cf_import_middle_sector_share} plots counterfactual-minus-observed changes in resident sector-of-employment shares; the other panels plot counterfactual-to-observed ratios. In the sectoral panels, red circles denote non-tradable services (including base employment in the tertiary share); blue squares and green triangles denote manufacturing and agriculture. Lines are linear fits, and the boxes report slopes and standard errors. The counterfactual sets $\Phi_i=L_{iB}=0$ and $\varepsilon_M=10$, with base land unavailable for private use.}
    \label{fig:cf_import_middle_spatial}
\end{figure}

\clearpage

% \paragraph{Decomposition.}

% Let ``OBS'' denote the observed outcomes. In addition, let ``CF1'' and ``CF2'' denote the counterfactual scenario in Sections \ref{sec:counterfactual_main} and Section \ref{sec:counterfactual_land}, respectively. 
% Then, the \textit{total effect} of the US bases on the growth of population density, calculated in the baseline counterfactual exercise, is decomposed as follows:
% \begin{equation}
% \label{eq:popdens_decomp}
%     \underbrace{\ln\left( \frac{PopDensity_i^{OBS}}{PopDensity_i^{CF2}} \right)}_{\text{total effect}}=
%     \underbrace{\ln\left( \frac{PopDensity_i^{OBS}}{PopDensity_i^{CF1}} \right)}_{\text{US-base sourced conusumption and employment effect}}+
%     \underbrace{\ln\left( \frac{PopDensity_i^{CF1}}{PopDensity_i^{CF2}} \right)}_{\text{land-occupying effect}}.
% \end{equation}
% The decomposition of the effects on the residential population can be done in the same way.

% Similarly, changes in the employment share of residents in each municipality can be decomposed as follows:
% \begin{equation}
% \label{eq:empshare_decomp}
%     \underbrace{Share_{i,s}^{OBS}-Share_{i,s}^{CF2}}_{\text{total effect}}=
%     \underbrace{Share_{i,s}^{OBS}-Share_{i,s}^{CF1}}_{\text{US-base sourced consumption and employment effect}}+
%     \underbrace{Share_{i,s}^{CF1}-Share_{i,s}^{CF2}}_{\text{land-occupying effect}} ,
% \end{equation}
% where the subscript $s \in \{ \text{primary, secondary, tertiary} \}$ indicates the sector. Here, the tertiary sector includes the non-tradable service and US base sectors\footnote{
% Note that $Share_{i,tertiary}^{CF1}$ and $Share_{i,tertiary}^{CF2}$ are the same as the employment share of the non-tradable service sector because there is no US base employment in these counterfactual equilibria.
% }
% The primary and secondary sectors correspond to the agricultural and manufacturing sectors in the model, respectively.

% The difference between the observed and CF1 equilibria is whether the model contains the US-base sourced consumption and employment, and whether the area of land developable for private uses is the same across these two equilibria. Thus, the first terms of the right-hand sides in equations (\ref{eq:popdens_decomp}) and (\ref{eq:empshare_decomp}) are regarded as capturing the US-base sourced consumption and employment effect of the US bases. On the other hand, the difference between the CF1 and CF2 counterfactual equilibria is whether the land for the US bases becomes available for private uses. Therefore, the second terms of the right-hand sides in equations (\ref{eq:popdens_decomp}) and (\ref{eq:empshare_decomp}) reflect the land-occupying effect.

% Given that the total effect is additively separable into the US-base sourced consumption and employment and land-occupying effects, we can measure how much of the variation in the total effect is due to variations in the two effects. Specifically, following \citeapp{hottman2016quantifying_app}, we regress each of the US-base sourced consumption and employment and land-occupying effects on the total effect. Then, the estimated slope represents the share of the variation in the total effect that can be attributed to each of the two effects. By construction, the two slopes add up to one.\footnote{
% Note that the slope estimated in a univariate regression is expressed as $\mathrm{Cov}(x_i,y)/\mathrm{Var}(y)$ when we take $x_i$ and $y$ as the dependent and explanatory variables, respectively. If $y=x_1+x_2$ holds, we have $\mathrm{Var}(y)=\mathrm{Var}(x_1)+\mathrm{Var}(x_2)+2\mathrm{Cov}(x_1,x_2)$ and $\mathrm{Cov}(x_i,y)=\mathrm{Var}(x_i)+\mathrm{Cov}(x_i,x_j)$, where $i \in \{1,2\}$ and $j \neq i$. Thus, the slope corresponds to the fraction of the variance of $y$ explained by the variance of $x_i$ and its covariance with the other determinant of $y$.
% }

% \begin{table}[t]
% \centering
% \caption{（Outdated）Contributions of the pecuniary and land-occupying effects}
% \label{tab:var_decomp}
% \resizebox{\textwidth}{!}{
% \begin{tabular}[t]{lcccccccccc}
% \toprule\toprule
% \multicolumn{1}{c}{} & \multicolumn{2}{c}{Population} & \multicolumn{2}{c}{Density} & \multicolumn{2}{c}{Primary} & \multicolumn{2}{c}{Secondary} & \multicolumn{2}{c}{Tertiary} \\
% \cmidrule(l{3pt}r{3pt}){2-3} \cmidrule(l{3pt}r{3pt}){4-5} \cmidrule(l{3pt}r{3pt}){6-7} \cmidrule(l{3pt}r{3pt}){8-9} \cmidrule(l{3pt}r{3pt}){10-11}
%   & UBCEE & LOE & UBCEE  & LOE  & UBCEE   & LOE   & UBCEE    & LOE    & UBCEE     & LOE    \\
% \midrule
% slope & $1.040$*** & $-0.040$ & $0.383$*** & $0.617$*** & $0.984$*** & $0.016$*** & $1.005$*** & $-0.005$ & $0.992$*** & $0.008$*\\
%  & ($0.133$) & ($0.133$) & ($0.037$) & ($0.037$) & ($0.005$) & ($0.005$) & ($0.003$) & ($0.003$) & ($0.004$) & ($0.004$)\\
% \midrule
% Observations & $35$ & $35$ & $35$ & $35$ & $35$ & $35$ & $35$ & $35$ & $35$ & $35$\\
% R${}^2$ & $0.649$ & $0.003$ & $0.763$ & $0.893$ & $0.999$ & $0.217$ & $1.000$ & $0.070$ & $0.999$ & $0.098$\\
% \bottomrule\bottomrule
% \multicolumn{9}{l}{\footnotesize
% Standard errors are reported in parentheses.
% }\\
% \multicolumn{9}{l}{\footnotesize
% \rule{0pt}{1em}* p $<$ 0.1, ** p $<$ 0.05, *** p $<$ 0.01
% }\\
% \end{tabular}
% }
%    \floatfoot{\textbf{Note}: The table reports the results of regressing the total effects on either of the US-base sourced consumption and employment or land-occupying effects. The ten columns in the table are classified into five groups, with each having two columns. From the left to right, each group presents the results of the variance decomposition for the residential population, population density, primary employment share, secondary employment share, and tertiary employment share. In each group, the left and right columns show the contributions of the US-base income and land-occupying effects, respectively.
%     }
% \end{table}

% Table \ref{tab:var_decomp} presents the result of this variance decomposition. The third and fourth columns show the decomposition of the variation in the total effect on population density growth. The value of the slope reported in the column labeled with UBCEE and LOE represents the fraction of variation in the total effect attributed to the variation in the US-base sourced consumption and employment effect and the land-occupying effect, respectively. According to the results, the US-base income effect accounts for nearly 40\% of the observed variation in the total effect. In contrast, for the population level, almost all variance is explained by the US-base income effect.

% We also report the decomposition of the effects on changes in the primary, secondary, and tertiary sectors' employment shares. Note that in each sector, the slope for the US-base sourced consumption and employment effect is close to one and that for the land-occupying effect is negligible, indicating that the land-occupying effect contributes little. Thus, the US-base sourced consumption and employment effect is considered to be the main driver of structural transformation.


% \clearpage

\clearpage

{
\footnotesize
\bibliographystyleapp{aer}
\bibliographyapp{reference.bib}


}

\end{document}
